\input{preamble}

% OK, start here.
%
\begin{document}

\title{Faisceaux de modules}


\maketitle

\phantomsection
\label{section-phantom}

\tableofcontents

\section{Introduction}
\label{section-introduction}

\noindent
Dans ce chapitre, nous développons les notions de base relatives aux faisceaux
de modules. Cela comprend en particulier le cas des faisceaux abéliens, puisque
ceux-ci peuvent être considérés comme des faisceaux de
$\underline{\mathbf{Z}}$-modules. Les références de base sont
\cite{FAC}, \cite{EGA} et \cite{SGA4}.

\medskip\noindent
Nous étudions dans un autre chapitre le cas des faisceaux de modules sur les
topos annelés (voir Modules sur les sites, Section
\ref{sites-modules-section-introduction}), bien que nous y reprenions
essentiellement les considérations du présent chapitre.





\section{Pathologie}
\label{section-pathology}

\noindent
Un espace annelé est un couple formé d'un espace topologique $X$ et d'un
faisceau d'anneaux $\mathcal{O}$. Nous autorisons $\mathcal{O} = 0$ dans
la définition. Dans ce cas, la catégorie des modules possède un seul objet
(à savoir $0$). Elle demeure une catégorie abélienne, etc., mais elle est
quelque peu dégénérée. De même, le faisceau $\mathcal{O}$ peut être nul
sur certains ouverts de $X$, etc.

\medskip\noindent
Ce phénomène ne se produit pas pour les espaces localement annelés (que nous
considérerons plus loin).








\section{La catégorie abélienne des faisceaux de modules}
\label{section-kernels}

\noindent
Soit $(X, \mathcal{O}_X)$ un espace annelé, voir Faisceaux, Définition
\ref{sheaves-definition-ringed-space}. Soient $\mathcal{F}$ et $\mathcal{G}$
des faisceaux de $\mathcal{O}_X$-modules, voir Faisceaux, Définition
\ref{sheaves-definition-sheaf-modules}. Soient
$\varphi, \psi : \mathcal{F} \to \mathcal{G}$ des morphismes de faisceaux
de $\mathcal{O}_X$-modules. Nous définissons
$\varphi + \psi : \mathcal{F} \to \mathcal{G}$ comme l'application qui,
sur chaque ouvert $U \subset X$, est la somme des applications induites par
$\varphi$ et $\psi$. Il s'agit manifestement encore d'une application de
faisceaux de $\mathcal{O}_X$-modules. Il est également clair que la composition
des applications de $\mathcal{O}_X$-modules est bilinéaire pour cette addition.
Ainsi, $\textit{Mod}(\mathcal{O}_X)$ est une catégorie préadditive, voir
Homologie, Définition \ref{homology-definition-preadditive}.

\medskip\noindent
Nous noterons $0$ le faisceau de $\mathcal{O}_X$-modules qui prend la valeur
constante $\{0\}$ sur tout ouvert $U \subset X$. Il s'agit manifestement à la
fois d'un objet final et d'un objet initial de
$\textit{Mod}(\mathcal{O}_X)$. Étant donné un morphisme de
$\mathcal{O}_X$-modules $\varphi : \mathcal{F} \to \mathcal{G}$, les
assertions suivantes sont équivalentes :
(a) $\varphi$ est nul, (b) $\varphi$ se factorise par $0$,
(c) $\varphi$ est nul sur les sections au-dessus de tout ouvert $U$, et
(d) $\varphi_x = 0$ pour tout $x \in X$. Voir Faisceaux, Lemme
\ref{sheaves-lemma-points-exactness}.

\medskip\noindent
En outre, étant donnés deux faisceaux $\mathcal{F}$, $\mathcal{G}$ de
$\mathcal{O}_X$-modules, nous pouvons définir leur somme directe par
$$
\mathcal{F} \oplus \mathcal{G} = \mathcal{F} \times \mathcal{G}
$$
munie des applications évidentes $(i, j, p, q)$ comme dans Homologie,
Définition \ref{homology-definition-direct-sum}. Ainsi
$\textit{Mod}(\mathcal{O}_X)$ est une catégorie additive, voir Homologie,
Définition \ref{homology-definition-additive-category}.

\medskip\noindent
Soit $\varphi : \mathcal{F} \to \mathcal{G}$ un morphisme de
$\mathcal{O}_X$-modules. Nous pouvons définir $\Ker(\varphi)$ comme le
sous-faisceau de $\mathcal{F}$ dont les sections sont
$$
\Ker(\varphi)(U) =
\{ s \in \mathcal{F}(U) \mid \varphi(s) = 0 \text{ dans } \mathcal{G}(U)\}
$$
pour tout ouvert $U \subset X$. On vérifie aisément qu'il s'agit bien d'un
noyau dans la catégorie des $\mathcal{O}_X$-modules. Autrement dit, un
morphisme $\alpha : \mathcal{H} \to \mathcal{F}$ se factorise par
$\Ker(\varphi)$ si et seulement si $\varphi \circ \alpha = 0$. De plus, au
niveau des germes, nous avons $\Ker(\varphi)_x = \Ker(\varphi_x)$.

\medskip\noindent
D'autre part, nous définissons $\Coker(\varphi)$ comme le faisceau de
$\mathcal{O}_X$-modules associé au préfaisceau de
$\mathcal{O}_X$-modules défini par la règle
$$
U
\longmapsto
\Coker(\mathcal{F}(U)\to \mathcal{G}(U)) =
\mathcal{G}(U)/\varphi(\mathcal{F}(U)).
$$
Comme la prise des germes commute au passage au faisceau associé, voir
Faisceaux, Lemme \ref{sheaves-lemma-stalk-sheafification}, nous obtenons
$\Coker(\varphi)_x = \Coker(\varphi_x)$. L'application
$\mathcal{G} \to \Coker(\varphi)$ est donc surjective (comme application de
faisceaux d'ensembles), voir Faisceaux, Section
\ref{sheaves-section-exactness-points}. Pour montrer qu'il s'agit d'un
conoyau, remarquons que si
$\beta : \mathcal{G} \to \mathcal{H}$ est un morphisme de
$\mathcal{O}_X$-modules tel que $\beta \circ \varphi$ soit nul, alors, pour
tout ouvert $U \subset X$, $\beta$ induit une application de
$\mathcal{G}(U)/\varphi(\mathcal{F}(U))$ dans $\mathcal{H}(U)$. Par la
propriété universelle du faisceau associé (voir Faisceaux, Lemme
\ref{sheaves-lemma-sheafification-presheaf-modules}), nous obtenons une
application canonique $\Coker(\varphi) \to \mathcal{H}$ telle que
l'application $\beta$ initiale soit égale à la composée
$\mathcal{G} \to \Coker(\varphi) \to \mathcal{H}$. Le morphisme
$\Coker(\varphi) \to \mathcal{H}$ est unique en vertu de la surjectivité
mentionnée ci-dessus.

\begin{lemma}
\label{lemma-abelian}
Soit $(X, \mathcal{O}_X)$ un espace annelé. La catégorie
$\textit{Mod}(\mathcal{O}_X)$ est une catégorie abélienne. De plus, un complexe
$$
\mathcal{F} \to \mathcal{G} \to \mathcal{H}
$$
est exact en $\mathcal{G}$ si et seulement si, pour tout $x \in X$, le complexe
$$
\mathcal{F}_x \to \mathcal{G}_x \to \mathcal{H}_x
$$
est exact en $\mathcal{G}_x$.
\end{lemma}

\begin{proof}
D'après Homologie, Définition \ref{homology-definition-abelian-category},
il faut montrer que l'image et la coimage coïncident. D'après Faisceaux,
Lemme \ref{sheaves-lemma-points-exactness}, il suffit de montrer qu'elles ont
le même germe en tout $x \in X$. Par les constructions des noyaux et conoyaux
ci-dessus, ces germes sont respectivement la coimage et l'image dans la
catégorie des $\mathcal{O}_{X, x}$-modules. Le résultat découle donc du fait
que la catégorie des modules sur un anneau est abélienne.
\end{proof}

\noindent
En fait, la catégorie $\textit{Mod}(\mathcal{O}_X)$ possède bien d'autres
propriétés. Voici deux constructions que nous pouvons effectuer.
\begin{enumerate}
\item Étant donnés un ensemble $I$ et, pour chaque $i \in I$, un
$\mathcal{O}_X$-module, nous pouvons former le produit
$$
\prod\nolimits_{i \in I} \mathcal{F}_i
$$
qui est le faisceau associant à chaque ouvert $U$ le produit des modules
$\mathcal{F}_i(U)$. C'est aussi le produit catégorique, comme dans Catégories,
Définition \ref{categories-definition-product}.
\item Étant donnés un ensemble $I$ et, pour chaque $i \in I$, un
$\mathcal{O}_X$-module, nous pouvons former la somme directe
$$
\bigoplus\nolimits_{i \in I} \mathcal{F}_i
$$
qui est le {\it faisceau associé} au préfaisceau associant à chaque ouvert
$U$ la somme directe des modules $\mathcal{F}_i(U)$. C'est aussi le coproduit
catégorique, comme dans Catégories, Définition
\ref{categories-definition-coproduct}. Cela résulte de la propriété
universelle du faisceau associé.
\end{enumerate}
Nous en déduisons que toutes les limites projectives et inductives existent
dans $\textit{Mod}(\mathcal{O}_X)$.

\begin{lemma}
\label{lemma-limits-colimits}
Soit $(X, \mathcal{O}_X)$ un espace annelé.
\begin{enumerate}
\item Toutes les limites projectives existent dans
$\textit{Mod}(\mathcal{O}_X)$. Elles coïncident avec les limites projectives
correspondantes de préfaisceaux de $\mathcal{O}_X$-modules (autrement dit,
elles commutent à la prise des sections sur les ouverts).
\item Toutes les limites inductives existent dans
$\textit{Mod}(\mathcal{O}_X)$. Elles s'obtiennent en prenant le faisceau
associé à la limite inductive correspondante dans la catégorie des
préfaisceaux. La prise des limites inductives commute à la prise des germes.
\item Les limites inductives filtrantes sont exactes.
\item Les sommes directes finies coïncident avec les sommes directes finies
correspondantes de préfaisceaux de $\mathcal{O}_X$-modules.
\end{enumerate}
\end{lemma}

\begin{proof}
Comme $\textit{Mod}(\mathcal{O}_X)$ est abélienne (Lemme
\ref{lemma-abelian}), elle admet toutes les limites projectives et inductives
finies (Homologie, Lemme \ref{homology-lemma-colimit-abelian-category}).
L'existence et la description des limites projectives et inductives découlent
donc de celles des produits et coproduits (voir la discussion ci-dessus) et
des lemmes de Catégories \ref{categories-lemma-limits-products-equalizers} et
\ref{categories-lemma-colimits-coproducts-coequalizers}. Puisque le passage au
faisceau associé commute à la prise des germes, les limites inductives
commutent à la prise des germes. L'assertion (3) signifie que, pour un système
$0 \to \mathcal{F}_i \to \mathcal{G}_i \to \mathcal{H}_i \to 0$ de suites
exactes de $\mathcal{O}_X$-modules indexé par un ensemble filtrant $I$, la
suite $0 \to \colim \mathcal{F}_i \to \colim \mathcal{G}_i \to
\colim \mathcal{H}_i \to 0$ est encore exacte. Comme l'exactitude se vérifie
sur les germes (Lemme \ref{lemma-abelian}), cela découle du cas des modules,
qui est Algèbre, Lemme \ref{algebra-lemma-directed-colimit-exact}. Nous
omettons la démonstration de (4).
\end{proof}

\noindent
L'existence des limites projectives et inductives permet d'exprimer les
propriétés d'exactitude des foncteurs définis sur la catégorie des
$\mathcal{O}$-modules en termes de ces limites, comme dans Catégories,
Section \ref{categories-section-exact-functor}. Voir Homologie, Lemme
\ref{homology-lemma-exact-functor}, pour une description des propriétés
d'exactitude en termes de suites exactes courtes.

\begin{lemma}
\label{lemma-exactness-pushforward-pullback}
Soit $f : (X, \mathcal{O}_X) \to (Y, \mathcal{O}_Y)$ un morphisme
d'espaces annelés.
\begin{enumerate}
\item Le foncteur
$f_* : \textit{Mod}(\mathcal{O}_X) \to \textit{Mod}(\mathcal{O}_Y)$
est exact à gauche. En fait, il commute à toutes les limites projectives.
\item Le foncteur
$f^* : \textit{Mod}(\mathcal{O}_Y) \to \textit{Mod}(\mathcal{O}_X)$
est exact à droite. En fait, il commute à toutes les limites inductives.
\item Le foncteur image inverse
$f^{-1} : \textit{Ab}(Y) \to \textit{Ab}(X)$ sur les faisceaux abéliens est
exact.
\end{enumerate}
\end{lemma}

\begin{proof}
Les assertions (1) et (2) résultent de ce que $(f^*, f_*)$ est un couple de
foncteurs adjoints, voir Faisceaux, Lemme
\ref{sheaves-lemma-adjoint-pullback-pushforward-modules}, et Catégories,
Section \ref{categories-section-adjoint}. L'assertion (3) résulte de ce que
l'exactitude se vérifie sur les germes (Lemme \ref{lemma-abelian}) et de la
description des germes de l'image inverse, voir Faisceaux, Lemme
\ref{sheaves-lemma-pullback-abelian-stalk}.
\end{proof}

\begin{lemma}
\label{lemma-j-shriek-exact}
Soit $j : U \to X$ une immersion ouverte d'espaces topologiques. Le foncteur
$j_! : \textit{Ab}(U) \to \textit{Ab}(X)$ est exact.
\end{lemma}

\begin{proof}
Cela découle de la description des germes donnée dans Faisceaux, Lemme
\ref{sheaves-lemma-j-shriek-abelian}.
\end{proof}

\begin{lemma}
\label{lemma-section-direct-sum-quasi-compact}
Soit $(X, \mathcal{O}_X)$ un espace annelé. Soit $I$ un ensemble. Pour
$i \in I$, soit $\mathcal{F}_i$ un faisceau de $\mathcal{O}_X$-modules.
Pour tout ouvert quasi-compact $U \subset X$, l'application
$$
\bigoplus\nolimits_{i \in I} \mathcal{F}_i(U)
\longrightarrow
\left(\bigoplus\nolimits_{i \in I} \mathcal{F}_i\right)(U)
$$
est bijective.
\end{lemma}

\begin{proof}
Si $s$ est un élément du membre de droite, il existe un recouvrement ouvert
$U = \bigcup_{j \in J} U_j$ tel que $s|_{U_j}$ soit une somme finie
$\sum_{i \in I_j} s_{ji}$, où $s_{ji} \in \mathcal{F}_i(U_j)$. Comme $U$
est quasi-compact, nous pouvons supposer le recouvrement fini, c'est-à-dire
$J$ fini. Alors $I' = \bigcup_{j \in J} I_j$ est une partie finie de $I$.
Manifestement, $s$ est une section du sous-faisceau
$\bigoplus_{i \in I'} \mathcal{F}_i$. Le résultat découle du fait que, pour
une somme directe finie, le passage au faisceau associé n'est pas nécessaire ;
voir le Lemme \ref{lemma-limits-colimits} ci-dessus.
\end{proof}



\section{Sections des faisceaux de modules}
\label{section-sections}

\noindent
Soit $(X, \mathcal{O}_X)$ un espace annelé. Soit $\mathcal{F}$ un faisceau de
$\mathcal{O}_X$-modules. Soit
$s \in \Gamma(X, \mathcal{F}) = \mathcal{F}(X)$ une section globale. Il existe
une unique {\it application de $\mathcal{O}_X$-modules}
$$
\mathcal{O}_X \longrightarrow \mathcal{F}, \ f \longmapsto fs
$$
{\it associée à $s$}. La notation ci-dessus signifie qu'une section locale
$f$ de $\mathcal{O}_X$, c'est-à-dire une section $f$ sur un ouvert $U$, est
envoyée sur le produit de $f$ par la restriction de $s$ à $U$.
Réciproquement, toute application
$\varphi : \mathcal{O}_X \to \mathcal{F}$ donne une section
$s = \varphi(1)$ telle que $\varphi$ soit le morphisme associé à $s$.

\begin{definition}
\label{definition-globally-generated}
Soit $(X, \mathcal{O}_X)$ un espace annelé. Soit $\mathcal{F}$ un faisceau de
$\mathcal{O}_X$-modules. Nous disons que $\mathcal{F}$ est {\it engendré par
ses sections globales} s'il existe un ensemble $I$ et des sections globales
$s_i \in \Gamma(X, \mathcal{F})$, $i \in I$, tels que l'application
$$
\bigoplus\nolimits_{i \in I}
\mathcal{O}_X \longrightarrow \mathcal{F}
$$
qui est l'application associée à $s_i$ sur le facteur correspondant à $i$,
soit surjective. Dans ce cas, nous disons que les sections $s_i$
{\it engendrent} $\mathcal{F}$.
\end{definition}

\noindent
Nous emploierons souvent l'abus de notation introduit dans Faisceaux,
Section \ref{sheaves-section-stalks} : si $s$ est une section locale de
$\mathcal{F}$ définie dans un voisinage ouvert d'un point $x \in X$, nous
notons $s_x$, voire $s$, l'image de $s$ dans le germe $\mathcal{F}_x$.

\begin{lemma}
\label{lemma-globally-generated}
Soit $(X, \mathcal{O}_X)$ un espace annelé. Soit $\mathcal{F}$ un faisceau de
$\mathcal{O}_X$-modules. Soit $I$ un ensemble. Soient
$s_i \in \Gamma(X, \mathcal{F})$, $i \in I$, des sections globales. Les
sections $s_i$ engendrent $\mathcal{F}$ si et seulement si, pour tout
$x\in X$, les éléments $s_{i, x} \in \mathcal{F}_x$ engendrent le
$\mathcal{O}_{X, x}$-module $\mathcal{F}_x$.
\end{lemma}

\begin{proof}
Omis.
\end{proof}

\begin{lemma}
\label{lemma-tensor-product-globally-generated}
\begin{slogan}
Le produit tensoriel de faisceaux de modules engendrés par leurs sections
globales est engendré par ses sections globales.
\end{slogan}
Soit $(X, \mathcal{O}_X)$ un espace annelé. Soient $\mathcal{F}$ et
$\mathcal{G}$ des faisceaux de $\mathcal{O}_X$-modules. Si $\mathcal{F}$ et
$\mathcal{G}$ sont engendrés par leurs sections globales, il en est de même de
$\mathcal{F} \otimes_{\mathcal{O}_X} \mathcal{G}$.
\end{lemma}

\begin{proof}
Omis.
\end{proof}

\begin{lemma}
\label{lemma-generated-by-local-sections}
Soit $(X, \mathcal{O}_X)$ un espace annelé. Soit $\mathcal{F}$ un faisceau de
$\mathcal{O}_X$-modules. Soit $I$ un ensemble. Soient $s_i$, $i \in I$, une
famille de sections locales de $\mathcal{F}$, c'est-à-dire
$s_i \in \mathcal{F}(U_i)$ pour certains ouverts $U_i \subset X$. Il existe
un unique plus petit sous-faisceau de $\mathcal{O}_X$-modules $\mathcal{G}$
tel que chaque $s_i$ corresponde à une section locale de $\mathcal{G}$.
\end{lemma}

\begin{proof}
Considérons le sous-préfaisceau de $\mathcal{O}_X$-modules défini par la règle
$$
U
\longmapsto
\{
\text{sommes } \sum\nolimits_{i \in J} f_i (s_i|_U)
\text{ où } J \text{ est fini, }
U \subset U_i \text{ pour } i\in J, \text{ et }
f_i \in \mathcal{O}_X(U)
\}
$$
Soit $\mathcal{G}$ le faisceau associé à ce sous-préfaisceau. C'est un
sous-faisceau de $\mathcal{F}$ d'après Faisceaux, Lemme
\ref{sheaves-lemma-sheafify-epi-mono}. Puisque toutes les sommes finies
doivent manifestement appartenir à $\mathcal{G}$, c'est bien le plus petit
sous-faisceau voulu.
\end{proof}

\begin{definition}
\label{definition-generated-by-local-sections}
Soit $(X, \mathcal{O}_X)$ un espace annelé. Soit $\mathcal{F}$ un faisceau de
$\mathcal{O}_X$-modules. Étant donnés un ensemble $I$ et des sections locales
$s_i$, $i \in I$, de $\mathcal{F}$, nous disons que le sous-faisceau
$\mathcal{G}$ du Lemme \ref{lemma-generated-by-local-sections} ci-dessus est
le {\it sous-faisceau engendré par les $s_i$}.
\end{definition}

\begin{lemma}
\label{lemma-generated-by-local-sections-stalk}
Soit $(X, \mathcal{O}_X)$ un espace annelé. Soit $\mathcal{F}$ un faisceau de
$\mathcal{O}_X$-modules. Soient un ensemble $I$ et des sections locales
$s_i$, $i \in I$, de $\mathcal{F}$. Soit $\mathcal{G}$ le sous-faisceau
engendré par les $s_i$, et soit $x\in X$. Alors $\mathcal{G}_x$ est le
sous-$\mathcal{O}_{X, x}$-module de $\mathcal{F}_x$ engendré par les éléments
$s_{i, x}$ pour les indices $i$ tels que $s_i$ soit définie en $x$.
\end{lemma}

\begin{proof}
Cela résulte immédiatement de la construction de $\mathcal{G}$ dans la
démonstration du Lemme \ref{lemma-generated-by-local-sections}.
\end{proof}










\section{Supports des modules et des sections}
\label{section-support}

\begin{definition}
\label{definition-support}
Soit $(X, \mathcal{O}_X)$ un espace annelé. Soit $\mathcal{F}$ un faisceau de
$\mathcal{O}_X$-modules.
\begin{enumerate}
\item Le {\it support de $\mathcal{F}$} est l'ensemble des points
$x \in X$ tels que $\mathcal{F}_x \not = 0$.
\item Nous notons $\text{Supp}(\mathcal{F})$ le support de $\mathcal{F}$.
\item Soit $s \in \Gamma(X, \mathcal{F})$ une section globale. Le
{\it support de $s$} est l'ensemble des points $x \in X$ tels que l'image
$s_x \in \mathcal{F}_x$ de $s$ soit non nulle.
\end{enumerate}
\end{definition}

\noindent
Le support d'une section locale se trouve bien entendu ainsi défini, puisqu'une
section locale est une section globale de la restriction de $\mathcal{F}$.

\begin{lemma}
\label{lemma-support-section-closed}
Soit $(X, \mathcal{O}_X)$ un espace annelé. Soit $\mathcal{F}$ un faisceau de
$\mathcal{O}_X$-modules. Soit $U \subset X$ un ouvert.
\begin{enumerate}
\item Le support de $s \in \mathcal{F}(U)$ est fermé dans $U$.
\item Le support de $fs$ est contenu dans l'intersection des supports de
$f \in \mathcal{O}_X(U)$ et de $s \in \mathcal{F}(U)$.
\item Le support de $s + s'$ est contenu dans la réunion des supports de
$s, s' \in \mathcal{F}(U)$.
\item Le support de $\mathcal{F}$ est la réunion des supports de toutes les
sections locales de $\mathcal{F}$.
\item Si $\varphi : \mathcal{F} \to \mathcal{G}$ est un morphisme de
$\mathcal{O}_X$-modules, alors le support de $\varphi(s)$ est contenu dans
celui de $s \in \mathcal{F}(U)$.
\end{enumerate}
\end{lemma}

\begin{proof}
En effet, si $s_x = 0$, alors $s$ est nul dans un voisinage ouvert de $x$ par
définition des germes. Il en va de même pour $f$. Nous omettons les détails.
\end{proof}

\noindent
En général, le support d'un faisceau de modules n'est pas fermé. En effet, on
peut prendre sur $\mathbf{R}$ (muni de sa topologie archimédienne usuelle) le
faisceau abélien somme directe d'une infinité de faisceaux gratte-ciel non nuls,
chacun supporté en un seul point $p_i$ de $\mathbf{R}$. Son support est alors
l'ensemble des points $p_i$, qui peut ne pas être fermé.

\medskip\noindent
Un autre exemple s'obtient en considérant l'immersion ouverte
$j : U = (0 , \infty) \to \mathbf{R} = X$ et le faisceau abélien
$j_!\underline{\mathbf{Z}}_U$. D'après Faisceaux, Section
\ref{sheaves-section-open-immersions}, le support de ce faisceau est exactement
$U$.

\begin{lemma}
\label{lemma-support-sheaf-rings-closed}
Soit $X$ un espace topologique. Le support d'un faisceau d'anneaux est fermé.
\end{lemma}

\begin{proof}
Cela résulte de ce que, selon nos conventions, un anneau est $0$ si et seulement
si $1 = 0$ ; le support d'un faisceau d'anneaux est donc le support de la
section unité.
\end{proof}






\section{Immersions fermées et faisceaux abéliens}
\label{section-closed-immersions}

\noindent
Rappelons que nous considérons un faisceau abélien sur un espace topologique
$X$ comme un faisceau de $\underline{\mathbf{Z}}_X$-modules. Nous pouvons donc
appliquer aux faisceaux abéliens tous les résultats et définitions concernant
les faisceaux de modules.

\begin{lemma}
\label{lemma-i-star-exact}
Soit $X$ un espace topologique et soit $Z \subset X$ une partie fermée.
Notons $i : Z \to X$ l'inclusion. Le foncteur
$$
i_* : \textit{Ab}(Z) \longrightarrow \textit{Ab}(X)
$$
est exact et pleinement fidèle ; son image essentielle est exactement
constituée des faisceaux abéliens dont le support est contenu dans $Z$. Le
foncteur $i^{-1}$ est un inverse à gauche de $i_*$.
\end{lemma}

\begin{proof}
L'exactitude résulte de la description des germes dans Faisceaux, Lemme
\ref{sheaves-lemma-stalks-closed-pushforward}, et du Lemme
\ref{lemma-abelian}. Le reste a été démontré dans Faisceaux, Lemme
\ref{sheaves-lemma-equivalence-categories-closed-abelian}.
\end{proof}

\noindent
Soit $\mathcal{F}$ un faisceau abélien sur l'espace topologique $X$. Étant
donnée une partie fermée $Z$, il existe un sous-faisceau abélien canonique de
$\mathcal{F}$ formé exactement des sections dont le support est contenu dans
$Z$. Voici l'énoncé précis.

\begin{remark}
\label{remark-sections-support-in-closed}
Soit $X$ un espace topologique, soit $Z \subset X$ une partie fermée et soit
$\mathcal{F}$ un faisceau abélien sur $X$. Pour tout ouvert $U \subset X$,
posons
$$
\mathcal{H}_Z(\mathcal{F})(U) =
\{s \in \mathcal{F}(U) \mid
\text{ le support de }s\text{ est contenu dans }Z \cap U\}.
$$
Alors $\mathcal{H}_Z(\mathcal{F})$ est un sous-faisceau abélien de
$\mathcal{F}$. C'est le plus grand sous-faisceau abélien de $\mathcal{F}$ dont
le support est contenu dans $Z$. D'après le Lemme \ref{lemma-i-star-exact},
nous pouvons (et allons) considérer $\mathcal{H}_Z(\mathcal{F})$ comme un
faisceau abélien sur $Z$. Nous obtenons ainsi un foncteur exact à gauche
$$
\textit{Ab}(X) \longrightarrow \textit{Ab}(Z),\quad
\mathcal{F} \longmapsto \mathcal{H}_Z(\mathcal{F})
\text{ considéré comme faisceau abélien sur }Z
$$
Toutes les assertions précédentes découlent directement du Lemme
\ref{lemma-support-section-closed}.
\end{remark}

\noindent
C'est l'occasion de montrer que le foncteur $i_*$ possède un adjoint à droite
sur les faisceaux abéliens.

\begin{lemma}
\label{lemma-i-star-right-adjoint}
Soit $i : Z \to X$ l'inclusion d'une partie fermée dans l'espace topologique
$X$. Le foncteur $\textit{Ab}(X) \to \textit{Ab}(Z)$,
$\mathcal{F} \mapsto \mathcal{H}_Z(\mathcal{F})$ de la Remarque
\ref{remark-sections-support-in-closed} est adjoint à droite de
$i_* : \textit{Ab}(Z) \to \textit{Ab}(X)$. En particulier, $i_*$ commute aux
limites inductives arbitraires.
\end{lemma}

\begin{proof}
Il faut montrer que, pour tout faisceau abélien $\mathcal{F}$ sur $X$ et tout
faisceau abélien $\mathcal{G}$ sur $Z$, nous avons
$$
\Hom_{\textit{Ab}(X)}(i_*\mathcal{G}, \mathcal{F}) =
\Hom_{\textit{Ab}(Z)}(\mathcal{G}, \mathcal{H}_Z(\mathcal{F}))
$$
C'est clair puisque toute section de $i_*\mathcal{G}$ a son support dans $Z$.
Nous omettons les détails.
\end{proof}

\begin{remark}
\label{remark-i-star-right-adjoint}
Dans Faisceaux, Remarque \ref{sheaves-remark-i-star-not-exact}, nous avons
montré que $i_*$, comme foncteur entre catégories de faisceaux d'ensembles,
n'admet pas d'adjoint à droite, simplement parce qu'il n'est pas exact.
Toutefois, cette assertion est presque vraie : le foncteur $i_*$ est en effet
exact sur les faisceaux d'ensembles pointés, on peut définir pour ceux-ci les
sections à support dans $Z$, et $\mathcal{H}_Z$ a un sens et est adjoint à
droite de $i_*$.
\end{remark}









\section{Une suite exacte canonique}
\label{section-canonical-exact-sequence}

\noindent
Nous consacrons une section à cette suite exacte.

\begin{lemma}
\label{lemma-canonical-exact-sequence}
Soit $X$ un espace topologique. Soit $U \subset X$ une partie ouverte de
complémentaire $Z \subset X$. Notons $j : U \to X$ l'immersion ouverte et
$i : Z \to X$ l'immersion fermée. Pour tout faisceau de groupes abéliens
$\mathcal{F}$ sur $X$, les morphismes d'adjonction
$j_{!}j^{-1}\mathcal{F} \to \mathcal{F}$ et
$\mathcal{F} \to i_*i^{-1}\mathcal{F}$ donnent une suite exacte courte
$$
0 \to j_{!}j^{-1}\mathcal{F} \to \mathcal{F} \to i_*i^{-1}\mathcal{F} \to 0
$$
de faisceaux de groupes abéliens. Pour tout morphisme
$\varphi : \mathcal{F} \to \mathcal{G}$ de faisceaux abéliens sur $X$, nous
obtenons un morphisme de suites exactes courtes
$$
\xymatrix{
0 \ar[r] &
j_{!}j^{-1}\mathcal{F} \ar[r] \ar[d] &
\mathcal{F} \ar[r] \ar[d] &
i_*i^{-1}\mathcal{F} \ar[r] \ar[d] &
0 \\
0 \ar[r] &
j_{!}j^{-1}\mathcal{G} \ar[r] &
\mathcal{G} \ar[r] &
i_*i^{-1}\mathcal{G} \ar[r] &
0
}
$$
\end{lemma}

\begin{proof}
La fonctorialité de la suite exacte courte résulte immédiatement de la
naturalité des morphismes d'adjonction. Nous pouvons vérifier l'exactitude sur
les germes (Lemme \ref{lemma-abelian}). Pour une description des germes en
question, voir Faisceaux, Lemmes \ref{sheaves-lemma-j-shriek-abelian} et
\ref{sheaves-lemma-stalks-closed-pushforward}.
\end{proof}








\section{Modules localement engendrés par des sections}
\label{section-locally-generated}

\noindent
Soit $(X, \mathcal{O}_X)$ un espace annelé. Dans cette section et la suivante,
nous restreindrons souvent les faisceaux aux sous-espaces ouverts
$U \subset X$, voir Faisceaux, Section
\ref{sheaves-section-open-immersions}. En particulier, nous noterons souvent
le sous-espace ouvert $(U, \mathcal{O}_U)$ plutôt que d'employer la notation
plus correcte $(U, \mathcal{O}_X|_U)$, voir Faisceaux, Définition
\ref{sheaves-definition-restriction}.

\medskip\noindent
Considérons l'immersion ouverte
$j : U = (0 , \infty) \to \mathbf{R} = X$ et le faisceau abélien
$j_!\underline{\mathbf{Z}}_U$. D'après Faisceaux, Section
\ref{sheaves-section-open-immersions}, le germe de
$j_!\underline{\mathbf{Z}}_U$ en $x = 0$ est $0$. En fait, les sections de ce
faisceau sur tout intervalle ouvert contenant $0$ sont $0$. Il n'existe donc
aucun voisinage ouvert du point $0$ sur lequel le faisceau puisse être engendré
par des sections.

\begin{definition}
\label{definition-locally-generated}
Soit $(X, \mathcal{O}_X)$ un espace annelé. Soit $\mathcal{F}$ un faisceau de
$\mathcal{O}_X$-modules. Nous disons que $\mathcal{F}$ est
{\it localement engendré par des sections} si, pour tout $x \in X$, il existe
un voisinage ouvert $U$ de $x$ tel que $\mathcal{F}|_U$ soit engendré par ses
sections globales comme faisceau de $\mathcal{O}_U$-modules.
\end{definition}

\noindent
Autrement dit, il existe un ensemble $I$ et, pour chaque $i$, une section
$s_i \in \mathcal{F}(U)$ telle que l'application associée
$$
\bigoplus\nolimits_{i \in I} \mathcal{O}_U
\longrightarrow
\mathcal{F}|_U
$$
soit surjective.

\begin{lemma}
\label{lemma-pullback-locally-generated}
Soit $f : (X, \mathcal{O}_X) \to (Y, \mathcal{O}_Y)$ un morphisme d'espaces
annelés. L'image inverse $f^*\mathcal{G}$ est localement engendrée par des
sections si $\mathcal{G}$ l'est.
\end{lemma}

\begin{proof}
Étant donné un sous-espace ouvert $V$ de $Y$, nous pouvons considérer le
diagramme commutatif d'espaces annelés
$$
\xymatrix{
(f^{-1}V, \mathcal{O}_{f^{-1}V}) \ar[r]_{j'} \ar[d]_{f'} &
(X, \mathcal{O}_X) \ar[d]^f \\
(V, \mathcal{O}_V) \ar[r]^j &
(Y, \mathcal{O}_Y)
}
$$
Nous savons que
$f^*\mathcal{G}|_{f^{-1}V} \cong (f')^*(\mathcal{G}|_V)$, voir Faisceaux,
Lemme \ref{sheaves-lemma-push-pull-composition-modules}. Nous pouvons donc
supposer que $\mathcal{G}$ est engendré par ses sections globales.

\medskip\noindent
Nous avons vu que $f^*$ commute à toutes les limites inductives et est exact à
droite, voir le Lemme \ref{lemma-exactness-pushforward-pullback}. Ainsi, si
nous avons une surjection
$$
\bigoplus\nolimits_{i \in I}
\mathcal{O}_Y
\to
\mathcal{G}
\to
0
$$
alors, en appliquant $f^*$, nous obtenons la surjection
$$
\bigoplus\nolimits_{i \in I}
\mathcal{O}_X
\to
f^*\mathcal{G}
\to
0.
$$
Cela démontre le lemme.
\end{proof}












\section{Modules de type fini}
\label{section-finite-type}

\begin{definition}
\label{definition-finite-type}
Soit $(X, \mathcal{O}_X)$ un espace annelé. Soit $\mathcal{F}$ un faisceau de
$\mathcal{O}_X$-modules. Nous disons que $\mathcal{F}$ est de
{\it type fini} si, pour tout $x \in X$, il existe un voisinage ouvert $U$ tel
que $\mathcal{F}|_U$ soit engendré par un nombre fini de sections.
\end{definition}

\begin{lemma}
\label{lemma-pullback-finite-type}
Soit $f : (X, \mathcal{O}_X) \to (Y, \mathcal{O}_Y)$ un morphisme d'espaces
annelés. L'image inverse $f^*\mathcal{G}$ d'un $\mathcal{O}_Y$-module de type
fini est un $\mathcal{O}_X$-module de type fini.
\end{lemma}

\begin{proof}
En raisonnant comme dans la démonstration du Lemme
\ref{lemma-pullback-locally-generated}, nous pouvons supposer que
$\mathcal{G}$ est engendré par un nombre fini de sections globales. Nous avons
vu que $f^*$ commute à toutes les limites inductives et est exact à droite,
voir le Lemme \ref{lemma-exactness-pushforward-pullback}. Ainsi, si nous avons
une surjection
$$
\bigoplus\nolimits_{i = 1, \ldots, n}
\mathcal{O}_Y
\to
\mathcal{G}
\to
0
$$
alors, en appliquant $f^*$, nous obtenons la surjection
$$
\bigoplus\nolimits_{i = 1, \ldots, n}
\mathcal{O}_X
\to
f^*\mathcal{G}
\to
0.
$$
Cela démontre le lemme.
\end{proof}

\begin{lemma}
\label{lemma-extension-finite-type}
Soit $X$ un espace annelé. L'image d'un morphisme de
$\mathcal{O}_X$-modules de type fini est de type fini. Soit
$0 \to \mathcal{F}_1 \to \mathcal{F}_2 \to \mathcal{F}_3 \to 0$
une suite exacte courte de $\mathcal{O}_X$-modules. Si $\mathcal{F}_1$ et
$\mathcal{F}_3$ sont de type fini, alors $\mathcal{F}_2$ l'est aussi.
\end{lemma}

\begin{proof}
L'assertion sur les images est triviale. Celle sur les suites exactes courtes
résulte de ce que les sections de $\mathcal{F}_3$ se relèvent localement en
des sections de $\mathcal{F}_2$, ainsi que du résultat correspondant dans la
catégorie des modules sur un anneau (appliqué, par exemple, aux germes).
\end{proof}

\begin{lemma}
\label{lemma-finite-type-surjective-on-stalk}
Soit $X$ un espace annelé. Soit
$\varphi : \mathcal{G} \to \mathcal{F}$ un homomorphisme de
$\mathcal{O}_X$-modules. Soit $x \in X$. Supposons $\mathcal{F}$ de type fini
et l'application sur les germes
$\varphi_x : \mathcal{G}_x \to \mathcal{F}_x$ surjective. Il existe alors un
voisinage ouvert $x \in U \subset X$ tel que $\varphi|_U$ soit surjective.
\end{lemma}

\begin{proof}
Choisissons un voisinage ouvert $U \subset X$ de $x$ tel que $\mathcal{F}$
soit engendré par $s_1, \ldots, s_n \in \mathcal{F}(U)$ sur $U$. Par
l'hypothèse de surjectivité de $\varphi_x$, après avoir rétréci $U$, nous
pouvons supposer que $s_i = \varphi(t_i)$ pour certains
$t_i \in \mathcal{G}(U)$. Alors $U$ convient.
\end{proof}

\begin{lemma}
\label{lemma-finite-type-stalk-zero}
Soit $X$ un espace annelé. Soit $\mathcal{F}$ un
$\mathcal{O}_X$-module et soit $x \in X$. Supposons $\mathcal{F}$ de type
fini et $\mathcal{F}_x = 0$. Il existe alors un voisinage ouvert
$x \in U \subset X$ tel que $\mathcal{F}|_U$ soit nul.
\end{lemma}

\begin{proof}
C'est un cas particulier du Lemme
\ref{lemma-finite-type-surjective-on-stalk}, appliqué au morphisme
$0 \to \mathcal{F}$.
\end{proof}

\begin{lemma}
\label{lemma-support-finite-type-closed}
\begin{slogan}
Sur tout espace annelé, les faisceaux de modules de type fini ont un support
fermé.
\end{slogan}
Soit $(X, \mathcal{O}_X)$ un espace annelé. Soit $\mathcal{F}$ un faisceau de
$\mathcal{O}_X$-modules. Si $\mathcal{F}$ est de type fini, le support de
$\mathcal{F}$ est fermé.
\end{lemma}

\begin{proof}
C'est une reformulation du Lemme \ref{lemma-finite-type-stalk-zero}.
\end{proof}

\begin{lemma}
\label{lemma-finite-type-quasi-compact-colimit}
Soit $X$ un espace annelé. Soit $I$ un ensemble préordonné et soit
$(\mathcal{F}_i, f_{ii'})$ un système indexé par $I$ et formé de faisceaux de
$\mathcal{O}_X$-modules (voir Catégories, Section
\ref{categories-section-posets-limits}). Soit
$\mathcal{F} = \colim \mathcal{F}_i$ sa limite inductive. Supposons que
(a) $I$ est filtrant, (b) $\mathcal{F}$ est un $\mathcal{O}_X$-module de type
fini, et (c) $X$ est quasi-compact. Il existe alors un $i$ tel que
$\mathcal{F}_i \to \mathcal{F}$ soit surjective. Si les applications de
transition $f_{ii'}$ sont injectives, nous en déduisons que
$\mathcal{F} = \mathcal{F}_i$ pour un certain $i \in I$.
\end{lemma}

\begin{proof}
Soit $x \in X$. Il existe un voisinage ouvert $U \subset X$ de $x$ et un
nombre fini de sections $s_j \in \mathcal{F}(U)$, $j = 1, \ldots, m$, tels
que $s_1, \ldots, s_m$ engendrent $\mathcal{F}$ comme
$\mathcal{O}_U$-module. Après avoir éventuellement rétréci $U$ en un voisinage
ouvert plus petit de $x$, nous pouvons supposer que chaque $s_j$ provient
d'une section de $\mathcal{F}_i$ pour un certain $i \in I$. Puisque $X$ est
quasi-compact, nous pouvons donc trouver un recouvrement ouvert fini
$X = \bigcup_{j = 1, \ldots, m} U_j$ et, pour chaque $j$, un indice $i_j$
ainsi qu'un nombre fini de sections
$s_{jl} \in \mathcal{F}_{i_j}(U_j)$ dont les images engendrent la restriction
de $\mathcal{F}$ à $U_j$. Le lemme vaut manifestement pour tout indice
$i \in I$ qui est $\geq$ à tous les $i_j$.
\end{proof}

\begin{lemma}
\label{lemma-set-isomorphism-classes-finite-type-modules}
Soit $X$ un espace annelé. Il existe un ensemble de
$\mathcal{O}_X$-modules de type fini $\{\mathcal{F}_i\}_{i \in I}$ tel que
chaque $\mathcal{O}_X$-module de type fini sur $X$ soit isomorphe à exactement
un des $\mathcal{F}_i$.
\end{lemma}

\begin{proof}
Pour chaque recouvrement ouvert $\mathcal{U} : X = \bigcup U_j$, considérons
les faisceaux de $\mathcal{O}_X$-modules $\mathcal{F}$ tels que chaque
restriction $\mathcal{F}|_{U_j}$ soit un quotient de
$\mathcal{O}_{U_j}^{\oplus r_j}$ pour un certain $r_j \geq 0$. Ils sont
paramétrés par des sous-faisceaux
$\mathcal{K}_j \subset \mathcal{O}_{U_j}^{\oplus r_j}$ et des données de
recollement
$$
\varphi_{jj'} :
\mathcal{O}_{U_j \cap U_{j'}}^{\oplus r_j}/
(\mathcal{K}_j|_{U_j \cap U_{j'}})
\longrightarrow
\mathcal{O}_{U_j \cap U_{j'}}^{\oplus r_{j'}}/
(\mathcal{K}_{j'}|_{U_j \cap U_{j'}})
$$
voir Faisceaux, Section \ref{sheaves-section-glueing-sheaves}. Remarquons que
la collection de toutes les données de recollement forme un ensemble. La
collection de tous les recouvrements
$\mathcal{U} : X = \bigcup_{j \in J} U_i$ pour lesquels
$J \to \mathcal{P}(X)$, $j \mapsto U_j$, est injective forme également un
ensemble. Par conséquent, la collection de tous les faisceaux de
$\mathcal{O}_X$-modules obtenus par recollement de quotients comme ci-dessus
forme un ensemble $\mathcal{I}$. Par définition, tout
$\mathcal{O}_X$-module de type fini est isomorphe à un élément de
$\mathcal{I}$. En choisissant un élément dans chaque classe d'isomorphisme de
$\mathcal{I}$, on obtient l'ensemble de faisceaux voulu (ce qui utilise
l'axiome du choix).
\end{proof}













\section{Modules quasi-cohérents}
\label{section-quasi-coherent}

\noindent
Dans cette section, nous introduisons une notion abstraite de
$\mathcal{O}_X$-module quasi-cohérent. Cette notion est très utile en géométrie
algébrique, puisque les modules quasi-cohérents sur un schéma admettent une
bonne description sur tout ouvert affine. Nous avertissons toutefois le
lecteur que, dans le cadre général des espaces (localement) annelés, cette
notion se comporte très mal. La catégorie des faisceaux quasi-cohérents n'est
pas abélienne en général, les sommes directes infinies de faisceaux
quasi-cohérents ne sont pas quasi-cohérentes, etc.

\begin{definition}
\label{definition-quasi-coherent}
Soit $(X, \mathcal{O}_X)$ un espace annelé. Soit $\mathcal{F}$ un faisceau de
$\mathcal{O}_X$-modules. Nous disons que $\mathcal{F}$ est un
{\it faisceau quasi-cohérent de $\mathcal{O}_X$-modules} si, pour tout point
$x \in X$, il existe un voisinage ouvert $x\in U \subset X$ tel que
$\mathcal{F}|_U$ soit isomorphe au conoyau d'une application
$$
\bigoplus\nolimits_{j \in J}
\mathcal{O}_U
\longrightarrow
\bigoplus\nolimits_{i \in I}
\mathcal{O}_U
$$
La catégorie des $\mathcal{O}_X$-modules quasi-cohérents est notée
$\QCoh(\mathcal{O}_X)$.
\end{definition}

\noindent
La définition signifie que $X$ est recouvert par des ouverts $U$ tels que
$\mathcal{F}|_U$ admette une {\it présentation} de la forme
$$
\bigoplus\nolimits_{j \in J}
\mathcal{O}_U
\longrightarrow
\bigoplus\nolimits_{i \in I}
\mathcal{O}_U
\longrightarrow
\mathcal{F}|_U
\longrightarrow
0.
$$
Ici, « présentation » signifie que la suite affichée est exacte. Autrement dit,
\begin{enumerate}
\item pour tout point $x$ de $X$, il existe un voisinage ouvert tel que
$\mathcal{F}|_U$ soit engendré par ses sections globales, et
\item pour un choix convenable de ces sections, le noyau de la surjection
associée est lui aussi engendré par ses sections globales.
\end{enumerate}

\begin{lemma}
\label{lemma-direct-sum-quasi-coherent}
Soit $(X, \mathcal{O}_X)$ un espace annelé. La somme directe de deux
$\mathcal{O}_X$-modules quasi-cohérents est un $\mathcal{O}_X$-module
quasi-cohérent.
\end{lemma}

\begin{proof}
Omis.
\end{proof}

\begin{remark}
\label{remark-infinite-direct-sum-quasi-coherent-not}
Attention : il n'est pas vrai, en général, qu'une somme directe infinie de
$\mathcal{O}_X$-modules quasi-cohérents soit quasi-cohérente. Pour des
comportements plus exotiques des modules quasi-cohérents, voir l'Exemple
\ref{example-quasi-coherent}.
\end{remark}

\begin{lemma}
\label{lemma-pullback-quasi-coherent}
Soit $f : (X, \mathcal{O}_X) \to (Y, \mathcal{O}_Y)$ un morphisme d'espaces
annelés. L'image inverse $f^*\mathcal{G}$ d'un $\mathcal{O}_Y$-module
quasi-cohérent est quasi-cohérente.
\end{lemma}

\begin{proof}
En raisonnant comme dans la démonstration du Lemme
\ref{lemma-pullback-locally-generated}, nous pouvons supposer que
$\mathcal{G}$ admet une présentation globale par des sommes directes de copies
de $\mathcal{O}_Y$. Nous avons vu que $f^*$ commute à toutes les limites
inductives et est exact à droite, voir le Lemme
\ref{lemma-exactness-pushforward-pullback}. Ainsi, si nous avons une suite
exacte
$$
\bigoplus\nolimits_{j \in J}
\mathcal{O}_Y
\longrightarrow
\bigoplus\nolimits_{i \in I}
\mathcal{O}_Y
\longrightarrow
\mathcal{G}
\longrightarrow
0
$$
alors, en appliquant $f^*$, nous obtenons la suite exacte
$$
\bigoplus\nolimits_{j \in J}
\mathcal{O}_X
\longrightarrow
\bigoplus\nolimits_{i \in I}
\mathcal{O}_X
\longrightarrow
f^*\mathcal{G}
\longrightarrow
0.
$$
Cela démontre le lemme.
\end{proof}

\noindent
Cela fournit de nombreux exemples de faisceaux quasi-cohérents.

\begin{lemma}
\label{lemma-construct-quasi-coherent-sheaves}
Soit $(X, \mathcal{O}_X)$ un espace annelé.
Soit $\alpha : R \to \Gamma(X, \mathcal{O}_X)$ un homomorphisme d'un anneau
$R$ vers l'anneau des sections globales sur $X$.
Soit $M$ un $R$-module.
Les trois constructions suivantes donnent des faisceaux de
$\mathcal{O}_X$-modules canoniquement isomorphes :
\begin{enumerate}
\item Soit $\pi : (X, \mathcal{O}_X) \longrightarrow (\{*\}, R)$
le morphisme d'espaces annelés dont $\pi : X \to \{*\}$ est l'unique
application et dont l'application associée à $\pi$ est $\pi^\sharp$,
l'application donnée
$\alpha : R \to \Gamma(X, \mathcal{O}_X)$. Posons $\mathcal{F}_1 = \pi^*M$.
\item Choisissons une présentation
$\bigoplus_{j \in J} R \to \bigoplus_{i \in I} R \to M \to 0$.
Posons
$$
\mathcal{F}_2 = \Coker\left(
\bigoplus\nolimits_{j \in J} \mathcal{O}_X
\to
\bigoplus\nolimits_{i \in I} \mathcal{O}_X
\right).
$$
Ici, l'application sur le facteur $\mathcal{O}_X$ correspondant à $j \in J$
est donnée par la section $\sum_i \alpha(r_{ij})$, où les $r_{ij}$
sont les coefficients matriciels de l'application dans la présentation de $M$.
\item Posons $\mathcal{F}_3$ égal au faisceau associé au préfaisceau
$U \mapsto \mathcal{O}_X(U) \otimes_R M$, où l'application
$R \to \mathcal{O}_X(U)$ est la composée de $\alpha$ et de
l'application de restriction $\mathcal{O}_X(X) \to \mathcal{O}_X(U)$.
\end{enumerate}
Cette construction possède les propriétés suivantes :
\begin{enumerate}
\item Le faisceau de $\mathcal{O}_X$-modules ainsi obtenu
$\mathcal{F}_M = \mathcal{F}_1 = \mathcal{F}_2 = \mathcal{F}_3$
est quasi-cohérent.
\item La construction définit un foncteur de la catégorie des $R$-modules
vers la catégorie des faisceaux quasi-cohérents sur $X$, qui commute aux
limites inductives arbitraires.
\item Pour tout $x \in X$, nous avons
$\mathcal{F}_{M, x} = \mathcal{O}_{X, x} \otimes_R M$
fonctoriellement en $M$.
\item Pour tout $\mathcal{O}_X$-module
$\mathcal{G}$, nous avons
$$
\Mor_{\mathcal{O}_X}(\mathcal{F}_M, \mathcal{G})
=
\Hom_R(M, \Gamma(X, \mathcal{G}))
$$
où la structure de $R$-module sur $\Gamma(X, \mathcal{G})$
provient de la structure de $\Gamma(X, \mathcal{O}_X)$-module via
$\alpha$.
\end{enumerate}
\end{lemma}

\begin{proof}
L'isomorphisme entre $\mathcal{F}_1$ et $\mathcal{F}_3$
provient du fait que $\pi^*$ est défini comme la faisceautisation
du préfaisceau de (3), voir Faisceaux, Section
\ref{sheaves-section-ringed-spaces-functoriality-modules}.
L'isomorphisme entre les constructions (2) et (1) provient
du fait que le foncteur $\pi^*$ est exact à droite, de sorte que
$\pi^*(\bigoplus_{j \in J} R) \to \pi^*(\bigoplus_{i \in I} R) \to
\pi^*M \to 0$ est exacte, que $\pi^*$ commute aux sommes directes
arbitraires, voir le Lemme \ref{lemma-exactness-pushforward-pullback},
et enfin du fait que $\pi^*(R) = \mathcal{O}_X$.

\medskip\noindent
L'assertion (1) est claire d'après la construction (2).
L'assertion (2) est claire puisque $\pi^*$ possède ces propriétés.
L'assertion (3) résulte de la description des germes des faisceaux images
inverses, voir Faisceaux, Lemme
\ref{sheaves-lemma-stalk-pullback-modules}.
L'assertion (4) résulte de l'adjonction entre $\pi_*$ et
$\pi^*$.
\end{proof}

\begin{definition}
\label{definition-sheaf-associated}
Dans la situation du Lemme \ref{lemma-construct-quasi-coherent-sheaves},
nous disons que $\mathcal{F}_M$ est le {\it faisceau associé au module $M$
et à l'homomorphisme d'anneaux $\alpha$}. Si $R = \Gamma(X, \mathcal{O}_X)$
et $\alpha = \text{id}_R$, nous disons simplement que $\mathcal{F}_M$ est le
{\it faisceau associé au module $M$}.
\end{definition}


\begin{lemma}
\label{lemma-restrict-quasi-coherent}
Soit $(X, \mathcal{O}_X)$ un espace annelé.
Posons $R = \Gamma(X, \mathcal{O}_X)$.
Soit $M$ un $R$-module.
Soit $\mathcal{F}_M$ le faisceau quasi-cohérent de
$\mathcal{O}_X$-modules associé à $M$.
Si $g : (Y, \mathcal{O}_Y) \to (X, \mathcal{O}_X)$
est un morphisme d'espaces annelés, alors
$g^*\mathcal{F}_M$ est le faisceau associé
au $\Gamma(Y, \mathcal{O}_Y)$-module
$\Gamma(Y, \mathcal{O}_Y) \otimes_R M$.
\end{lemma}

\begin{proof}
L'assertion résulte de la première description de
$\mathcal{F}_M$ dans le Lemme \ref{lemma-construct-quasi-coherent-sheaves}
comme $\pi^*M$, ainsi que du diagramme commutatif suivant
d'espaces annelés
$$
\xymatrix{
(Y, \mathcal{O}_Y) \ar[r]_-\pi \ar[d]_g &
(\{*\}, \Gamma(Y, \mathcal{O}_Y)) \ar[d]^{\text{induite par }g^\sharp} \\
(X, \mathcal{O}_X) \ar[r]^-\pi &
(\{*\}, \Gamma(X, \mathcal{O}_X))
}
$$
(On utilise aussi Faisceaux, Lemme
\ref{sheaves-lemma-push-pull-composition-modules}.)
\end{proof}

\begin{lemma}
\label{lemma-quasi-coherent-module}
Soit $(X, \mathcal{O}_X)$ un espace annelé.
Soit $x \in X$ un point.
Supposons que $x$ possède un système fondamental de voisinages quasi-compacts.
Considérons un $\mathcal{O}_X$-module quasi-cohérent quelconque $\mathcal{F}$.
Il existe alors un voisinage ouvert $U$ de $x$
tel que $\mathcal{F}|_U$ soit isomorphe au
faisceau de modules $\mathcal{F}_M$ sur $(U, \mathcal{O}_U)$
associé à un certain $\Gamma(U, \mathcal{O}_U)$-module $M$.
\end{lemma}

\begin{proof}
Nous pouvons d'abord remplacer $X$ par un voisinage ouvert de $x$
et supposer que $\mathcal{F}$ est isomorphe au
conoyau d'une application
$$
\Psi :
\bigoplus\nolimits_{j \in J}
\mathcal{O}_X
\longrightarrow
\bigoplus\nolimits_{i \in I}
\mathcal{O}_X.
$$
Le problème est que cette application n'est pas nécessairement donnée par
une « matrice », car le module des sections globales d'une somme directe
diffère en général de la somme directe des modules
de sections globales.

\medskip\noindent
Soit $x \in E \subset X$ un voisinage quasi-compact
de $x$ (remarquons que $E$ n'est pas nécessairement ouvert).
Soit $x \in U \subset E$ un voisinage ouvert de
$x$ contenu dans $E$.
Nous procédons ensuite comme dans la démonstration du
Lemme \ref{lemma-section-direct-sum-quasi-compact}.
Pour chaque $j \in J$, notons
$s_j \in \Gamma(X, \bigoplus\nolimits_{i \in I} \mathcal{O}_X)$
l'image de la section $1$ dans le facteur $\mathcal{O}_X$
correspondant à $j$. Il existe une famille finie d'ouverts
$U_{jk}$, $k \in K_j$, telle que $E \subset \bigcup_{k \in K_j} U_{jk}$
et telle que chaque restriction $s_j|_{U_{jk}}$
soit une somme finie $\sum_{i \in I_{jk}} f_{jki}$
avec $I_{jk} \subset I$, et $f_{jki}$ dans le facteur
$\mathcal{O}_X$ correspondant à $i \in I$. Posons
$I_j = \bigcup_{k \in K_j} I_{jk}$. C'est un ensemble fini.
Puisque $U \subset E \subset \bigcup_{k \in K_j} U_{jk}$,
la section $s_j|_U$ est une section de la somme directe finie
$\bigoplus_{i \in I_j} \mathcal{O}_X$.
D'après le Lemme \ref{lemma-limits-colimits},
nous voyons qu'en fait $s_j|_U$ est une somme
$\sum_{i \in I_j} f_{ij}$ et que
$f_{ij} \in \mathcal{O}_X(U) = \Gamma(U, \mathcal{O}_U)$.

\medskip\noindent
Nous pouvons alors définir un module $M$ comme le conoyau de l'application
$$
\bigoplus\nolimits_{j \in J}
\Gamma(U, \mathcal{O}_U)
\longrightarrow
\bigoplus\nolimits_{i \in I}
\Gamma(U, \mathcal{O}_U)
$$
dont la matrice est $(f_{ij})$. D'après la construction (2) du
Lemme \ref{lemma-construct-quasi-coherent-sheaves}, nous voyons que
$\mathcal{F}_M$ possède la même présentation que $\mathcal{F}|_U$
et donc que $\mathcal{F}_M \cong \mathcal{F}|_U$.
\end{proof}

\begin{example}
\label{example-quasi-coherent}
Soit $X$ l'espace formé d'une infinité dénombrable de copies
$L_1, L_2, L_3, \ldots$ de la droite réelle, toutes recollées en $0$ ;
un système fondamental de voisinages de $0$ est la famille
$\{U_n\}_{n \in \mathbf{N}}$, où $U_n \cap L_i = (-1/n, 1/n)$.
Soit $\mathcal{O}_X$ le faisceau des fonctions continues à valeurs réelles.
Soit $f : \mathbf{R} \to \mathbf{R}$ une fonction continue
identiquement nulle sur $(-1, 1)$ et identiquement égale à $1$
sur $(-\infty, -2) \cup (2, \infty)$. Notons $f_n$ la fonction continue
sur $X$ qui vaut $x \mapsto f(nx)$ sur chaque
$L_j = \mathbf{R}$. Soit $1_{L_j}$ la fonction caractéristique
de $L_j$. Considérons l'application
$$
\bigoplus\nolimits_{j \in \mathbf{N}}
\mathcal{O}_X
\longrightarrow
\bigoplus\nolimits_{j, i \in \mathbf{N}}
\mathcal{O}_X, \quad
e_j \longmapsto \sum\nolimits_{i \in \mathbf{N}} f_j 1_{L_i} e_{ij}
$$
avec les notations évidentes. Cela a un sens, car cette somme est localement
finie puisque $f_j$ est nulle dans un voisinage de $0$. Sur $U_n$,
l'image de $e_j$, pour $j > 2n$, n'est pas une combinaison linéaire finie
$\sum g_{ij} e_{ij}$ avec les $g_{ij}$ continues. Il n'existe donc
aucun voisinage de $0 \in X$ sur lequel l'application affichée soit donnée par
une « matrice » comme dans la démonstration du Lemme
\ref{lemma-quasi-coherent-module} ci-dessus.

\medskip\noindent
Remarquons que $\bigoplus\nolimits_{j \in \mathbf{N}} \mathcal{O}_X$
est le faisceau associé au module libre de base $e_j$,
et de même pour l'autre somme directe. Nous voyons donc qu'en général,
même localement sur $X$, un morphisme entre faisceaux associés à des modules
ne provient pas d'un morphisme de modules.
Il devrait de même exister un exemple d'espace annelé $X$
et de $\mathcal{O}_X$-module quasi-cohérent $\mathcal{F}$
tel que $\mathcal{F}$ ne soit pas localement de la forme $\mathcal{F}_M$.
(Merci de nous écrire si vous en trouvez un.)
En outre, il devrait exister des exemples d'espaces localement compacts
$X$ et d'applications $\mathcal{F}_M \to \mathcal{F}_N$ qui, elles non plus,
ne proviennent pas localement d'applications de modules (la démonstration du
Lemme \ref{lemma-quasi-coherent-module} montre que cela ne peut se produire
si $N$ est libre).
\end{example}












\section{Modules de présentation finie}
\label{section-finite-presentation}

\noindent
Voici la définition.

\begin{definition}
\label{definition-finite-presentation}
Soit $(X, \mathcal{O}_X)$ un espace annelé.
Soit $\mathcal{F}$ un faisceau de $\mathcal{O}_X$-modules.
Nous disons que $\mathcal{F}$ est de {\it présentation finie}
si, pour tout point $x \in X$, il existe un voisinage ouvert
$x\in U \subset X$ et des entiers $n, m \in \mathbf{N}$ tels que
$\mathcal{F}|_U$ soit isomorphe au conoyau d'une application
$$
\bigoplus\nolimits_{j = 1, \ldots, m}
\mathcal{O}_U
\longrightarrow
\bigoplus\nolimits_{i = 1, \ldots, n}
\mathcal{O}_U
$$
\end{definition}

\noindent
Cela signifie que $X$ est recouvert par des ouverts $U$
tels que $\mathcal{F}|_U$ possède une {\it présentation}
de la forme
$$
\bigoplus\nolimits_{j = 1, \ldots, m}
\mathcal{O}_U
\longrightarrow
\bigoplus\nolimits_{i = 1, \ldots, n}
\mathcal{O}_U
\to
\mathcal{F}|_U
\to
0.
$$
Ici, « présentation » signifie que la suite affichée
est exacte. Autrement dit,
\begin{enumerate}
\item pour tout point $x$ de $X$, il existe
un voisinage ouvert tel que $\mathcal{F}|_U$
soit engendré par un nombre fini de sections globales, et
\item pour un choix convenable de ces sections,
le noyau de la surjection associée est lui aussi
engendré par un nombre fini de sections globales.
\end{enumerate}

\begin{lemma}
\label{lemma-finite-presentation-quasi-coherent}
Soit $(X, \mathcal{O}_X)$ un espace annelé.
Tout $\mathcal{O}_X$-module de présentation finie
est quasi-cohérent.
\end{lemma}

\begin{proof}
Cela découle immédiatement des définitions.
\end{proof}

\begin{lemma}
\label{lemma-cokernel-finite-finite-presentation}
Soit $(X,\mathcal{O}_X)$ un espace annelé. Soit $\mathcal{F}$ un
$\mathcal{O}_X$-module de présentation finie. Soit
$\varphi : \mathcal{G} \to \mathcal{F}$ un morphisme de
$\mathcal{O}_X$-modules. Si $\mathcal{G}$ est de type fini, alors
$\Coker(\varphi)$ est de présentation finie.
\end{lemma}

\begin{proof}
Localement sur $X$, nous pouvons écrire
$\mathcal{F} = \mathcal{O}_X^{\oplus n} / \mathcal{M}$, où
$\mathcal{M} \subset \mathcal{O}_X^{\oplus n}$ est un sous-module de
$\mathcal{O}_X$ de type fini. Par conséquent,
$\Im(\varphi) = \mathcal{N} / \mathcal{M}$, où
$\mathcal{N} \subset \mathcal{O}_X^{\oplus n}$ est un sous-module de
$\mathcal{O}_X$ contenant $\mathcal{M}$.
Le $\mathcal{O}_X$-module $\Im(\varphi)$ est de type fini
parce que $\mathcal{G} \to \Im(\varphi)$ est surjective et que
$\mathcal{G}$ est de type fini. D'après le Lemme
\ref{lemma-extension-finite-type}, $\mathcal{N}$ est de type fini. Ainsi,
$\Coker(\varphi) = \mathcal{O}_X^{\oplus n} / \mathcal{N}$
est de présentation finie.
\end{proof}

\begin{lemma}
\label{lemma-kernel-surjection-finite-free-onto-finite-presentation}
Soit $(X, \mathcal{O}_X)$ un espace annelé.
Soit $\mathcal{F}$ un $\mathcal{O}_X$-module de présentation finie.
\begin{enumerate}
\item Si $\psi : \mathcal{O}_X^{\oplus r} \to \mathcal{F}$ est surjective,
alors $\Ker(\psi)$ est de type fini.
\item Si $\theta : \mathcal{G} \to \mathcal{F}$ est surjective et si
$\mathcal{G}$ est de type fini, alors $\Ker(\theta)$ est de type fini.
\end{enumerate}
\end{lemma}

\begin{proof}
Démonstration de (1).
Soit $x \in X$. Choisissons un voisinage ouvert $U \subset X$ de $x$
sur lequel il existe une présentation
$$
\mathcal{O}_U^{\oplus m}
\xrightarrow{\chi}
\mathcal{O}_U^{\oplus n}
\xrightarrow{\varphi}
\mathcal{F}|_U
\to
0.
$$
Soit $e_k$ la section qui engendre le $k$-ième facteur de
$\mathcal{O}_X^{\oplus r}$. Pour tout $k = 1, \ldots, r$,
nous pouvons, après avoir rétréci $U$ en un petit voisinage de $x$,
relever $\psi(e_k)$ en une section $\tilde e_k$ de
$\mathcal{O}_U^{\oplus n}$ sur $U$. Cela définit un morphisme
de faisceaux $\alpha : \mathcal{O}_U^{\oplus r} \to \mathcal{O}_U^{\oplus n}$
tel que $\varphi \circ \alpha = \psi$.
De même, après avoir rétréci $U$, nous pouvons trouver un morphisme
$\beta : \mathcal{O}_U^{\oplus n} \to \mathcal{O}_U^{\oplus r}$
tel que $\psi \circ \beta = \varphi$. Alors l'application
$$
\mathcal{O}_U^{\oplus m} \oplus
\mathcal{O}_U^{\oplus r}
\xrightarrow{\beta \circ \chi, 1 - \beta \circ \alpha}
\mathcal{O}_U^{\oplus r}
$$
est une surjection sur le noyau de $\psi$.

\medskip\noindent
Pour démontrer (2), nous pouvons choisir localement une surjection
$\eta : \mathcal{O}_X^{\oplus r} \to \mathcal{G}$.
D'après (1), $\Ker(\theta \circ \eta)$ est de type fini.
Comme $\Ker(\theta) = \eta(\Ker(\theta \circ \eta))$, le résultat s'ensuit.
\end{proof}

\begin{lemma}
\label{lemma-pullback-finite-presentation}
Soit $f : (X, \mathcal{O}_X) \to (Y, \mathcal{O}_Y)$
un morphisme d'espaces annelés.
L'image inverse $f^*\mathcal{G}$ d'un module de présentation finie
est de présentation finie.
\end{lemma}

\begin{proof}
La démonstration est exactement la même que celle du Lemme
\ref{lemma-pullback-quasi-coherent}, mais avec des ensembles d'indices finis.
\end{proof}

\begin{lemma}
\label{lemma-quasi-coherent-limit-finite-presentation}
Soit $(X, \mathcal{O}_X)$ un espace annelé.
Posons $R = \Gamma(X, \mathcal{O}_X)$.
Soit $M$ un $R$-module.
Le $\mathcal{O}_X$-module $\mathcal{F}_M$ associé à $M$
est une limite inductive filtrante de $\mathcal{O}_X$-modules de présentation
finie.
\end{lemma}

\begin{proof}
Cela résulte immédiatement du
Lemme \ref{lemma-construct-quasi-coherent-sheaves}
et du fait que tout module est une limite inductive filtrante
de modules de présentation finie, voir
Algèbre, Lemme \ref{algebra-lemma-module-colimit-fp}.
\end{proof}

\begin{lemma}
\label{lemma-finite-presentation-stalk-free}
Soit $(X, \mathcal{O}_X)$ un espace annelé. Soit $\mathcal{F}$ un
$\mathcal{O}_X$-module de présentation finie. Soit $x \in X$ tel que
$\mathcal{F}_x \cong \mathcal{O}_{X, x}^{\oplus r}$. Il existe alors
un voisinage ouvert $U$ de $x$ tel que
$\mathcal{F}|_U \cong \mathcal{O}_U^{\oplus r}$.
\end{lemma}

\begin{proof}
Choisissons $s_1, \ldots, s_r \in \mathcal{F}_x$ qui correspondent, par
l'isomorphisme, à une base de $\mathcal{O}_{X, x}^{\oplus r}$. Choisissons un
voisinage ouvert $U$ de $x$ tel que $s_i$ se relève en
$s_i \in \mathcal{F}(U)$. Après avoir rétréci $U$, nous voyons que
l'application induite $\psi : \mathcal{O}_U^{\oplus r} \to \mathcal{F}|_U$
est surjective (Lemme \ref{lemma-finite-type-surjective-on-stalk}).
D'après le Lemme
\ref{lemma-kernel-surjection-finite-free-onto-finite-presentation},
$\Ker(\psi)$ est de type fini.
Alors $\Ker(\psi)_x = 0$ implique que $\Ker(\psi)$ devient nul
après avoir encore rétréci $U$ (Lemme \ref{lemma-finite-type-stalk-zero}).
\end{proof}
\section{Modules cohérents}
\label{section-coherent}

\noindent
Une référence pour cette section est \cite{FAC}.

\medskip\noindent
La catégorie des faisceaux cohérents sur un espace annelé $X$
est un objet plus raisonnable
que la catégorie des faisceaux quasi-cohérents, en ce sens
qu'elle est au moins une sous-catégorie abélienne de $\textit{Mod}(\mathcal{O}_X)$
quel que soit $X$. En revanche, l'image réciproque d'un
module cohérent n'est « presque jamais » cohérente dans le cadre général
des espaces annelés.

\begin{definition}
\label{definition-coherent}
Soit $(X, \mathcal{O}_X)$ un espace annelé.
Soit $\mathcal{F}$ un faisceau de $\mathcal{O}_X$-modules.
On dit que $\mathcal{F}$ est un {\it $\mathcal{O}_X$-module cohérent}
si les deux conditions suivantes sont satisfaites :
\begin{enumerate}
\item $\mathcal{F}$ est de type fini, et
\item pour tout ouvert $U \subset X$ et toute famille finie
$s_i \in \mathcal{F}(U)$, $i = 1, \ldots, n$,
le noyau de l'application associée
$\bigoplus_{i = 1, \ldots, n} \mathcal{O}_U \to \mathcal{F}|_U$
est de type fini.
\end{enumerate}
La catégorie des $\mathcal{O}_X$-modules cohérents est notée
$\textit{Coh}(\mathcal{O}_X)$.
\end{definition}

\begin{lemma}
\label{lemma-coherent-finite-presentation}
Soit $(X, \mathcal{O}_X)$ un espace annelé.
Tout $\mathcal{O}_X$-module cohérent est de présentation finie
et, par conséquent, quasi-cohérent.
\end{lemma}

\begin{proof}
Soit $\mathcal{F}$ un faisceau cohérent sur $X$.
Choisissons un point $x \in X$.
D'après (1) de la définition de la cohérence, il existe un voisinage ouvert $U$
et des sections $s_i$, $i = 1, \ldots, n$, de $\mathcal{F}$ sur $U$
tels que $\Psi : \bigoplus_{i = 1, \ldots, n} \mathcal{O}_U \to \mathcal{F}$
soit surjective. D'après (2) de la définition de la cohérence, il existe
un voisinage ouvert $V$, $x \in V \subset U$, et des sections
$t_1, \ldots, t_m$ de $\bigoplus_{i = 1, \ldots, n} \mathcal{O}_V$
qui engendrent le noyau de $\Psi|_V$. On obtient alors sur $V$ la
présentation
$$
\bigoplus\nolimits_{j = 1, \ldots, m}
\mathcal{O}_V
\longrightarrow
\bigoplus\nolimits_{i = 1, \ldots, n}
\mathcal{O}_V
\to
\mathcal{F}|_V
\to
0
$$
voulue.
\end{proof}

\begin{example}
\label{example-coherent-not-Noetherian}
Supposons que $X$ soit un point. Dans ce cas, la définition
ci-dessus donne une notion pour les modules sur les anneaux.
Que signifie la définition de la cohérence ?
Elle est étroitement liée à la notion de noethérianité,
mais ne lui est pas identique : en effet, l'anneau
$R = \mathbf{C}[x_1, x_2, x_3, \ldots]$ est cohérent
comme module sur lui-même, mais n'est pas noethérien comme module
sur lui-même. Voir Algèbre, section \ref{algebra-section-coherent},
pour une discussion plus approfondie.
\end{example}

\begin{lemma}
\label{lemma-coherent-abelian}
Soit $(X, \mathcal{O}_X)$ un espace annelé.
\begin{enumerate}
\item Tout sous-faisceau de type fini d'un faisceau cohérent est cohérent.
\item Soit $\varphi : \mathcal{F} \to \mathcal{G}$
un morphisme d'un faisceau de type fini $\mathcal{F}$
vers un faisceau cohérent $\mathcal{G}$. Alors $\Ker(\varphi)$ est de type fini.
\item Soit $\varphi : \mathcal{F} \to \mathcal{G}$ un morphisme
de $\mathcal{O}_X$-modules cohérents. Alors
$\Ker(\varphi)$ et
$\Coker(\varphi)$ sont cohérents.
\item Étant donnée une suite exacte courte de $\mathcal{O}_X$-modules
$0 \to \mathcal{F}_1 \to \mathcal{F}_2 \to \mathcal{F}_3 \to 0$
si deux des trois termes sont cohérents, le troisième l'est aussi.
\item La catégorie $\textit{Coh}(\mathcal{O}_X)$ est une sous-catégorie de Serre faible
de $\textit{Mod}(\mathcal{O}_X)$. En particulier, la catégorie des
modules cohérents est abélienne et le foncteur d'inclusion
$\textit{Coh}(\mathcal{O}_X) \to \textit{Mod}(\mathcal{O}_X)$
est exact.
\end{enumerate}
\end{lemma}

\begin{proof}
La condition (2) de la définition \ref{definition-coherent}
est satisfaite par tout sous-faisceau d'un faisceau cohérent. On obtient donc (1).

\medskip\noindent
Supposons vérifiées les hypothèses de (2).
Montrons que $\Ker(\varphi)$ est de type fini. Choisissons $x \in X$.
Choisissons un voisinage ouvert $U$ de $x$ dans $X$ tel
que $\mathcal{F}|_U$ soit engendré par $s_1, \ldots, s_n$.
D'après la définition \ref{definition-coherent}, le noyau $\mathcal{K}$
de l'application induite
$\bigoplus_{i = 1}^n \mathcal{O}_U \to \mathcal{G}$,
$e_i \mapsto \varphi(s_i)$ est de type fini.
Par conséquent, $\Ker(\varphi)$, qui est l'image de la
composée
$\mathcal{K} \to \bigoplus_{i = 1}^n \mathcal{O}_U \to \mathcal{F}$
est de type fini.

\medskip\noindent
Supposons vérifiées les hypothèses de (3).
D'après (2), le noyau de $\varphi$ est de type fini et
il est donc cohérent d'après (1).

\medskip\noindent
Sous les mêmes hypothèses, montrons que $\Coker(\varphi)$ est cohérent.
Puisque $\mathcal{G}$ est de type fini, $\Coker(\varphi)$ l'est aussi.
Soit $U \subset X$ un ouvert et soient $\overline{s}_i \in \Coker(\varphi)(U)$,
$i = 1, \ldots, n$, des sections. Nous devons montrer que le noyau du
morphisme associé
$\overline{\Psi} : \bigoplus_{i = 1}^n \mathcal{O}_U \to \Coker(\varphi)$
est de type fini. Il existe un recouvrement ouvert
de $U$ tel que, sur chaque ouvert, toutes les sections $\overline{s}_i$
se relèvent en des sections $s_i$ de $\mathcal{G}$. Nous pouvons donc supposer
que tel est le cas sur $U$. Nous pouvons en outre supposer qu'il existe des
sections $t_j$, $j = 1, \ldots, m$, de $\Im(\varphi)$ sur $U$
qui engendrent $\Im(\varphi)$ sur $U$.
Soit $\Phi : \bigoplus_{j = 1}^m \mathcal{O}_U \to \Im(\varphi)$
définie à l'aide des $t_j$, et
$\Psi :
\bigoplus_{j = 1}^m \mathcal{O}_U \oplus
\bigoplus_{i = 1}^n \mathcal{O}_U \to \mathcal{G}$
définie à l'aide des $t_j$ et des $s_i$.
Considérons le diagramme commutatif suivant
$$
\xymatrix{
0 \ar[r] &
\bigoplus_{j = 1}^m \mathcal{O}_U \ar[d]_\Phi \ar[r] &
\bigoplus_{j = 1}^m \mathcal{O}_U \oplus
\bigoplus_{i = 1}^n \mathcal{O}_U \ar[d]_\Psi \ar[r] &
\bigoplus_{i = 1}^n \mathcal{O}_U \ar[d]_{\overline{\Psi}} \ar[r] &
0 \\
0 \ar[r] &
\Im(\varphi) \ar[r] &
\mathcal{G} \ar[r] &
\Coker(\varphi) \ar[r] &
0
}
$$
Le lemme du serpent donne une suite exacte
$\Ker(\Psi) \to \Ker(\overline{\Psi}) \to 0$. Puisque $\Ker(\Psi)$
est un module de type fini, on voit que $\Ker(\overline{\Psi})$ est de type fini.

\medskip\noindent
Démonstration de la partie (4).
Soit $0 \to \mathcal{F}_1 \to \mathcal{F}_2 \to \mathcal{F}_3 \to 0$
une suite exacte courte de $\mathcal{O}_X$-modules. D'après la partie
(3), il suffit
de prouver que, si $\mathcal{F}_1$ et $\mathcal{F}_3$ sont cohérents,
$\mathcal{F}_2$ l'est aussi. Le lemme \ref{lemma-extension-finite-type}
montre que $\mathcal{F}_2$ est de type fini. Soient
$s_1, \ldots, s_n$ un nombre fini de sections locales
de $\mathcal{F}_2$ définies sur un même ouvert $U$ de $X$.
Nous devons montrer que le module des relations $\mathcal{K}$
entre celles-ci est de type fini.
Considérons le diagramme commutatif suivant
$$
\xymatrix{
0 \ar[r] &
0 \ar[r] \ar[d] &
\bigoplus_{i = 1}^{n} \mathcal{O}_U \ar[r] \ar[d] &
\bigoplus_{i = 1}^{n} \mathcal{O}_U \ar[r] \ar[d] &
0 \\
0 \ar[r] &
\mathcal{F}_1 \ar[r] &
\mathcal{F}_2 \ar[r] &
\mathcal{F}_3 \ar[r] &
0
}
$$
avec les notations évidentes. Le lemme du serpent
donne une suite exacte courte
$0 \to \mathcal{K} \to \mathcal{K}_3 \to \mathcal{F}_1$
où $\mathcal{K}_3$ est le module des relations entre
les images des sections $s_i$ dans $\mathcal{F}_3$.
Puisque $\mathcal{F}_1$ est cohérent, on voit que
$\mathcal{K}$ est le noyau d'une application d'un module de type fini
vers un module cohérent, et est donc de type fini d'après (2).

\medskip\noindent
Démonstration de (5). Cela résulte de (3) et (4), qui montrent que
le lemme \ref{homology-lemma-characterize-weak-serre-subcategory} d'Homologie
s'applique.
\end{proof}

\begin{lemma}
\label{lemma-coherent-structure-sheaf}
Soit $(X, \mathcal{O}_X)$ un espace annelé.
Soit $\mathcal{F}$ un $\mathcal{O}_X$-module.
Supposons que $\mathcal{O}_X$ soit un $\mathcal{O}_X$-module cohérent.
Alors $\mathcal{F}$ est cohérent si et seulement s'il est
de présentation finie.
\end{lemma}

\begin{proof}
Omis.
\end{proof}

\begin{lemma}
\label{lemma-finite-type-to-coherent-injective-on-stalk}
Soit $X$ un espace annelé.
Soit $\varphi : \mathcal{G} \to \mathcal{F}$ un homomorphisme
de $\mathcal{O}_X$-modules.
Soit $x \in X$. Supposons $\mathcal{G}$ de type fini,
$\mathcal{F}$ cohérent et l'application sur les germes
$\varphi_x : \mathcal{G}_x \to \mathcal{F}_x$ injective.
Alors il existe un voisinage ouvert
$x \in U \subset X$ tel que $\varphi|_U$ soit injective.
\end{lemma}

\begin{proof}
Notons $\mathcal{K} \subset \mathcal{G}$ le noyau de $\varphi$.
Le lemme \ref{lemma-coherent-abelian} montre que $\mathcal{K}$ est
un $\mathcal{O}_X$-module de type fini. Notre hypothèse est que
$\mathcal{K}_x = 0$. D'après le lemme \ref{lemma-finite-type-stalk-zero},
il existe un voisinage ouvert $U$ de $x$ tel que $\mathcal{K}|_U = 0$.
Cet ouvert $U$ convient.
\end{proof}













\section{Immersions fermées d'espaces annelés}
\label{section-closed-immersion}

\noindent
Quand déclarons-nous qu'un morphisme d'espaces annelés
$i : (Z, \mathcal{O}_Z) \to (X, \mathcal{O}_X)$
est une immersion fermée ?

\medskip\noindent
Motivés par l'exemple d'une immersion fermée d'espaces topologiques normaux
(annelés par le faisceau des fonctions continues), ou de variétés différentielles
(annelées par le faisceau des fonctions différentiables), il semble naturel de
supposer au moins que :
\begin{enumerate}
\item L'application $i$ est une immersion fermée d'espaces topologiques.
\item L'application associée $\mathcal{O}_X \to i_*\mathcal{O}_Z$
est surjective. Notons son noyau $\mathcal{I}$.
\end{enumerate}
Ces conditions impliquent déjà plusieurs résultats agréables : par exemple,
nous démontrons que la catégorie des $\mathcal{O}_Z$-modules est équivalente à
la catégorie des $\mathcal{O}_X$-modules annulés par $\mathcal{I}$,
ce qui généralise le résultat sur les faisceaux abéliens de la
section \ref{section-closed-immersions}.

\medskip\noindent
Cependant, dans le projet Stacks nous choisissons la
définition qui garantit que, si $i$ est une immersion fermée
et $(X, \mathcal{O}_X)$ un schéma, alors $(Z, \mathcal{O}_Z)$
est également un schéma. De plus, dans cette situation, nous voulons que $i_*$ et $i^*$
induisent une équivalence entre la catégorie des
$\mathcal{O}_Z$-modules quasi-cohérents et la catégorie des
$\mathcal{O}_X$-modules quasi-cohérents annulés par $\mathcal{I}$.
Une condition minimale est que $i_*\mathcal{O}_Z$ soit un
faisceau quasi-cohérent de $\mathcal{O}_X$-modules.
Une bonne manière de garantir que $i_*\mathcal{O}_Z$ est un
$\mathcal{O}_X$-module quasi-cohérent consiste à supposer que
$\mathcal{I}$ est localement engendré par des sections.
On peut interpréter cette condition en disant que « $(Z, \mathcal{O}_Z)$ est
localement défini sur $(X, \mathcal{O}_X)$ en annulant certaines fonctions régulières
$f_i$, c'est-à-dire des sections locales de $\mathcal{O}_X$ ».
Ceci conduit à la définition suivante.

\begin{definition}
\label{definition-closed-immersion}
Une {\it immersion fermée d'espaces annelés}\footnote{Cette terminologie n'est
pas standard ; voir la discussion ci-dessus.} est un morphisme
$i : (Z, \mathcal{O}_Z) \to (X, \mathcal{O}_X)$
ayant les propriétés suivantes :
\begin{enumerate}
\item L'application $i$ est une immersion fermée d'espaces topologiques.
\item L'application associée $\mathcal{O}_X \to i_*\mathcal{O}_Z$
est surjective. Notons son noyau $\mathcal{I}$.
\item Le $\mathcal{O}_X$-module $\mathcal{I}$ est localement
engendré par des sections.
\end{enumerate}
\end{definition}

\noindent
En fait, cette définition ne garantit toujours pas que
l'image directe par $i_*$ d'un $\mathcal{O}_Z$-module quasi-cohérent soit un
$\mathcal{O}_X$-module quasi-cohérent. Le problème est que
l'on ne voit pas clairement comment transformer une présentation locale d'un
$\mathcal{O}_Z$-module quasi-cohérent en une présentation locale
de son image directe. Toutefois, l'énoncé suivant
est immédiat.

\begin{lemma}
\label{lemma-i-star-quasi-coherent}
Soit $i : (Z, \mathcal{O}_Z) \to (X, \mathcal{O}_X)$
une immersion fermée d'espaces annelés.
Soit $\mathcal{F}$ un $\mathcal{O}_Z$-module quasi-cohérent.
Alors $i_*\mathcal{F}$ est localement sur $X$ le conoyau d'une
application de $\mathcal{O}_X$-modules quasi-cohérents.
\end{lemma}

\begin{proof}
Ceci est vrai parce que $i_*\mathcal{O}_Z$ est quasi-cohérent
par définition. Et, localement sur $Z$, le faisceau $\mathcal{F}$
est le conoyau d'une application entre sommes directes de copies
de $\mathcal{O}_Z$. De plus, toute somme directe de copies d'un
{\it même} faisceau quasi-cohérent est quasi-cohérente.
Enfin, $i_*$ commute aux colimites arbitraires,
voir le lemme \ref{lemma-i-star-right-adjoint}. Certains détails sont omis.
\end{proof}

\begin{lemma}
\label{lemma-i-star-reflects-finite-type}
Soit $i : (Z, \mathcal{O}_Z) \to (X, \mathcal{O}_X)$ un morphisme d'espaces
annelés. Supposons que $i$ soit un homéomorphisme sur un fermé de $X$ et
que $\mathcal{O}_X \to i_*\mathcal{O}_Z$ soit surjective.
Soit $\mathcal{F}$ un $\mathcal{O}_Z$-module.
Alors $i_*\mathcal{F}$ est de type fini si et seulement si
$\mathcal{F}$ est de type fini.
\end{lemma}

\begin{proof}
Supposons que $\mathcal{F}$ soit de type fini.
Choisissons $x \in X$. Si $x \not \in Z$, alors $i_*\mathcal{F}$
est nul dans un voisinage de $x$, et donc engendré par un nombre fini d'éléments dans
un voisinage de $x$. Si $x = i(z)$, choisissons un voisinage ouvert
$z \in V \subset Z$ et des sections $s_1, \ldots, s_n \in \mathcal{F}(V)$
qui engendrent $\mathcal{F}$ sur $V$. Écrivons $V = Z \cap U$ pour un ouvert
$U \subset X$. Remarquons que $U$ est un voisinage de $x$. Les
sections $s_i$ donnent clairement des sections $s_i$ de $i_*\mathcal{F}$ sur $U$.
L'application qui en résulte
$$
\bigoplus\nolimits_{i = 1, \ldots, n} \mathcal{O}_U
\longrightarrow
i_*\mathcal{F}|_U
$$
est surjective, comme on le voit sur les germes
(on utilise ici la surjectivité de $\mathcal{O}_X \to i_*\mathcal{O}_Z$).
Par conséquent, $i_*\mathcal{F}$
est de type fini.

\medskip\noindent
Réciproquement, supposons que $i_*\mathcal{F}$ soit de type fini.
Choisissons $z \in Z$. Posons $x = i(z)$. Par hypothèse, il existe
un voisinage ouvert $U \subset X$ de $x$, et des sections
$s_1, \ldots, s_n \in (i_*\mathcal{F})(U)$ qui engendrent $i_*\mathcal{F}$
sur $U$. Posons $V = Z \cap U$. Par définition de $i_*$, les sections
$s_i$ correspondent à des sections $s_i$ de $\mathcal{F}$ sur $V$.
L'application qui en résulte
$$
\bigoplus\nolimits_{i = 1, \ldots, n} \mathcal{O}_V
\longrightarrow
\mathcal{F}|_V
$$
est surjective, comme on le voit sur les germes. Par conséquent,
$\mathcal{F}$ est de type fini.
\end{proof}

\begin{lemma}
\label{lemma-i-star-equivalence}
Soit $i : (Z, \mathcal{O}_Z) \to (X, \mathcal{O}_X)$ un morphisme
d'espaces annelés. Supposons que $i$ soit un homéomorphisme sur un fermé de $X$
et que $i^\sharp : \mathcal{O}_X \to i_*\mathcal{O}_Z$ soit surjective.
Notons $\mathcal{I} \subset \mathcal{O}_X$ le noyau de $i^\sharp$.
Le foncteur
$$
i_* :
\textit{Mod}(\mathcal{O}_Z)
\longrightarrow
\textit{Mod}(\mathcal{O}_X)
$$
est exact, pleinement fidèle, et a pour image essentielle les
$\mathcal{O}_X$-modules $\mathcal{G}$ tels que $\mathcal{I}\mathcal{G} = 0$.
\end{lemma}

\begin{proof}
Nous affirmons que, pour un $\mathcal{O}_Z$-module $\mathcal{F}$, l'application canonique
$$
i^*i_*\mathcal{F} \longrightarrow \mathcal{F}
$$
est un isomorphisme. Vérifions-le sur les germes. Soient $z \in Z$ et $x = i(z)$.
Nous avons
$$
(i^*i_*\mathcal{F})_z =
(i_*\mathcal{F})_x \otimes_{\mathcal{O}_{X, x}} \mathcal{O}_{Z, z} =
\mathcal{F}_z \otimes_{\mathcal{O}_{X, x}} \mathcal{O}_{Z, z} =
\mathcal{F}_z
$$
d'après le lemme \ref{sheaves-lemma-stalk-pullback-modules} de Faisceaux,
le fait que $\mathcal{O}_{Z, z}$ est un quotient de $\mathcal{O}_{X, x}$, et
le lemme \ref{sheaves-lemma-stalks-closed-pushforward} de Faisceaux.
Il s'ensuit que $i_*$ est pleinement fidèle.

\medskip\noindent
Soit $\mathcal{G}$ un $\mathcal{O}_X$-module tel que
$\mathcal{I}\mathcal{G} = 0$. Nous allons démontrer que l'application canonique
$$
\mathcal{G} \longrightarrow i_*i^*\mathcal{G}
$$
est un isomorphisme. Ceci prouve que $\mathcal{G} = i_*\mathcal{F}$
avec $\mathcal{F} = i^*\mathcal{G}$, ce qui achève la démonstration.
Vérifions que l'application affichée induit un isomorphisme sur les germes.
Si $x \in X$, $x \not \in i(Z)$, alors $\mathcal{G}_x = 0$
car $\mathcal{I}_x = \mathcal{O}_{X, x}$ dans ce
cas. Comme ci-dessus, $(i_*i^*\mathcal{G})_x = 0$ d'après le
lemme \ref{sheaves-lemma-stalks-closed-pushforward} de Faisceaux.
En revanche, si $x \in Z$, nous obtenons l'application
$$
\mathcal{G}_x
\longrightarrow
\mathcal{G}_x \otimes_{\mathcal{O}_{X, x}} \mathcal{O}_{Z, x}
$$
d'après les lemmes \ref{sheaves-lemma-stalk-pullback-modules} et
\ref{sheaves-lemma-stalks-closed-pushforward} de Faisceaux. Cette application est un isomorphisme
car $\mathcal{O}_{Z, x} = \mathcal{O}_{X, x}/\mathcal{I}_x$
et parce que $\mathcal{G}_x$ est annulé par $\mathcal{I}_x$ par hypothèse.
\end{proof}

\begin{remark}
\label{remark-sections-support-in-closed-modules}
Soit $(X, \mathcal{O}_X)$ un espace annelé. Soit $Z \subset X$
un fermé. Pour un $\mathcal{O}_X$-module $\mathcal{F}$, on peut
considérer le {\it sous-module des sections à support dans $Z$}, noté
$\mathcal{H}_Z(\mathcal{F})$, défini par la règle
$$
\mathcal{H}_Z(\mathcal{F})(U) =
\{s \in \mathcal{F}(U) \mid \text{Supp}(s) \subset U \cap Z\}
$$
Observons que $\mathcal{H}_Z(\mathcal{F})(U)$ est un module sur
$\mathcal{O}_X(U)$, c'est-à-dire que $\mathcal{H}_Z(\mathcal{F})$ est un
$\mathcal{O}_X$-module. Par construction, $\mathcal{H}_Z(\mathcal{F})$
est le plus grand $\mathcal{O}_X$-sous-module de $\mathcal{F}$ dont le support est
contenu dans $Z$.
En appliquant le lemme \ref{lemma-i-star-equivalence} au
morphisme d'espaces annelés $(Z, \mathcal{O}_X|_Z) \to (X, \mathcal{O}_X)$, nous
pouvons (et nous allons) considérer
$\mathcal{H}_Z(\mathcal{F})$ comme un $\mathcal{O}_X|_Z$-module sur $Z$.
Nous obtenons ainsi un foncteur
$$
\textit{Mod}(\mathcal{O}_X) \longrightarrow \textit{Mod}(\mathcal{O}_X|_Z),
\quad
\mathcal{F} \longmapsto \mathcal{H}_Z(\mathcal{F})
\text{ considéré comme }\mathcal{O}_X|_Z\text{-module sur }Z
$$
Ce foncteur est exact à gauche, mais n'est pas exact en général.
Tous les énoncés ci-dessus résultent directement du
lemme \ref{lemma-support-section-closed}.
Cette construction est manifestement compatible avec celle de la
remarque \ref{remark-sections-support-in-closed}.
\end{remark}

\begin{lemma}
\label{lemma-adjoint-section-with-support}
Soit $(X, \mathcal{O}_X)$ un espace annelé. Soit $i : Z \to X$
l'inclusion d'un fermé. Le foncteur
$\mathcal{H}_Z : \textit{Mod}(\mathcal{O}_X) \to
\textit{Mod}(\mathcal{O}_X|_Z)$ de la
remarque \ref{remark-sections-support-in-closed-modules}
est adjoint à droite de
$i_* : \textit{Mod}(\mathcal{O}_X|_Z) \to \textit{Mod}(\mathcal{O}_X)$.
\end{lemma}

\begin{proof}
Nous devons montrer que, pour tout $\mathcal{O}_X$-module $\mathcal{F}$
et tout $\mathcal{O}_X|_Z$-module $\mathcal{G}$, nous avons
$$
\Hom_{\mathcal{O}_X|_Z}(\mathcal{G}, \mathcal{H}_Z(\mathcal{F})) =
\Hom_{\mathcal{O}_X}(i_*\mathcal{G}, \mathcal{F})
$$
C'est clair,
car toute section de $i_*\mathcal{G}$ est après tout à support dans $Z$.
Les détails sont omis.
\end{proof}










\section{Faisceaux localement libres}
\label{section-locally-free}

\noindent
Soit $(X, \mathcal{O}_X)$ un espace annelé.
Nos conventions permettent que (certains) germes $\mathcal{O}_{X, x}$
soient l'anneau nul. Cela signifie qu'il faut prendre quelques précautions
lors de la définition du rang d'un faisceau localement libre.

\begin{definition}
\label{definition-locally-free}
Soit $(X, \mathcal{O}_X)$ un espace annelé.
Soit $\mathcal{F}$ un faisceau de $\mathcal{O}_X$-modules.
\begin{enumerate}
\item On dit que $\mathcal{F}$ est {\it localement libre} si, pour tout
point $x \in X$, il existe un ensemble $I$ et un
voisinage ouvert $x \in U \subset X$
tels que $\mathcal{F}|_U$ soit isomorphe à
$\bigoplus_{i \in I} \mathcal{O}_X|_U$ comme $\mathcal{O}_X|_U$-module.
\item On dit que $\mathcal{F}$ est {\it localement libre de type fini} si l'on peut
choisir les ensembles d'indices $I$ finis.
\item On dit que $\mathcal{F}$ est {\it localement libre de type fini et de rang $r$}
si l'on peut choisir les ensembles d'indices $I$ de cardinal $r$.
\end{enumerate}
\end{definition}

\noindent
Une somme directe finie de faisceaux localement libres (de type fini) est
localement libre (de type fini). Toutefois, une somme directe infinie
de faisceaux localement libres peut ne pas être localement libre.

\begin{lemma}
\label{lemma-locally-free-quasi-coherent}
Soit $(X, \mathcal{O}_X)$ un espace annelé.
Soit $\mathcal{F}$ un faisceau de $\mathcal{O}_X$-modules.
Si $\mathcal{F}$ est localement libre, alors il est quasi-cohérent.
\end{lemma}

\begin{proof}
Omis.
\end{proof}

\begin{lemma}
\label{lemma-pullback-locally-free}
Soit $f : (X, \mathcal{O}_X) \to (Y, \mathcal{O}_Y)$
un morphisme d'espaces annelés. Si $\mathcal{G}$ est
un $\mathcal{O}_Y$-module localement libre, alors
$f^*\mathcal{G}$ est un $\mathcal{O}_X$-module localement libre.
\end{lemma}

\begin{proof}
Omis.
\end{proof}

\begin{lemma}
\label{lemma-rank}
Soit $(X, \mathcal{O}_X)$ un espace annelé.
Supposons que le support de $\mathcal{O}_X$ soit $X$,
c'est-à-dire que tous les germes de $\mathcal{O}_X$ soient des anneaux non nuls.
Soit $\mathcal{F}$ un faisceau localement libre de $\mathcal{O}_X$-modules.
Il existe une fonction localement constante
$$
\text{rang}_\mathcal{F} :
X \longrightarrow \{0, 1, 2, \ldots\}\cup\{\infty\}
$$
telle que, pour tout point $x \in X$, le cardinal de tout
ensemble $I$ tel que $\mathcal{F}$ soit isomorphe à
$\bigoplus_{i\in I} \mathcal{O}_X$ dans un voisinage
de $x$ soit $\text{rang}_\mathcal{F}(x)$.
\end{lemma}

\begin{proof}
Sous l'hypothèse du lemme, le cardinal de $I$ se lit
sur le rang du module libre $\mathcal{F}_x$ sur l'anneau non nul
$\mathcal{O}_{X, x}$, et il est constant dans un voisinage de $x$.
\end{proof}

\begin{lemma}
\label{lemma-map-finite-locally-free}
Soit $(X, \mathcal{O}_X)$ un espace annelé. Soit $r \geq 0$.
Soit $\varphi : \mathcal{F} \to \mathcal{G}$ une application de
$\mathcal{O}_X$-modules localement libres de type fini et de rang $r$.
Alors $\varphi$ est un isomorphisme si et seulement si $\varphi$
est surjective.
\end{lemma}

\begin{proof}
Supposons $\varphi$ surjective. Choisissons $x \in X$.
Il existe un voisinage ouvert $U$ de $x$
tel que $\mathcal{F}|_U$ et $\mathcal{G}|_U$ soient tous deux
isomorphes à $\mathcal{O}_U^{\oplus r}$.
Choisissons des relèvements des générateurs libres de $\mathcal{G}|_U$
pour obtenir une application $\psi : \mathcal{G}|_U \to \mathcal{F}|_U$
telle que $\varphi|_U \circ \psi = \text{id}$.
Nous en déduisons que l'application
$\Gamma(U, \mathcal{F}) \to \Gamma(U, \mathcal{G})$
induite par $\varphi$ est surjective. Puisque
$\Gamma(U, \mathcal{F})$ et $\Gamma(U, \mathcal{G})$
sont isomorphes à $\Gamma(U, \mathcal{O}_U)^{\oplus r}$
comme $\Gamma(U, \mathcal{O}_U)$-modules, nous pouvons
appliquer le lemme \ref{algebra-lemma-fun} d'Algèbre pour voir que
$\Gamma(U, \mathcal{F}) \to \Gamma(U, \mathcal{G})$
est injective. Ceci achève la démonstration.
\end{proof}

\begin{lemma}
\label{lemma-direct-summand-of-locally-free-is-locally-free}
Soit $(X, \mathcal{O}_X)$ un espace annelé. Si tous les germes $\mathcal{O}_{X, x}$
sont des anneaux locaux, alors tout facteur direct d'un
$\mathcal{O}_X$-module localement libre de type fini est localement libre de type fini.
\end{lemma}

\begin{proof}
Supposons que $\mathcal{F}$ soit un facteur direct du
$\mathcal{O}_X$-module localement libre de type fini $\mathcal{H}$.
Soit $x \in X$ un point. Alors $\mathcal{H}_x$ est un
$\mathcal{O}_{X, x}$-module libre de type fini.
Comme $\mathcal{O}_{X, x}$ est local, on voit que
$\mathcal{F}_x \cong \mathcal{O}_{X, x}^{\oplus r}$ pour un certain $r$, voir le
lemme \ref{algebra-lemma-finite-projective} d'Algèbre.
D'après le lemme \ref{lemma-finite-presentation-stalk-free},
$\mathcal{F}$ est libre de rang $r$ dans un voisinage ouvert de $x$.
(Remarquons que $\mathcal{F}$ est de présentation finie comme facteur direct de
$\mathcal{H}$.)
\end{proof}





\section{Applications bilinéaires}
\label{section-bilinear}


\noindent
Soit $(X, \mathcal{O}_X)$ un espace annelé.
Soient $\mathcal{F}$, $\mathcal{G}$ et $\mathcal{H}$ des $\mathcal{O}_X$-modules.
Une {\it application bilinéaire} $f : \mathcal{F} \times \mathcal{G} \to \mathcal{H}$
de faisceaux de $\mathcal{O}_X$-modules est une application de faisceaux d'ensembles comme indiqué
telle que, pour tout ouvert $U \subset X$, l'application induite
$$
\mathcal{F}(U) \times \mathcal{G}(U) \to \mathcal{H}(U)
$$
soit une application bilinéaire de $\mathcal{O}_X(U)$-modules. De manière équivalente, on peut demander
que certains diagrammes d'applications de faisceaux d'ensembles commutent, en imitant les
axiomes usuels des applications bilinéaires de modules. Par exemple, l'axiome
$f(x + y, z) = f(x, z) + f(y, z)$ est représenté par la commutativité
du diagramme
$$
\xymatrix{
\mathcal{F} \times \mathcal{F} \times \mathcal{G}
\ar[rrr]_{(f \circ \text{pr}_{13}, f \circ \text{pr}_{23})}
\ar[d]_{(+ \circ \text{pr}_{12}, \text{pr}_3)} & & &
\mathcal{H} \times \mathcal{H}
\ar[d]^{+} \\
\mathcal{F} \times \mathcal{G} \ar[rrr]^f & & &
\mathcal{H}
}
$$
Une autre caractérisation est la suivante : si
$f : \mathcal{F} \times \mathcal{G} \to \mathcal{H}$
est une application de faisceaux d'ensembles et
induit une application bilinéaire de modules sur les germes en tout point de $X$, alors
$f$ est une application bilinéaire de faisceaux de modules. En effet, on peut vérifier
l'égalité de sections locales sur les germes.

\medskip\noindent
Notons $\text{Mor}( - , - )$ les morphismes dans la catégorie des faisceaux
d'ensembles sur $X$. Voici une autre caractérisation d'une application bilinéaire :
une application de faisceaux d'ensembles
$f : \mathcal{F} \times \mathcal{G} \to \mathcal{H}$
est bilinéaire si, pour tout faisceau d'ensembles $\mathcal{S}$, la règle
$$
\text{Mor}(\mathcal{S}, \mathcal{F}) \times
\text{Mor}(\mathcal{S}, \mathcal{G}) \to
\text{Mor}(\mathcal{S}, \mathcal{H}),\quad
(a, b) \mapsto f \circ (a \times b)
$$
est une application bilinéaire de modules sur l'anneau
$\text{Mor}(\mathcal{S}, \mathcal{O}_X)$.
Nous n'adoptons généralement pas ce point de vue, car il est plus facile de raisonner sur les
ensembles de sections locales, et les deux points de vue sont manifestement équivalents.

\medskip\noindent
Enfin, voici encore une autre formulation de la définition :
$\mathcal{O}_X$ est un objet en anneaux dans la catégorie des faisceaux d'ensembles,
et $\mathcal{F}$, $\mathcal{G}$, $\mathcal{H}$ sont des objets en modules
sur cet anneau. On peut alors définir une application bilinéaire entre objets en modules
sur un objet en anneaux dans toute catégorie.
Pour formuler ce qu'est un objet en anneaux, ce qu'est un objet en modules
sur un objet en anneaux, et ce qu'est une application bilinéaire entre de tels objets dans une
catégorie, il est commode (mais pas strictement nécessaire)
de supposer que la catégorie possède des produits finis ;
c'est le cas de la catégorie des faisceaux d'ensembles.







\section{Produit tensoriel}
\label{section-tensor-product}

\noindent
Nous avons déjà brièvement abordé le produit tensoriel dans le cadre
du changement d'anneaux dans Faisceaux, sections
\ref{sheaves-section-presheaves-modules} et
\ref{sheaves-section-sheafification-presheaves-modules}.
Généralisons cette discussion aux produits tensoriels de modules.

\medskip\noindent
Soit $(X, \mathcal{O}_X)$ un espace annelé et soient
$\mathcal{F}$ et $\mathcal{G}$ des $\mathcal{O}_X$-modules.
Nous définissons d'abord le {\it préfaisceau produit tensoriel}
$$
\mathcal{F} \otimes_{p, \mathcal{O}_X} \mathcal{G}
$$
comme la règle qui associe à tout ouvert $U \subset X$ le
$\mathcal{O}_X(U)$-module
$\mathcal{F}(U) \otimes_{\mathcal{O}_X(U)} \mathcal{G}(U)$.
Cela étant défini, nous appelons {\it faisceau produit tensoriel}
le faisceau associé au préfaisceau précédent :
$$
\mathcal{F} \otimes_{\mathcal{O}_X} \mathcal{G}
=
(\mathcal{F} \otimes_{p, \mathcal{O}_X} \mathcal{G})^\#
$$
On peut le caractériser comme le faisceau de
$\mathcal{O}_X$-modules tel que, pour tout troisième
faisceau de $\mathcal{O}_X$-modules $\mathcal{H}$,
on ait
$$
\Hom_{\mathcal{O}_X}
(\mathcal{F} \otimes_{\mathcal{O}_X} \mathcal{G}, \mathcal{H})
=
\text{Bilin}_{\mathcal{O}_X}(\mathcal{F} \times \mathcal{G}, \mathcal{H}).
$$
Ici, le membre de droite désigne l'ensemble des applications bilinéaires de faisceaux
de $\mathcal{O}_X$-modules définies dans la section \ref{section-bilinear}.

\medskip\noindent
Le produit tensoriel de modules $M, N$ sur un anneau $R$ est symétrique,
à savoir $M \otimes_R N = N \otimes_R M$ ; il en va donc de même pour
les produits tensoriels de faisceaux de modules, c'est-à-dire que l'on a
$$
\mathcal{F} \otimes_{\mathcal{O}_X} \mathcal{G}
=
\mathcal{G} \otimes_{\mathcal{O}_X} \mathcal{F}
$$
fonctoriellement en $\mathcal{F}$, $\mathcal{G}$.
Puisque le produit tensoriel de modules est associatif, nous
obtenons aussi des isomorphismes canoniques fonctoriels
$$
(\mathcal{F} \otimes_{\mathcal{O}_X} \mathcal{G})
\otimes_{\mathcal{O}_X} \mathcal{H}
=
\mathcal{F} \otimes_{\mathcal{O}_X}
(\mathcal{G} \otimes_{\mathcal{O}_X} \mathcal{H})
$$
en $\mathcal{F}$, $\mathcal{G}$ et $\mathcal{H}$.

\begin{lemma}
\label{lemma-stalk-tensor-product}
Soit $(X, \mathcal{O}_X)$ un espace annelé.
Soient $\mathcal{F}$, $\mathcal{G}$ des $\mathcal{O}_X$-modules.
Soit $x \in X$. Il existe un isomorphisme canonique
de $\mathcal{O}_{X, x}$-modules
$$
(\mathcal{F} \otimes_{\mathcal{O}_X} \mathcal{G})_x
=
\mathcal{F}_x \otimes_{\mathcal{O}_{X, x}} \mathcal{G}_x
$$
fonctoriel en $\mathcal{F}$ et $\mathcal{G}$.
\end{lemma}

\begin{proof}
La démonstration est omise.
\end{proof}

\begin{lemma}
\label{lemma-tensor-product-sheafification}
Soit $(X, \mathcal{O}_X)$ un espace annelé.
Soient $\mathcal{F}'$, $\mathcal{G}'$ des préfaisceaux de $\mathcal{O}_X$-modules
dont les faisceaux associés sont $\mathcal{F}$, $\mathcal{G}$. Alors
$\mathcal{F} \otimes_{\mathcal{O}_X} \mathcal{G} =
(\mathcal{F}' \otimes_{p, \mathcal{O}_X} \mathcal{G}')^\#$.
\end{lemma}

\begin{proof}
La démonstration est omise.
\end{proof}


\begin{lemma}
\label{lemma-tensor-product-exact}
Soit $(X, \mathcal{O}_X)$ un espace annelé.
Soit $\mathcal{G}$ un $\mathcal{O}_X$-module.
Si
$\mathcal{F}_1
\to \mathcal{F}_2
\to \mathcal{F}_3
\to 0$
est une suite exacte de $\mathcal{O}_X$-modules, alors
la suite induite
$$
\mathcal{F}_1 \otimes_{\mathcal{O}_X} \mathcal{G} \to
\mathcal{F}_2 \otimes_{\mathcal{O}_X} \mathcal{G} \to
\mathcal{F}_3 \otimes_{\mathcal{O}_X} \mathcal{G} \to
0
$$
est exacte.
\end{lemma}

\begin{proof}
Cela résulte du fait que l'exactitude peut se vérifier sur les germes
(lemme \ref{lemma-abelian}), de la description des germes
(lemme \ref{lemma-stalk-tensor-product}) et du résultat correspondant
pour les produits tensoriels de modules
(Algèbre, lemme \ref{algebra-lemma-tensor-product-exact}).
\end{proof}

\begin{lemma}
\label{lemma-tensor-product-pullback}
Soit $f : (X, \mathcal{O}_X) \to (Y, \mathcal{O}_Y)$
un morphisme d'espaces annelés. Soient $\mathcal{F}$, $\mathcal{G}$
des $\mathcal{O}_Y$-modules. Alors
$f^*(\mathcal{F} \otimes_{\mathcal{O}_Y} \mathcal{G})
= f^*\mathcal{F} \otimes_{\mathcal{O}_X} f^*\mathcal{G}$
fonctoriellement en $\mathcal{F}$, $\mathcal{G}$.
\end{lemma}

\begin{proof}
La démonstration est omise.
\end{proof}

\begin{lemma}
\label{lemma-tensor-commute-colimits}
Soit $(X, \mathcal{O}_X)$ un espace annelé.
Pour tout $\mathcal{O}_X$-module $\mathcal{F}$, le foncteur
$$
\textit{Mod}(\mathcal{O}_X) \longrightarrow \textit{Mod}(\mathcal{O}_X)
, \quad
\mathcal{G} \longmapsto \mathcal{F} \otimes_{\mathcal{O}_X} \mathcal{G}
$$
commute aux colimites quelconques.
\end{lemma}

\begin{proof}
Soit $I$ un ensemble préordonné et soit $\{\mathcal{G}_i\}$
un système indexé par $I$. Posons $\mathcal{G} = \colim_i \mathcal{G}_i$.
Rappelons que $\mathcal{G}$ est le faisceau associé au préfaisceau
$\mathcal{G}' : U \mapsto \colim_i \mathcal{G}_i(U)$ ; voir
Faisceaux, section \ref{sheaves-section-limits-sheaves}.
D'après le
lemme \ref{lemma-tensor-product-sheafification},
le produit tensoriel $\mathcal{F} \otimes_{\mathcal{O}_X} \mathcal{G}$
est le faisceau associé au préfaisceau
$$
U \longmapsto
\mathcal{F}(U) \otimes_{\mathcal{O}_X(U)} \colim_i \mathcal{G}_i(U) =
\colim_i \mathcal{F}(U) \otimes_{\mathcal{O}_X(U)} \mathcal{G}_i(U)
$$
où l'égalité est donnée par
Algèbre, lemme \ref{algebra-lemma-tensor-products-commute-with-limits}.
Le lemme résulte donc de la description des colimites dans
$\textit{Mod}(\mathcal{O}_X)$ ; voir le lemme \ref{lemma-limits-colimits}.
\end{proof}

\begin{lemma}
\label{lemma-tensor-product-permanence}
Soit $(X, \mathcal{O}_X)$ un espace annelé.
Soient $\mathcal{F}$, $\mathcal{G}$ des $\mathcal{O}_X$-modules.
\begin{enumerate}
\item Si $\mathcal{F}$, $\mathcal{G}$ sont localement engendrés
par leurs sections, il en va de même pour $\mathcal{F} \otimes_{\mathcal{O}_X} \mathcal{G}$.
\item Si $\mathcal{F}$, $\mathcal{G}$ sont de type fini,
il en va de même pour $\mathcal{F} \otimes_{\mathcal{O}_X} \mathcal{G}$.
\item Si $\mathcal{F}$, $\mathcal{G}$ sont quasi-cohérents,
il en va de même pour $\mathcal{F} \otimes_{\mathcal{O}_X} \mathcal{G}$.
\item Si $\mathcal{F}$, $\mathcal{G}$ sont de présentation finie,
il en va de même pour $\mathcal{F} \otimes_{\mathcal{O}_X} \mathcal{G}$.
\item Si $\mathcal{F}$ est de présentation finie et $\mathcal{G}$ est cohérent,
alors $\mathcal{F} \otimes_{\mathcal{O}_X} \mathcal{G}$ est cohérent.
\item Si $\mathcal{F}$, $\mathcal{G}$ sont cohérents,
il en va de même pour $\mathcal{F} \otimes_{\mathcal{O}_X} \mathcal{G}$.
\item Si $\mathcal{F}$, $\mathcal{G}$ sont localement libres,
il en va de même pour $\mathcal{F} \otimes_{\mathcal{O}_X} \mathcal{G}$.
\end{enumerate}
\end{lemma}

\begin{proof}
Montrons d'abord que le produit tensoriel de
$\mathcal{O}_X$-modules localement libres est localement libre. Il suffit de montrer
que
$(\bigoplus_{i \in I} \mathcal{O}_X) \otimes_{\mathcal{O}_X}
(\bigoplus_{j \in J} \mathcal{O}_X) \cong
\bigoplus_{(i, j) \in I \times J} \mathcal{O}_X$.
Le faisceau $\bigoplus_{i \in I} \mathcal{O}_X$ est le faisceau associé
au préfaisceau $U \mapsto \bigoplus_{i \in I} \mathcal{O}_X(U)$.
Par conséquent, le produit tensoriel est le faisceau associé
au préfaisceau
$$
U \longmapsto
(\bigoplus\nolimits_{i \in I} \mathcal{O}_X(U))
\otimes_{\mathcal{O}_X(U)}
(\bigoplus\nolimits_{j \in J} \mathcal{O}_X(U)).
$$
On en déduit le résultat voulu, puisque pour tout anneau $R$ on a
$(\bigoplus_{i \in I} R) \otimes_R (\bigoplus_{j \in J} R) =
\bigoplus_{(i, j) \in I \times J} R$.

\medskip\noindent
Si $\mathcal{F}_2 \to \mathcal{F}_1 \to \mathcal{F} \to 0$
est exacte, alors, d'après le lemme \ref{lemma-tensor-product-exact},
le complexe
$\mathcal{F}_2 \otimes_{\mathcal{O}_X} \mathcal{G} \to
\mathcal{F}_1 \otimes_{\mathcal{O}_X} \mathcal{G} \to
\mathcal{F} \otimes_{\mathcal{O}_X} \mathcal{G} \to 0$
est exact. Cela permet de démontrer (5). En effet, dans ce cas, il existe
localement une telle suite exacte où les $\mathcal{F}_i$, $i = 1, 2$, sont
libres de type fini. Les deux termes
$\mathcal{F}_2 \otimes_{\mathcal{O}_X} \mathcal{G}$ et
$\mathcal{F}_1 \otimes_{\mathcal{O}_X} \mathcal{G}$
sont donc isomorphes à des sommes directes finies de $\mathcal{G}$ (par exemple
d'après le lemme \ref{lemma-tensor-commute-colimits}).
Comme les sommes directes finies sont des faisceaux cohérents, ces termes sont cohérents,
ainsi que le conoyau de l'application ; voir le lemme \ref{lemma-coherent-abelian}.

\medskip\noindent
Et si, de plus,
$\mathcal{G}_2 \to \mathcal{G}_1 \to \mathcal{G} \to 0$
est exacte, alors on voit que
$$
\mathcal{F}_2 \otimes_{\mathcal{O}_X} \mathcal{G}_1
\oplus
\mathcal{F}_1 \otimes_{\mathcal{O}_X} \mathcal{G}_2
\to
\mathcal{F}_1 \otimes_{\mathcal{O}_X} \mathcal{G}_1
\to
\mathcal{F} \otimes_{\mathcal{O}_X} \mathcal{G}
\to 0
$$
est exacte. Cela permet, par exemple, de démontrer (3).
En effet, l'hypothèse signifie que l'on peut trouver localement des présentations
comme ci-dessus où les $\mathcal{F}_i$ et les $\mathcal{G}_i$ sont des
$\mathcal{O}_X$-modules libres. La présentation affichée est donc aussi
une présentation du produit tensoriel par des faisceaux libres.

\medskip\noindent
La démonstration des autres assertions est omise.
\end{proof}





\section{Modules plats}
\label{section-flat}

\noindent
On peut définir les modules plats exactement comme dans le cas des modules sur un anneau.

\begin{definition}
\label{definition-flat}
Soit $(X, \mathcal{O}_X)$ un espace annelé.
Un $\mathcal{O}_X$-module $\mathcal{F}$ est {\it plat} si le foncteur
$$
\textit{Mod}(\mathcal{O}_X)
\longrightarrow
\textit{Mod}(\mathcal{O}_X), \quad
\mathcal{G} \mapsto \mathcal{G} \otimes_{\mathcal{O}_X} \mathcal{F}
$$
est exact.
\end{definition}

\noindent
On peut caractériser la platitude au moyen des germes.

\begin{lemma}
\label{lemma-flat-stalks-flat}
Soit $(X, \mathcal{O}_X)$ un espace annelé.
Un $\mathcal{O}_X$-module $\mathcal{F}$ est plat
si et seulement si le germe $\mathcal{F}_x$ est un
$\mathcal{O}_{X, x}$-module plat pour tout $x \in X$.
\end{lemma}

\begin{proof}
Supposons que $\mathcal{F}_x$ soit un $\mathcal{O}_{X, x}$-module plat pour tout
$x \in X$. Dans ce cas, si $\mathcal{G} \to \mathcal{H} \to \mathcal{K}$
est exacte, alors
$\mathcal{G} \otimes_{\mathcal{O}_X} \mathcal{F} \to
\mathcal{H} \otimes_{\mathcal{O}_X} \mathcal{F} \to
\mathcal{K} \otimes_{\mathcal{O}_X} \mathcal{F}$
est également exacte, car on peut vérifier l'exactitude sur les germes et
le produit tensoriel commute à la prise des germes ; voir le
lemme \ref{lemma-stalk-tensor-product}.
Réciproquement, supposons que $\mathcal{F}$ soit plat, et soit $x \in X$.
Considérons les faisceaux gratte-ciel $i_{x, *} M$, où $M$ est un
$\mathcal{O}_{X, x}$-module. Remarquons que
$$
M \otimes_{\mathcal{O}_{X, x}} \mathcal{F}_x =
\left(i_{x, *} M \otimes_{\mathcal{O}_X} \mathcal{F}\right)_x
$$
de nouveau d'après le
lemme \ref{lemma-stalk-tensor-product}.
Puisque $i_{x, *}$ est exact, la platitude de $\mathcal{F}$ implique que
$M \mapsto M \otimes_{\mathcal{O}_{X, x}} \mathcal{F}_x$
est exact. Ainsi $\mathcal{F}_x$ est un $\mathcal{O}_{X, x}$-module plat.
\end{proof}

\noindent
La définition suivante a donc un sens.

\begin{definition}
\label{definition-flat-at-point}
Soit $(X, \mathcal{O}_X)$ un espace annelé. Soit $x \in X$.
Un $\mathcal{O}_X$-module $\mathcal{F}$ est
{\it plat en $x$} si $\mathcal{F}_x$ est un
$\mathcal{O}_{X, x}$-module plat.
\end{definition}

\noindent
Ainsi, $\mathcal{F}$ est un $\mathcal{O}_X$-module plat
si et seulement s'il est plat en tout point.

\begin{lemma}
\label{lemma-base-change-flat}
Soit $f : (X, \mathcal{O}_X) \to (Y, \mathcal{O}_Y)$
un morphisme d'espaces annelés. Si $\mathcal{G}$ est un
$\mathcal{O}_Y$-module plat, alors $f^*\mathcal{G}$ est un
$\mathcal{O}_X$-module plat.
\end{lemma}

\begin{proof}
Combiner le lemme \ref{lemma-flat-stalks-flat} avec
Faisceaux, lemme \ref{sheaves-lemma-stalk-pullback-modules}, et
Algèbre, lemme \ref{algebra-lemma-flat-base-change}.
\end{proof}

\begin{lemma}
\label{lemma-colimits-flat}
Soit $(X, \mathcal{O}_X)$ un espace annelé.
Une colimite filtrante de $\mathcal{O}_X$-modules plats est plate.
Une somme directe de $\mathcal{O}_X$-modules plats est plate.
\end{lemma}

\begin{proof}
Cela résulte du
lemme \ref{lemma-tensor-commute-colimits}, du
lemme \ref{lemma-stalk-tensor-product}, d'
Algèbre, lemme \ref{algebra-lemma-directed-colimit-exact},
et du fait que l'on peut vérifier l'exactitude sur les germes.
\end{proof}

\begin{lemma}
\label{lemma-j-shriek-flat}
Soit $(X, \mathcal{O}_X)$ un espace annelé.
Soit $U \subset X$ un ouvert. Le faisceau $j_{U!}\mathcal{O}_U$
est un faisceau plat de $\mathcal{O}_X$-modules.
\end{lemma}

\begin{proof}
Les germes de $j_{U!}\mathcal{O}_U$ sont soit nuls, soit égaux
à $\mathcal{O}_{X, x}$. Appliquer le
lemme \ref{lemma-flat-stalks-flat}.
\end{proof}

\begin{lemma}
\label{lemma-module-quotient-flat}
Soit $(X, \mathcal{O}_X)$ un espace annelé.
\begin{enumerate}
\item Tout faisceau de $\mathcal{O}_X$-modules est quotient d'une
somme directe $\bigoplus j_{U_i!}\mathcal{O}_{U_i}$.
\item Tout $\mathcal{O}_X$-module est quotient d'un
$\mathcal{O}_X$-module plat.
\end{enumerate}
\end{lemma}

\begin{proof}
Soit $\mathcal{F}$ un $\mathcal{O}_X$-module.
Pour tout ouvert $U \subset X$ et tout
$s \in \mathcal{F}(U)$, on obtient un morphisme
$j_{U!}\mathcal{O}_U \to \mathcal{F}$, à savoir l'adjoint du
morphisme $\mathcal{O}_U \to \mathcal{F}|_U$, $1 \mapsto s$.
L'application
$$
\bigoplus\nolimits_{(U, s)} j_{U!}\mathcal{O}_U
\longrightarrow
\mathcal{F}
$$
est manifestement surjective, et sa source est plate d'après les lemmes
\ref{lemma-colimits-flat} et \ref{lemma-j-shriek-flat}.
\end{proof}

\begin{lemma}
\label{lemma-flat-tor-zero}
Soit $(X, \mathcal{O}_X)$ un espace annelé.
Soit
$$
0 \to \mathcal{F}'' \to \mathcal{F}' \to \mathcal{F} \to 0
$$
une suite exacte courte de $\mathcal{O}_X$-modules.
Supposons $\mathcal{F}$ plat. Alors, pour tout $\mathcal{O}_X$-module
$\mathcal{G}$, la suite
$$
0 \to
\mathcal{F}'' \otimes_\mathcal{O} \mathcal{G} \to
\mathcal{F}' \otimes_\mathcal{O} \mathcal{G} \to
\mathcal{F} \otimes_\mathcal{O} \mathcal{G} \to 0
$$
est exacte.
\end{lemma}

\begin{proof}
Comme $\mathcal{F}_x$ est un $\mathcal{O}_{X, x}$-module plat
pour tout $x \in X$ et que l'exactitude se vérifie sur les germes, le résultat
découle d'
Algèbre, lemme \ref{algebra-lemma-flat-tor-zero}.
\end{proof}

\begin{lemma}
\label{lemma-flat-ses}
\begin{slogan}
Les noyaux d'épimorphismes et les extensions de faisceaux plats de modules sur
un espace annelé sont encore plats.
\end{slogan}
Soit $(X, \mathcal{O}_X)$ un espace annelé.
Soit
$$
0 \to
\mathcal{F}_2 \to
\mathcal{F}_1 \to
\mathcal{F}_0 \to 0
$$
une suite exacte courte de $\mathcal{O}_X$-modules.
\begin{enumerate}
\item Si $\mathcal{F}_2$ et $\mathcal{F}_0$ sont plats, il en va de même pour
$\mathcal{F}_1$.
\item Si $\mathcal{F}_1$ et $\mathcal{F}_0$ sont plats, il en va de même pour
$\mathcal{F}_2$.
\end{enumerate}
\end{lemma}

\begin{proof}
Comme l'exactitude et la platitude se vérifient sur les germes,
cela résulte d'
Algèbre, lemme \ref{algebra-lemma-flat-ses}.
\end{proof}

\begin{lemma}
\label{lemma-flat-resolution-of-flat}
Soit $(X, \mathcal{O}_X)$ un espace annelé.
Soit
$$
\ldots \to
\mathcal{F}_2 \to
\mathcal{F}_1 \to
\mathcal{F}_0 \to
\mathcal{Q} \to 0
$$
un complexe exact de $\mathcal{O}_X$-modules.
Si $\mathcal{Q}$ et tous les $\mathcal{F}_i$ sont des $\mathcal{O}_X$-modules plats,
alors, pour tout $\mathcal{O}_X$-module $\mathcal{G}$, le complexe
$$
\ldots \to
\mathcal{F}_2 \otimes_{\mathcal{O}_X} \mathcal{G} \to
\mathcal{F}_1 \otimes_{\mathcal{O}_X} \mathcal{G} \to
\mathcal{F}_0 \otimes_{\mathcal{O}_X} \mathcal{G} \to
\mathcal{Q} \otimes_{\mathcal{O}_X} \mathcal{G} \to 0
$$
est également exact.
\end{lemma}

\begin{proof}
Cela résulte du lemme \ref{lemma-flat-tor-zero}, en décomposant le complexe
en suites exactes courtes et en utilisant le lemme \ref{lemma-flat-ses} pour
montrer par récurrence que $\Im(\mathcal{F}_{i + 1} \to \mathcal{F}_i)$
est plat.
\end{proof}

\noindent
Le lemme suivant donne une implication du critère équationnel de
platitude (Algèbre, lemme \ref{algebra-lemma-flat-eq}).

\begin{lemma}
\label{lemma-flat-eq}
Soit $(X, \mathcal{O}_X)$ un espace annelé. Soit $\mathcal{F}$ un
$\mathcal{O}_X$-module plat. Soit $U \subset X$ un ouvert et soit
$$
\mathcal{O}_U \xrightarrow{(f_1, \ldots, f_n)}
\mathcal{O}_U^{\oplus n} \xrightarrow{(s_1, \ldots, s_n)}
\mathcal{F}|_U
$$
un complexe de $\mathcal{O}_U$-modules. Pour tout $x \in U$, il
existe un voisinage ouvert $V \subset U$ de $x$ et une factorisation
$$
\mathcal{O}_V^{\oplus n}
\xrightarrow{A}
\mathcal{O}_V^{\oplus m} \xrightarrow{(t_1, \ldots, t_m)}
\mathcal{F}|_V
$$
de $(s_1, \ldots, s_n)|_V$ telle que $A \circ (f_1, \ldots, f_n)|_V = 0$.
\end{lemma}

\begin{proof}
Soit $\mathcal{I} \subset \mathcal{O}_U$ le faisceau d'idéaux
engendré par $f_1, \ldots, f_n$. Alors $\sum f_i \otimes s_i$ est
une section de $\mathcal{I} \otimes_{\mathcal{O}_U} \mathcal{F}|_U$
dont l'image dans $\mathcal{F}|_U$ est nulle. Comme $\mathcal{F}|_U$ est plat,
l'application
$\mathcal{I} \otimes_{\mathcal{O}_U} \mathcal{F}|_U \to \mathcal{F}|_U$
est injective. Puisque $\mathcal{I} \otimes_{\mathcal{O}_U} \mathcal{F}|_U$
est le faisceau associé au préfaisceau produit tensoriel, il existe
un voisinage ouvert $V \subset U$ de $x$ tel
que $\sum f_i|_V \otimes s_i|_V$ soit nul dans
$\mathcal{I}(V) \otimes_{\mathcal{O}(V)} \mathcal{F}(V)$.
En explicitant les définitions à l'aide d'Algèbre, lemme \ref{algebra-lemma-relations},
on trouve $t_1, \ldots, t_m \in \mathcal{F}(V)$ et $a_{ij} \in \mathcal{O}(V)$
tels que $\sum a_{ij}f_i|_V = 0$ et $s_i|_V = \sum a_{ij}t_j$.
\end{proof}














\section{Duaux}
\label{section-duals}

\noindent
Soit $(X, \mathcal{O}_X)$ un espace annelé. La catégorie des
$\mathcal{O}_X$-modules, munie du produit tensoriel
construit dans la section \ref{section-tensor-product},
est une catégorie monoïdale symétrique. Pour un $\mathcal{O}_X$-module
$\mathcal{F}$, les conditions suivantes sont équivalentes :
\begin{enumerate}
\item $\mathcal{F}$ admet un dual à gauche dans la catégorie monoïdale
des $\mathcal{O}_X$-modules ;
\item $\mathcal{F}$ est localement facteur direct d'un
$\mathcal{O}_X$-module libre de type fini ;
\item $\mathcal{F}$ est de présentation finie et plat comme
$\mathcal{O}_X$-module.
\end{enumerate}
Cela est démontré dans l'exemple \ref{example-dual} et les
lemmes \ref{lemma-left-dual-module} et
\ref{lemma-flat-locally-finite-presentation} de cette section.

\begin{example}
\label{example-dual}
Soit $(X, \mathcal{O}_X)$ un espace annelé. Soit $\mathcal{F}$
un $\mathcal{O}_X$-module qui est localement facteur direct d'un
$\mathcal{O}_X$-module libre de type fini. Alors l'application
$$
\mathcal{F} \otimes_{\mathcal{O}_X}
\SheafHom_{\mathcal{O}_X}(\mathcal{F}, \mathcal{O}_X)
\longrightarrow
\SheafHom_{\mathcal{O}_X}(\mathcal{F}, \mathcal{F})
$$
est un isomorphisme. En effet, la question est locale, l'assertion est vraie
si $\mathcal{F}$ est libre de type fini, et elle vaut pour tout facteur direct
d'un module pour lequel elle est vraie. Notons
$$
\eta :
\mathcal{O}_X
\longrightarrow
\mathcal{F} \otimes_{\mathcal{O}_X}
\SheafHom_{\mathcal{O}_X}(\mathcal{F}, \mathcal{O}_X)
$$
l'application qui envoie $1$ sur la section correspondant à
$\text{id}_\mathcal{F}$ par l'isomorphisme ci-dessus. Notons
$$
\epsilon : 
\SheafHom_{\mathcal{O}_X}(\mathcal{F}, \mathcal{O}_X)
\otimes_{\mathcal{O}_X} \mathcal{F}
\longrightarrow
\mathcal{O}_X
$$
l'application d'évaluation. Alors
$\SheafHom_{\mathcal{O}_X}(\mathcal{F}, \mathcal{O}_X), \eta, \epsilon$
est un dual à gauche de $\mathcal{F}$ au sens de
Catégories, définition \ref{categories-definition-dual}.
Nous omettons la vérification des égalités
$(1 \otimes \epsilon) \circ (\eta \otimes 1) = \text{id}_\mathcal{F}$
et
$(\epsilon \otimes 1) \circ (1 \otimes \eta) =
\text{id}_{\SheafHom_{\mathcal{O}_X}(\mathcal{F}, \mathcal{O}_X)}$.
\end{example}

\begin{lemma}
\label{lemma-left-dual-module}
Soit $(X, \mathcal{O}_X)$ un espace annelé. Soit $\mathcal{F}$ un
$\mathcal{O}_X$-module. Soit $\mathcal{G}, \eta, \epsilon$
un dual à gauche de $\mathcal{F}$ dans la catégorie monoïdale des
$\mathcal{O}_X$-modules ; voir
Catégories, définition \ref{categories-definition-dual}. Alors :
\begin{enumerate}
\item $\mathcal{F}$ est localement facteur direct d'un
$\mathcal{O}_X$-module libre de type fini ;
\item l'application
$e : \SheafHom_{\mathcal{O}_X}(\mathcal{F}, \mathcal{O}_X) \to \mathcal{G}$
qui envoie une section locale $\lambda$ sur $(\lambda \otimes 1)(\eta)$
est un isomorphisme ;
\item on a $\epsilon(f, g) = e^{-1}(g)(f)$ pour des sections locales
$f$ et $g$ de $\mathcal{F}$ et $\mathcal{G}$.
\end{enumerate}
\end{lemma}

\begin{proof}
Les hypothèses signifient que
$$
\mathcal{F} \xrightarrow{\eta \otimes 1}
\mathcal{F} \otimes_{\mathcal{O}_X} \mathcal{G}
\otimes_{\mathcal{O}_X} \mathcal{F}
\xrightarrow{1 \otimes \epsilon} \mathcal{F}
\quad\text{et}\quad
\mathcal{G} \xrightarrow{1 \otimes \eta}
\mathcal{G} \otimes_{\mathcal{O}_X} \mathcal{F}
\otimes_{\mathcal{O}_X} \mathcal{G}
\xrightarrow{\epsilon \otimes 1} \mathcal{G}
$$
sont les applications identiques. Soit $x \in X$. On peut trouver un voisinage ouvert
$U$ de $x$, un nombre fini de sections $f_1, \ldots, f_n$
et $g_1, \ldots, g_n$ de
$\mathcal{F}$ et $\mathcal{G}$ sur $U$ telles que
$\eta(1) = \sum f_i g_i$. Notons
$$
\mathcal{O}_U^{\oplus n} \to \mathcal{F}|_U
$$
l'application qui envoie le $i$-ième vecteur de base sur $f_i$. On peut alors
factoriser l'application $\eta|_U$ par une application
$\tilde \eta : \mathcal{O}_U \to
\mathcal{O}_U^{\oplus n} \otimes_{\mathcal{O}_U} \mathcal{G}|_U$.
On obtient un diagramme commutatif
$$
\xymatrix{
\mathcal{F}|_U
\ar[rr]_-{\eta \otimes 1} \ar[rrd]_-{\tilde \eta \otimes 1} & &
\mathcal{F}|_U \otimes \mathcal{G}|_U \otimes \mathcal{F}|_U
\ar[r]_-{1 \otimes \epsilon} &
\mathcal{F}|_U \\
& &
\mathcal{O}_U^{\oplus n} \otimes \mathcal{G}|_U \otimes \mathcal{F}|_U
\ar[u] \ar[r]^-{1 \otimes \epsilon} &
\mathcal{O}_U^{\oplus n} \ar[u]
}
$$
Cela montre que l'identité de $\mathcal{F}$, localement sur $X$,
se factorise par un module libre de type fini. Ceci démontre (1). L'assertion (2) résulte de
Catégories, lemme \ref{categories-lemma-left-dual}, et de sa démonstration.
L'assertion (3) résulte de la première égalité de la démonstration.
On peut aussi déduire (2) et (3) de l'unicité des duaux à gauche
(Catégories, remarque \ref{categories-remark-left-dual-adjoint})
et de la construction du dual à gauche dans
l'exemple \ref{example-dual}.
\end{proof}

\begin{lemma}
\label{lemma-flat-locally-finite-presentation}
Soit $(X, \mathcal{O}_X)$ un espace annelé. Soit $\mathcal{F}$ un
$\mathcal{O}_X$-module plat de présentation finie. Alors $\mathcal{F}$ est
localement facteur direct d'un $\mathcal{O}_X$-module libre de type fini.
\end{lemma}

\begin{proof}
Après avoir remplacé $X$ par les membres d'un recouvrement ouvert, on peut
supposer qu'il existe une présentation
$$
\mathcal{O}_X^{\oplus r} \to
\mathcal{O}_X^{\oplus n} \to \mathcal{F} \to 0
$$
Soit $x \in X$. D'après le lemme \ref{lemma-flat-eq},
on peut, après avoir restreint $X$ à un voisinage
ouvert de $x$, supposer qu'il existe une factorisation
$$
\mathcal{O}_X^{\oplus n} \to
\mathcal{O}_X^{\oplus n_1} \to \mathcal{F}
$$
telle que la composée
$\mathcal{O}_X^{\oplus r} \to \mathcal{O}_X^{\oplus n} \to
\mathcal{O}_X^{\oplus n_1}$
annule le premier facteur de $\mathcal{O}_X^{\oplus r}$.
En répétant cet argument encore $r - 1$ fois, on obtient une factorisation
$$
\mathcal{O}_X^{\oplus n} \to
\mathcal{O}_X^{\oplus n_r} \to \mathcal{F}
$$
telle que la composée
$\mathcal{O}_X^{\oplus r} \to \mathcal{O}_X^{\oplus n}
\to \mathcal{O}_X^{\oplus n_r}$ est nulle.
Cela signifie que la surjection $\mathcal{O}_X^{\oplus n_r} \to \mathcal{F}$
admet une section, ce qui conclut.
\end{proof}







\section{Faisceaux constructibles d'ensembles}
\label{section-constructible}

\noindent
Soit $X$ un espace topologique. Étant donné un ensemble $S$, rappelons que $\underline{S}$
ou $\underline{S}_X$ désigne le faisceau constant de valeur $S$ ; voir
Faisceaux, définition \ref{sheaves-definition-constant-sheaf}.
Soit $U \subset X$ un ouvert d'un espace topologique $X$.
Nous noterons $j_U$ le morphisme d'inclusion et
$j_{U!} : \Sh(U) \to \Sh(X)$ le prolongement par l'ensemble vide
décrit dans Faisceaux, section \ref{sheaves-section-open-immersions}.

\begin{lemma}
\label{lemma-surjection}
Soit $X$ un espace topologique. Soit $\mathcal{B}$ une base de la
topologie de $X$. Soit $\mathcal{F}$ un faisceau d'ensembles sur $X$.
Il existe un ensemble $I$ et, pour tout $i \in I$, un élément
$U_i \in \mathcal{B}$ et un ensemble fini $S_i$ tels qu'il existe
une surjection $\coprod_{i \in I} j_{U_i!}\underline{S_i} \to \mathcal{F}$.
\end{lemma}

\begin{proof}
Soit $S$ un singleton. Nous démontrerons le résultat avec $S_i = S$.
Pour tout $x \in X$ et tout élément $s \in \mathcal{F}_x$, on peut choisir
$U(x, s) \in \mathcal{B}$ et $s(x, s) \in \mathcal{F}(U(x, s))$
dont l'image est $s$ dans $\mathcal{F}_x$. D'après
Faisceaux, lemme \ref{sheaves-lemma-j-shriek},
la section $s(x, s)$
correspond à une application de faisceaux $j_{U(x, s)!}\underline{S} \to \mathcal{F}$.
Alors
$$
\coprod\nolimits_{(x, s)} j_{U(x, s)!}\underline{S} \to \mathcal{F}
$$
est surjective sur les germes, donc surjective.
\end{proof}

\begin{lemma}
\label{lemma-filtered-colimit-constructibles}
Soit $X$ un espace topologique. Soit $\mathcal{B}$ une base de la
topologie de $X$, et supposons que tout $U \in \mathcal{B}$ soit quasi-compact.
Alors tout faisceau d'ensembles sur $X$ est une colimite filtrante de faisceaux de la forme
\begin{equation}
\label{equation-towards-constructible-sets}
\text{Coégalisateur}\left(
\xymatrix{
\coprod\nolimits_{b = 1, \ldots, m} j_{V_b!}\underline{S_b}
\ar@<1ex>[r] \ar@<-1ex>[r] &
\coprod\nolimits_{a = 1, \ldots, n} j_{U_a!}\underline{S_a}
}
\right)
\end{equation}
où $U_a$ et $V_b$ appartiennent à $\mathcal{B}$, et où $S_a$ et $S_b$ sont des ensembles finis.
\end{lemma}

\begin{proof}
Par le lemme \ref{lemma-surjection}, tout faisceau d'ensembles $\mathcal{F}$
est le but d'une surjection dont la source $\mathcal{F}_0$ est un coproduit
de faisceaux de la forme $j_{U!}\underline{S}$, avec $U \in \mathcal{B}$
et $S$ fini. En appliquant ceci à
$\mathcal{F}_0 \times_\mathcal{F} \mathcal{F}_0$
on trouve que $\mathcal{F}$ est le coégalisateur d'une paire d'applications
$$
\xymatrix{
\coprod\nolimits_{b \in B} j_{V_b!}\underline{S_b}
\ar@<1ex>[r] \ar@<-1ex>[r] &
\coprod\nolimits_{a \in A} j_{U_a!}\underline{S_a}
}
$$
pour certains ensembles d'indices $A$, $B$, avec $V_b$ et $U_a$ dans $\mathcal{B}$ et
$S_a$ et $S_b$ finis. Pour toute partie finie $B' \subset B$,
il existe une partie finie $A' \subset A$ telle que, par les deux applications, le coproduit
indexé par $b \in B'$ s'envoie dans le coproduit indexé par $a \in A'$.
En effet, on peut voir le membre de droite comme une colimite filtrante dont les
applications de transition sont injectives. Par conséquent, la prise des sections sur les ouverts
quasi-compacts $V_b$, $b \in B'$, commute à ce coproduit ;
voir Faisceaux, lemme \ref{sheaves-lemma-directed-colimits-sections}.
Ainsi notre faisceau est la colimite des conoyaux de ces applications
entre coproduits finis.
\end{proof}

\begin{lemma}
\label{lemma-constructible-comes-from-finite}
Soit $X$ un espace topologique spectral. Soit $\mathcal{B}$
l'ensemble des ouverts quasi-compacts de $X$.
Soit $\mathcal{F}$ un faisceau d'ensembles comme dans
l'équation (\ref{equation-towards-constructible-sets}).
Alors il existe une application spectrale continue $f : X \to Y$
vers un espace topologique sobre fini $Y$ et un faisceau
d'ensembles $\mathcal{G}$ sur $Y$, à germes finis,
tels que $f^{-1}\mathcal{G} \cong \mathcal{F}$.
\end{lemma}

\begin{proof}
On peut écrire $X = \lim X_i$ comme une limite projective
d'espaces sobres finis ; voir Topologie, lemme
\ref{topology-lemma-spectral-inverse-limit-finite-sober-spaces}.
Bien sûr, les applications de transition $X_{i'} \to X_i$ sont spectrales ; par conséquent,
d'après Topologie, lemme \ref{topology-lemma-directed-inverse-limit-spectral-spaces},
les applications $p_i : X \to X_i$ sont spectrales.
Pour un certain $i$, on peut trouver des ouverts $U_{a, i}$ et $V_{b, i}$
de $X_i$ dont les images inverses sont $U_a$ et $V_b$ ; voir
Topologie, lemme \ref{topology-lemma-descend-opens}.
Les deux applications
$$
\beta, \gamma :
\coprod\nolimits_{b \in B} j_{V_b!}\underline{S_b}
\longrightarrow
\coprod\nolimits_{a \in A} j_{U_a!}\underline{S_a}
$$
dont le coégalisateur est $\mathcal{F}$ correspondent par adjonction à deux familles
$$
\beta_b, \gamma_b :
S_b
\longrightarrow
\Gamma(V_b, \coprod\nolimits_{a \in A} j_{U_a!}\underline{S_a}), \quad
b \in B
$$
d'applications d'ensembles. Observons que
$p_i^{-1}(j_{U_{a, i}!}\underline{S_a}) = j_{U_a!}\underline{S_a}$
et $(X_{i'} \to X_i)^{-1}(j_{U_{a, i}!}\underline{S_a}) =
j_{U_{a, i'}!}\underline{S_a}$. Il résulte de
Faisceaux, lemme \ref{sheaves-lemma-compute-pullback-to-limit}
(en utilisant aussi que $S_b$ et $B$ sont des ensembles finis) que,
après avoir augmenté $i$, on trouve des applications
$$
\beta_{b, i}, \gamma_{b, i} :
S_b
\longrightarrow
\Gamma(V_{b, i}, \coprod\nolimits_{a \in A} j_{U_{a, i}!}\underline{S_a})
, \quad b \in B
$$
qui donnent les applications $\beta_b$ et $\gamma_b$ après image inverse
par $p_i$. Ces applications correspondent à leur tour à des applications de faisceaux
$$
\beta_i, \gamma_i :
\coprod\nolimits_{b \in B} j_{V_{b, i}!}\underline{S_b}
\longrightarrow
\coprod\nolimits_{a \in A} j_{U_{a, i}!}\underline{S_a}
$$
sur $X_i$. On peut alors prendre $Y = X_i$ et
$$
\mathcal{G} =
\text{Coégalisateur}\left(
\xymatrix{
\coprod\nolimits_{b = 1, \ldots, m} j_{V_{b, i}!}\underline{S_b}
\ar@<1ex>[r] \ar@<-1ex>[r] &
\coprod\nolimits_{a = 1, \ldots, n} j_{U_{a, i}!}\underline{S_a}
}
\right)
$$
Nous omettons quelques détails.
\end{proof}

\begin{lemma}
\label{lemma-constructible-in-constant}
Soit $X$ un espace topologique spectral. Soit $\mathcal{B}$
l'ensemble des ouverts quasi-compacts de $X$.
Soit $\mathcal{F}$ un faisceau d'ensembles comme dans
l'équation (\ref{equation-towards-constructible-sets}).
Alors il existe un nombre fini de fermés constructibles
$Z_1, \ldots, Z_n \subset X$ et d'ensembles finis $S_i$
tels que $\mathcal{F}$ soit isomorphe à un sous-faisceau de
$\prod (Z_i \to X)_*\underline{S_i}$.
\end{lemma}

\begin{proof}
Par le lemme \ref{lemma-constructible-comes-from-finite},
on se ramène au cas d'un espace topologique sobre fini
et d'un faisceau à germes finis.
Dans ce cas, $\mathcal{F} \subset \prod_{x \in X} i_{x, *}\mathcal{F}_x$,
où $i_x : \{x\} \to X$ est l'inclusion. Nous omettons la démonstration
du fait que $i_{x, *}\mathcal{F}_x$ est un faisceau constant sur $\overline{\{x\}}$.
\end{proof}
\section{Morphismes plats d'espaces annelés}
\label{section-flat-morphisms}

\noindent
La définition ponctuelle est motivée par le
Lemme \ref{lemma-flat-stalks-flat}
et la
Définition \ref{definition-flat-at-point}
ci-dessus.

\begin{definition}
\label{definition-flat-morphism}
Soit $f : X \to Y$ un morphisme d'espaces annelés.
Soit $x \in X$. Nous disons que $f$ est {\it plat en $x$}
si le morphisme d'anneaux $\mathcal{O}_{Y, f(x)} \to \mathcal{O}_{X, x}$ est plat.
Nous disons que $f$ est {\it plat} si $f$ est plat en tout $x \in X$.
\end{definition}

\noindent
Considérons le morphisme de faisceaux d'anneaux
$f^\sharp : f^{-1}\mathcal{O}_Y \to \mathcal{O}_X$.
Nous voyons que son germe en $x$ est le morphisme d'anneaux
$f^\sharp_x : \mathcal{O}_{Y, f(x)} \to \mathcal{O}_{X, x}$.
Ainsi, $f$ est plat en $x$ si et seulement si $\mathcal{O}_X$ est plat en $x$
comme $f^{-1}\mathcal{O}_Y$-module. De plus, $f$ est plat si et seulement si
$\mathcal{O}_X$ est plat comme $f^{-1}\mathcal{O}_Y$-module.
Les immersions ouvertes fournissent un cas particulier de morphismes plats.

\begin{lemma}
\label{lemma-pullback-flat}
Soit $f : X \to Y$ un morphisme plat d'espaces annelés.
Alors le foncteur image inverse
$f^* : \textit{Mod}(\mathcal{O}_Y) \to \textit{Mod}(\mathcal{O}_X)$
est exact.
\end{lemma}

\begin{proof}
Le foncteur $f^*$ est la composée du foncteur exact
$f^{-1} : \textit{Mod}(\mathcal{O}_Y) \to \textit{Mod}(f^{-1}\mathcal{O}_Y)$
et du foncteur de changement d'anneaux
$$
\textit{Mod}(f^{-1}\mathcal{O}_Y) \to \textit{Mod}(\mathcal{O}_X), \quad
\mathcal{F} \longmapsto
\mathcal{F} \otimes_{f^{-1}\mathcal{O}_Y} \mathcal{O}_X.
$$
Le résultat découle donc de la discussion qui suit la
Définition \ref{definition-flat-morphism}.
\end{proof}

\begin{definition}
\label{definition-flat-module}
Soit $f : (X, \mathcal{O}_X) \to (Y, \mathcal{O}_Y)$ un morphisme
d'espaces annelés. Soit $\mathcal{F}$ un faisceau de $\mathcal{O}_X$-modules.
\begin{enumerate}
\item Nous disons que $\mathcal{F}$ est {\it plat sur $Y$ en un point $x \in X$}
si le germe $\mathcal{F}_x$ est un $\mathcal{O}_{Y, f(x)}$-module plat.
\item Nous disons que $\mathcal{F}$ est {\it plat sur $Y$} si
$\mathcal{F}$ est plat sur $Y$ en tout point $x$ de $X$.
\end{enumerate}
\end{definition}

\noindent
Avec cette définition, nous voyons que $\mathcal{F}$ est plat sur $Y$ en $x$
si et seulement si $\mathcal{F}$ est plat en $x$ comme
$f^{-1}\mathcal{O}_Y$-module, car
$(f^{-1}\mathcal{O}_Y)_x = \mathcal{O}_{Y, f(x)}$ d'après
Faisceaux, lemme \ref{sheaves-lemma-stalk-pullback}.

\begin{lemma}
\label{lemma-pullback-tensor-flat-module}
Soit $f : X \to Y$ un morphisme d'espaces annelés. Soit $\mathcal{F}$
un $\mathcal{O}_X$-module plat sur $Y$. Alors le foncteur
$$
\textit{Mod}(\mathcal{O}_Y) \to \textit{Mod}(\mathcal{O}_X),\quad
\mathcal{G} \longmapsto f^*\mathcal{G} \otimes_{\mathcal{O}_X} \mathcal{F}
$$
est exact.
\end{lemma}

\begin{proof}
Cela est vrai parce que
$f^*\mathcal{G} \otimes_{\mathcal{O}_X} \mathcal{F} =
f^{-1}\mathcal{G} \otimes_{f^{-1}\mathcal{O}_Y} \mathcal{F}$,
le foncteur $f^{-1}$ est exact et $\mathcal{F}$ est un
$f^{-1}\mathcal{O}_Y$-module plat.
\end{proof}














\section{Puissances symétriques et extérieures}
\label{section-symmetric-exterior}

\noindent
Soit $(X, \mathcal{O}_X)$ un espace annelé.
Soit $\mathcal{F}$ un $\mathcal{O}_X$-module.
Nous définissons l'{\it algèbre tensorielle} de $\mathcal{F}$ comme
le faisceau de $\mathcal{O}_X$-algèbres non commutatives
$$
\text{T}(\mathcal{F})
=
\text{T}_{\mathcal{O}_X}(\mathcal{F})
= \bigoplus\nolimits_{n \geq 0} \text{T}^n(\mathcal{F}).
$$
Ici, $\text{T}^0(\mathcal{F}) = \mathcal{O}_X$,
$\text{T}^1(\mathcal{F}) = \mathcal{F}$
et, pour $n \geq 2$, nous avons
$$
\text{T}^n(\mathcal{F}) =
\mathcal{F} \otimes_{\mathcal{O}_X} \ldots \otimes_{\mathcal{O}_X} \mathcal{F}
\ \ (n\text{ facteurs})
$$
Nous définissons $\wedge(\mathcal{F})$ comme le quotient de
$\text{T}(\mathcal{F})$ par l'idéal bilatère engendré par les
sections locales $s \otimes s$ de $\text{T}^2(\mathcal{F})$, où
$s$ est une section locale de $\mathcal{F}$. On l'appelle
l'{\it algèbre extérieure} de $\mathcal{F}$.
De même, nous définissons $\text{Sym}(\mathcal{F})$ comme
le quotient de $\text{T}(\mathcal{F})$ par l'idéal
bilatère engendré par les sections locales de la forme
$s \otimes t - t \otimes s$ de $\text{T}^2(\mathcal{F})$.

\medskip\noindent
Les algèbres $\text{T}(\mathcal{F})$, $\wedge(\mathcal{F})$ et
$\text{Sym}(\mathcal{F})$ sont des $\mathcal{O}_X$-algèbres graduées
au sens de Faisceaux différentiels gradués, définition
\ref{sdga-definition-ga}.
De plus, $\text{Sym}(\mathcal{F})$
est commutative et $\wedge(\mathcal{F})$ est commutative au sens gradué.

\begin{lemma}
\label{lemma-local-tensor-algebra}
Dans la situation décrite ci-dessus,
le faisceau $\wedge^n\mathcal{F}$ est le faisceau associé au
préfaisceau
$$
U \longmapsto \wedge^n_{\mathcal{O}_X(U)}(\mathcal{F}(U)).
$$
Voir Algèbre, section \ref{algebra-section-tensor-algebra}.
De même, le faisceau $\text{Sym}^n\mathcal{F}$ est le faisceau associé
au préfaisceau
$$
U \longmapsto \text{Sym}^n_{\mathcal{O}_X(U)}(\mathcal{F}(U)).
$$
\end{lemma}

\begin{proof}
Omis. Il peut être plus efficace de définir $\text{Sym}(\mathcal{F})$
et $\wedge(\mathcal{F})$ de cette manière plutôt que par la méthode
donnée ci-dessus.
\end{proof}

\begin{lemma}
\label{lemma-stalk-tensor-algebra}
Dans la situation décrite ci-dessus, soit $x \in X$.
Il existe des isomorphismes canoniques de $\mathcal{O}_{X, x}$-modules
$\text{T}(\mathcal{F})_x = \text{T}(\mathcal{F}_x)$,
$\text{Sym}(\mathcal{F})_x = \text{Sym}(\mathcal{F}_x)$ et
$\wedge(\mathcal{F})_x = \wedge(\mathcal{F}_x)$.
\end{lemma}

\begin{proof}
Cela résulte immédiatement du Lemme \ref{lemma-local-tensor-algebra} ci-dessus et de
Algèbre, lemme \ref{algebra-lemma-colimit-tensor-algebra}.
\end{proof}

\begin{lemma}
\label{lemma-pullback-tensor-algebra}
Soit $f : (X, \mathcal{O}_X) \to (Y, \mathcal{O}_Y)$ un morphisme
d'espaces annelés. Soit $\mathcal{F}$ un faisceau de $\mathcal{O}_Y$-modules.
Alors $f^*\text{T}(\mathcal{F}) = \text{T}(f^*\mathcal{F})$,
et il en va de même des algèbres extérieure et symétrique associées
à $\mathcal{F}$.
\end{lemma}

\begin{proof}
Omis.
\end{proof}

\begin{lemma}
\label{lemma-presentation-sym-exterior}
Soit $(X, \mathcal{O}_X)$ un espace annelé.
Soit $\mathcal{F}_2 \to \mathcal{F}_1 \to \mathcal{F} \to 0$
une suite exacte de faisceaux de $\mathcal{O}_X$-modules.
Pour tout $n \geq 1$, il existe une suite exacte
$$
\mathcal{F}_2 \otimes_{\mathcal{O}_X} \text{Sym}^{n - 1}(\mathcal{F}_1)
\to
\text{Sym}^n(\mathcal{F}_1)
\to
\text{Sym}^n(\mathcal{F})
\to
0
$$
et, de même, une suite exacte
$$
\mathcal{F}_2 \otimes_{\mathcal{O}_X} \wedge^{n - 1}(\mathcal{F}_1)
\to
\wedge^n(\mathcal{F}_1)
\to
\wedge^n(\mathcal{F})
\to
0
$$
\end{lemma}

\begin{proof}
Voir Algèbre, lemme \ref{algebra-lemma-presentation-sym-exterior}.
\end{proof}

\begin{lemma}
\label{lemma-tensor-algebra-permanence}
Soit $(X, \mathcal{O}_X)$ un espace annelé.
Soit $\mathcal{F}$ un faisceau de $\mathcal{O}_X$-modules.
\begin{enumerate}
\item Si $\mathcal{F}$ est localement engendré par des sections,
il en va de même de chaque $\text{T}^n(\mathcal{F})$,
$\wedge^n(\mathcal{F})$ et $\text{Sym}^n(\mathcal{F})$.
\item Si $\mathcal{F}$ est de type fini,
il en va de même de chaque $\text{T}^n(\mathcal{F})$,
$\wedge^n(\mathcal{F})$ et $\text{Sym}^n(\mathcal{F})$.
\item Si $\mathcal{F}$ est de présentation finie,
il en va de même de chaque $\text{T}^n(\mathcal{F})$,
$\wedge^n(\mathcal{F})$ et $\text{Sym}^n(\mathcal{F})$.
\item Si $\mathcal{F}$ est cohérent,
alors, pour $n > 0$, chacun de $\text{T}^n(\mathcal{F})$,
$\wedge^n(\mathcal{F})$ et $\text{Sym}^n(\mathcal{F})$
est cohérent.
\item Si $\mathcal{F}$ est quasi-cohérent,
il en va de même de chaque $\text{T}^n(\mathcal{F})$,
$\wedge^n(\mathcal{F})$ et $\text{Sym}^n(\mathcal{F})$.
\item Si $\mathcal{F}$ est localement libre,
il en va de même de chaque $\text{T}^n(\mathcal{F})$,
$\wedge^n(\mathcal{F})$ et $\text{Sym}^n(\mathcal{F})$.
\end{enumerate}
\end{lemma}

\begin{proof}
Ces assertions pour $\text{T}^n(\mathcal{F})$ résultent du
Lemme \ref{lemma-tensor-product-permanence}.

\medskip\noindent
Les assertions (1) et (2) résultent du fait que
$\wedge^n(\mathcal{F})$ et $\text{Sym}^n(\mathcal{F})$
sont des quotients de $\text{T}^n(\mathcal{F})$.

\medskip\noindent
L'assertion (6) résulte de
Algèbre, lemme \ref{algebra-lemma-free-tensor-algebra}.

\medskip\noindent
Pour (3) et (5), nous utiliserons le
Lemme \ref{lemma-presentation-sym-exterior} ci-dessus.
En choisissant localement une présentation
$\mathcal{F}_2 \to \mathcal{F}_1 \to \mathcal{F} \to 0$
où les $\mathcal{F}_i$ sont libres, ou libres de type fini, puis en appliquant le
lemme, nous voyons que $\text{Sym}^n(\mathcal{F})$ et $\wedge^n(\mathcal{F})$
admettent une présentation analogue ; nous utilisons ici (6) et le
Lemme \ref{lemma-tensor-product-permanence}.

\medskip\noindent
Pour démontrer (4), nous utiliserons
Algèbre, lemme \ref{algebra-lemma-present-sym-wedge}.
Nous pouvons localiser sur $X$ et supposer que
$\mathcal{F}$ est engendré par une famille finie
$(s_i)_{i \in I}$ de sections globales.
Le lemme mentionné ci-dessus,
combiné au Lemme \ref{lemma-local-tensor-algebra} ci-dessus,
implique que, pour $n \geq 2$,
il existe une suite exacte
$$
\bigoplus\nolimits_{j \in J}
\text{T}^{n - 2}(\mathcal{F})
\to
\text{T}^n(\mathcal{F})
\to
\text{Sym}^n(\mathcal{F})
\to
0
$$
où l'ensemble d'indices $J$ est fini. Or nous savons que
$\text{T}^{n - 2}(\mathcal{F})$ est de type fini ;
l'image de la première flèche est donc un sous-faisceau cohérent
de $\text{T}^n(\mathcal{F})$ ; voir le Lemme \ref{lemma-coherent-abelian}.
Ce même lemme permet de conclure que $\text{Sym}^n(\mathcal{F})$ est
cohérent.
\end{proof}

\begin{lemma}
\label{lemma-whole-tensor-algebra-permanence}
Soit $(X, \mathcal{O}_X)$ un espace annelé.
Soit $\mathcal{F}$ un faisceau de $\mathcal{O}_X$-modules.
\begin{enumerate}
\item Si $\mathcal{F}$ est quasi-cohérent,
il en va de même de chacun de $\text{T}(\mathcal{F})$,
$\wedge(\mathcal{F})$ et $\text{Sym}(\mathcal{F})$.
\item Si $\mathcal{F}$ est localement libre,
il en va de même de chacun de $\text{T}(\mathcal{F})$,
$\wedge(\mathcal{F})$ et $\text{Sym}(\mathcal{F})$.
\end{enumerate}
\end{lemma}

\begin{proof}
Il n'est pas vrai qu'une somme directe infinie $\bigoplus \mathcal{G}_i$ de
modules localement libres soit localement libre, ni qu'une
somme directe infinie de modules quasi-cohérents
soit quasi-cohérente. Le problème est que, pour un
point $x \in X$, les voisinages ouverts $U_i$ de $x$ sur lesquels $\mathcal{G}_i$
devient libre (resp.\ admet une présentation convenable) peuvent avoir une intersection
qui n'est pas un voisinage ouvert de $x$. Cependant, dans la
démonstration du Lemme \ref{lemma-tensor-algebra-permanence}, nous avons vu que,
dès qu'un voisinage ouvert convenable pour $\mathcal{F}$ est choisi,
ce voisinage ouvert convient à chacun des faisceaux
$\text{T}^n(\mathcal{F})$, $\wedge^n(\mathcal{F})$ et
$\text{Sym}^n(\mathcal{F})$.
Le lemme en résulte.
\end{proof}







\section{Hom interne}
\label{section-internal-hom}

\noindent
Soit $(X, \mathcal{O}_X)$ un espace annelé.
Soient $\mathcal{F}$ et $\mathcal{G}$ des $\mathcal{O}_X$-modules.
Considérons la règle
$$
U \longmapsto \Hom_{\mathcal{O}_X|_U}(\mathcal{F}|_U, \mathcal{G}|_U).
$$
Il résulte de la discussion de Faisceaux, section
\ref{sheaves-section-glueing-sheaves}, qu'il s'agit d'un faisceau de
groupes abéliens. De plus, étant donnés un élément
$\varphi \in \Hom_{\mathcal{O}_X|_U}(\mathcal{F}|_U, \mathcal{G}|_U)$
et une section $f \in \mathcal{O}_X(U)$, nous pouvons définir
$f\varphi \in \Hom_{\mathcal{O}_X|_U}(\mathcal{F}|_U, \mathcal{G}|_U)$
soit en précomposant avec la multiplication par $f$ sur $\mathcal{F}|_U$,
soit en postcomposant avec la multiplication par $f$ sur $\mathcal{G}|_U$ (le résultat
est le même). Nous obtenons donc en fait un faisceau de $\mathcal{O}_X$-modules.
Nous noterons ce faisceau
$\SheafHom_{\mathcal{O}_X}(\mathcal{F}, \mathcal{G})$.
Il existe un morphisme canonique d'``évaluation''
$$
\mathcal{F}
\otimes_{\mathcal{O}_X}
\SheafHom_{\mathcal{O}_X}(\mathcal{F}, \mathcal{G})
\longrightarrow
\mathcal{G}.
$$
Pour tout $x \in X$, il existe aussi un morphisme canonique
$$
\SheafHom_{\mathcal{O}_X}(\mathcal{F}, \mathcal{G})_x
\to
\Hom_{\mathcal{O}_{X, x}}(\mathcal{F}_x, \mathcal{G}_x)
$$
qui est rarement un isomorphisme.

\begin{lemma}
\label{lemma-internal-hom}
Soit $(X, \mathcal{O}_X)$ un espace annelé.
Soient $\mathcal{F}$, $\mathcal{G}$ et $\mathcal{H}$ des $\mathcal{O}_X$-modules.
Il existe un isomorphisme canonique
$$
\SheafHom_{\mathcal{O}_X}
(\mathcal{F} \otimes_{\mathcal{O}_X} \mathcal{G}, \mathcal{H})
\longrightarrow
\SheafHom_{\mathcal{O}_X}
(\mathcal{F}, \SheafHom_{\mathcal{O}_X}(\mathcal{G}, \mathcal{H}))
$$
qui est fonctoriel en ses trois arguments (il s'agit du Hom de faisceaux
aux trois emplacements). En particulier, se donner un
morphisme $\mathcal{F} \otimes_{\mathcal{O}_X} \mathcal{G} \to \mathcal{H}$
équivaut à se donner un morphisme
$\mathcal{F} \to \SheafHom_{\mathcal{O}_X}(\mathcal{G}, \mathcal{H})$.
\end{lemma}

\begin{proof}
C'est l'analogue de
Algèbre, lemme \ref{algebra-lemma-hom-from-tensor-product}.
La démonstration, identique, est omise.
\end{proof}

\begin{lemma}
\label{lemma-internal-hom-exact}
Soit $(X, \mathcal{O}_X)$ un espace annelé.
Soient $\mathcal{F}$ et $\mathcal{G}$ des $\mathcal{O}_X$-modules.
\begin{enumerate}
\item Si $\mathcal{F}_2 \to \mathcal{F}_1 \to \mathcal{F} \to 0$
est une suite exacte de $\mathcal{O}_X$-modules, alors
$$
0 \to
\SheafHom_{\mathcal{O}_X}(\mathcal{F}, \mathcal{G}) \to
\SheafHom_{\mathcal{O}_X}(\mathcal{F}_1, \mathcal{G}) \to
\SheafHom_{\mathcal{O}_X}(\mathcal{F}_2, \mathcal{G})
$$
est exacte.
\item Si $0 \to \mathcal{G} \to \mathcal{G}_1 \to \mathcal{G}_2$
est une suite exacte de $\mathcal{O}_X$-modules, alors
$$
0 \to
\SheafHom_{\mathcal{O}_X}(\mathcal{F}, \mathcal{G}) \to
\SheafHom_{\mathcal{O}_X}(\mathcal{F}, \mathcal{G}_1) \to
\SheafHom_{\mathcal{O}_X}(\mathcal{F}, \mathcal{G}_2)
$$
est exacte.
\end{enumerate}
\end{lemma}

\begin{proof}
Soit $\mathcal{F}_2 \to \mathcal{F}_1 \to \mathcal{F} \to 0$
comme en (1). Pour tout ouvert $U \subset X$, la suite
$$
0 \to \Hom_{\mathcal{O}_U}(\mathcal{F}|_U, \mathcal{G}|_U) \to
\Hom_{\mathcal{O}_U}(\mathcal{F}_1|_U, \mathcal{G}|_U) \to
\Hom_{\mathcal{O}_U}(\mathcal{F}_2|_U, \mathcal{G}|_U)
$$
est exacte d'après Homologie, lemme \ref{homology-lemma-check-exactness}.
Cela signifie que
prendre les sections sur $U$ de la suite de faisceaux de (1) donne
une suite exacte de groupes abéliens. La suite de (1) est donc exacte
par définition. La démonstration de (2) est exactement la même.
\end{proof}

\begin{lemma}
\label{lemma-adjoint-tensor-restrict}
Soit $X$ un espace topologique. Soit $\mathcal{O}_1 \to \mathcal{O}_2$
un morphisme de faisceaux d'anneaux. Alors nous avons
$$
\Hom_{\mathcal{O}_1}(\mathcal{F}_{\mathcal{O}_1}, \mathcal{G}) =
\Hom_{\mathcal{O}_2}(\mathcal{F},
\SheafHom_{\mathcal{O}_1}(\mathcal{O}_2, \mathcal{G}))
$$
bifonctoriellement en $\mathcal{F} \in \textit{Mod}(\mathcal{O}_2)$
et $\mathcal{G} \in \textit{Mod}(\mathcal{O}_1)$.
\end{lemma}

\begin{proof}
Omis. C'est l'analogue de
Algèbre, lemme \ref{algebra-lemma-adjoint-hom-restrict},
et la démonstration est exactement la même.
\end{proof}

\begin{lemma}
\label{lemma-stalk-internal-hom}
Soit $(X, \mathcal{O}_X)$ un espace annelé.
Soient $\mathcal{F}$ et $\mathcal{G}$ des $\mathcal{O}_X$-modules.
Si $\mathcal{F}$ est de type fini, alors le morphisme canonique
$$
\SheafHom_{\mathcal{O}_X}(\mathcal{F}, \mathcal{G})_x
\to
\Hom_{\mathcal{O}_{X, x}}(\mathcal{F}_x, \mathcal{G}_x)
$$
est injectif.
Si $\mathcal{F}$ est de présentation finie,
ce morphisme canonique est un isomorphisme.
\end{lemma}

\begin{proof}
Le morphisme envoie la classe d'équivalence de
$(U, \varphi)$ dans $\SheafHom_{\mathcal{O}_X}(\mathcal{F}, \mathcal{G})_x$,
où $x \in U \subset X$, l'ensemble intermédiaire étant ouvert, et
$\varphi \in \Hom_{\mathcal{O}_U}(\mathcal{F}|_U, \mathcal{G}|_U)$,
sur le morphisme induit sur les germes en $x$, à savoir
$\varphi_x : \mathcal{F}_x \to \mathcal{G}_x$.

\medskip\noindent
Supposons $\mathcal{F}$ de type fini.
Choisissons un représentant $(U, \varphi)$ d'un élément $\sigma$
du noyau du morphisme, c'est-à-dire tel que $\varphi_x = 0$.
En rétrécissant $U$ si nécessaire, choisissons des sections
$s^1, \ldots, s^n \in \mathcal{F}(U)$ qui engendrent $\mathcal{F}|_U$.
Comme $\varphi_x(s^i_x) = 0$
et que nous n'avons qu'un nombre fini de sections,
nous pouvons trouver un voisinage ouvert
$V \subset U$ de $x$ tel que $\varphi_V(s^i|_V)=0$
pour tout $i = 1, \ldots, n$. Puisque les $s^i|_V$, $i = 1, \ldots, n$, engendrent
$\mathcal{F}|_V$, cela signifie que $\varphi|_V = 0$.
Comme $(U, \varphi)$ est équivalent à $(V, \varphi|_V)$,
nous concluons que $\sigma = 0$, d'où l'injectivité du morphisme.

\medskip\noindent
Supposons maintenant $\mathcal{F}$ de présentation finie.
En localisant sur $X$, nous pouvons supposer que $\mathcal{F}$ admet une présentation
$$
\bigoplus\nolimits_{j = 1, \ldots, m}
\mathcal{O}_X
\longrightarrow
\bigoplus\nolimits_{i = 1, \ldots, n}
\mathcal{O}_X
\to
\mathcal{F}
\to
0.
$$
D'après le Lemme \ref{lemma-internal-hom-exact}, cela donne une suite exacte
$
0 \to
\SheafHom_{\mathcal{O}_X}(\mathcal{F}, \mathcal{G}) \to
\bigoplus\nolimits_{i = 1, \ldots, n} \mathcal{G}
\longrightarrow
\bigoplus\nolimits_{j = 1, \ldots, m} \mathcal{G}.
$
En prenant les germes, nous obtenons une suite exacte
$
0 \to
\SheafHom_{\mathcal{O}_X}(\mathcal{F}, \mathcal{G})_x \to
\bigoplus\nolimits_{i = 1, \ldots, n} \mathcal{G}_x
\longrightarrow
\bigoplus\nolimits_{j = 1, \ldots, m} \mathcal{G}_x
$
et le résultat en découle puisque $\mathcal{F}_x$ s'insère dans
une suite exacte
$
\bigoplus\nolimits_{j = 1, \ldots, m}
\mathcal{O}_{X, x}
\longrightarrow
\bigoplus\nolimits_{i = 1, \ldots, n}
\mathcal{O}_{X, x}
\to
\mathcal{F}_x
\to
0
$
qui induit la suite exacte
$
0 \to
\Hom_{\mathcal{O}_{X, x}}(\mathcal{F}_x, \mathcal{G}_x) \to
\bigoplus\nolimits_{i = 1, \ldots, n} \mathcal{G}_x
\longrightarrow
\bigoplus\nolimits_{j = 1, \ldots, m} \mathcal{G}_x
$
qui est la même que celle ci-dessus.
\end{proof}

\begin{lemma}
\label{lemma-pullback-internal-hom}
Soit $f : (X, \mathcal{O}_X) \to (Y, \mathcal{O}_Y)$ un morphisme
d'espaces annelés. Soient $\mathcal{F}$ et $\mathcal{G}$ des $\mathcal{O}_Y$-modules.
Si $\mathcal{F}$ est de présentation finie et si $f$ est plat,
alors le morphisme canonique
$$
f^*\SheafHom_{\mathcal{O}_Y}(\mathcal{F}, \mathcal{G})
\longrightarrow
\SheafHom_{\mathcal{O}_X}(f^*\mathcal{F}, f^*\mathcal{G})
$$
est un isomorphisme.
\end{lemma}

\begin{proof}
Remarquons que $f^*\mathcal{F}$ est aussi de présentation finie
(Lemme \ref{lemma-pullback-finite-presentation}).
Soit $x \in X$ d'image $y \in Y$. En considérant les germes
en $x$, nous obtenons un isomorphisme d'après le
Lemme \ref{lemma-stalk-internal-hom} et
Compléments d'algèbre, lemme
\ref{more-algebra-lemma-pseudo-coherence-and-base-change-ext}
qui montre que, dans ce cas, $\Hom$ commute au changement de base par
$\mathcal{O}_{Y, y} \to \mathcal{O}_{X, x}$.
Seconde démonstration : utiliser exactement le même argument que dans
la démonstration du Lemme \ref{lemma-stalk-internal-hom}.
\end{proof}

\begin{lemma}
\label{lemma-internal-hom-locally-kernel-direct-sum}
Soit $(X, \mathcal{O}_X)$ un espace annelé.
Soient $\mathcal{F}$ et $\mathcal{G}$ des $\mathcal{O}_X$-modules.
Si $\mathcal{F}$ est de présentation finie, alors le faisceau
$\SheafHom_{\mathcal{O}_X}(\mathcal{F}, \mathcal{G})$
est localement le noyau d'un morphisme entre des sommes directes finies
de copies de $\mathcal{G}$.
En particulier, si $\mathcal{G}$ est cohérent, alors
$\SheafHom_{\mathcal{O}_X}(\mathcal{F}, \mathcal{G})$
est lui aussi cohérent.
\end{lemma}

\begin{proof}
Nous avons vu la première assertion
dans la démonstration du Lemme \ref{lemma-stalk-internal-hom}.
Le résultat pour les faisceaux cohérents découle alors du
Lemme \ref{lemma-coherent-abelian}.
\end{proof}

\begin{lemma}
\label{lemma-hom-finite-presentation-colimit}
Soit $X$ un espace annelé. Soit $\mathcal{F}$ un $\mathcal{O}_X$-module
de présentation finie. Soit
$\mathcal{G} = \colim_{\lambda \in \Lambda} \mathcal{G}_\lambda$
une colimite filtrante de $\mathcal{O}_X$-modules. Alors le morphisme canonique
$$
\colim_\lambda \SheafHom_{\mathcal{O}_X}(\mathcal{F}, \mathcal{G}_\lambda)
\longrightarrow
\SheafHom_{\mathcal{O}_X}(\mathcal{F}, \mathcal{G})
$$
est un isomorphisme.
\end{lemma}

\begin{proof}
La formation des colimites de faisceaux de modules commute à la restriction aux ouverts ; voir
Faisceaux, section \ref{sheaves-section-limits-sheaves}. Nous pouvons donc supposer que
$\mathcal{F}$ admet une présentation globale
$$
\bigoplus\nolimits_{j = 1, \ldots, m}
\mathcal{O}_X
\longrightarrow
\bigoplus\nolimits_{i = 1, \ldots, n}
\mathcal{O}_X
\to
\mathcal{F}
\to
0
$$
Le foncteur $\SheafHom_{\mathcal{O}_X}(-, -)$ commute aux sommes directes finies
en chacune des variables et $\SheafHom_{\mathcal{O}_X}(\mathcal{O}_X, -)$
est le foncteur identité. D'après cela et le Lemme \ref{lemma-internal-hom-exact},
nous obtenons une suite exacte
$$
0 \to
\SheafHom_{\mathcal{O}_X}(\mathcal{F}, \mathcal{G}) \to
\bigoplus\nolimits_{i = 1, \ldots, n} \mathcal{G} \to
\bigoplus\nolimits_{j = 1, \ldots, m} \mathcal{G}
$$
Puisque les colimites filtrantes sont exactes dans $\textit{Mod}(\mathcal{O}_X)$,
la ligne supérieure du diagramme commutatif suivant est elle aussi exacte
$$
\xymatrix{
0 \ar[r] &
\colim_\lambda \SheafHom_{\mathcal{O}_X}(\mathcal{F}, \mathcal{G}_\lambda)
\ar[r] \ar[d] &
\colim_\lambda \bigoplus\nolimits_{i = 1, \ldots, n} \mathcal{G}_\lambda
\ar[r] \ar[d] &
\colim_\lambda \bigoplus\nolimits_{j = 1, \ldots, m} \mathcal{G}_\lambda
\ar[d] \\
0 \ar[r] &
\SheafHom_{\mathcal{O}_X}(\mathcal{F}, \mathcal{G}) \ar[r] &
\bigoplus\nolimits_{i = 1, \ldots, n} \mathcal{G} \ar[r] &
\bigoplus\nolimits_{j = 1, \ldots, m} \mathcal{G}
}
$$
Comme les deux flèches verticales de droite sont des isomorphismes, nous concluons.
\end{proof}

\begin{lemma}
\label{lemma-finite-presentation-quasi-compact-colimit}
Soit $X$ un espace annelé.
Soit $I$ un ensemble préordonné et
soit $(\mathcal{F}_i, \varphi_{ii'})$ un système sur $I$
formé de faisceaux de $\mathcal{O}_X$-modules
(voir Catégories, section \ref{categories-section-posets-limits}).
Supposons que
\begin{enumerate}
\item $I$ est filtrant,
\item $\mathcal{G}$ est un $\mathcal{O}_X$-module de présentation finie, et
\item $X$ admet un système cofinal de recouvrements ouverts
$\mathcal{U} : X = \bigcup_{j\in J} U_j$, avec
$J$ fini et $U_j \cap U_{j'}$ quasi-compact
pour tous $j, j' \in J$.
\end{enumerate}
Alors nous avons
$$
\colim_i \Hom_X(\mathcal{G}, \mathcal{F}_i)
=
\Hom_X(\mathcal{G}, \colim_i \mathcal{F}_i).
$$
\end{lemma}

\begin{proof}
Posons
$\mathcal{H} = \SheafHom_{\mathcal{O}_X}(\mathcal{G}, \colim \mathcal{F}_i)$
et $\mathcal{H}_i = \SheafHom_{\mathcal{O}_X}(\mathcal{G}, \mathcal{F}_i)$.
Rappelons que
$$
\Hom_X(\mathcal{G}, \mathcal{F}) = \Gamma(X, \mathcal{H})
\quad\text{et}\quad
\Hom_X(\mathcal{G}, \mathcal{F}_i) = \Gamma(X, \mathcal{H}_i)
$$
par construction. D'après le Lemme \ref{lemma-hom-finite-presentation-colimit}, nous avons
$\mathcal{H} = \colim \mathcal{H}_i$. Le lemme résulte donc de
Faisceaux, lemme \ref{sheaves-lemma-directed-colimits-sections}.
\end{proof}

\begin{remark}
\label{remark-condition-necessary}
Dans le lemme ci-dessus, il faut une condition supplémentaire à la
quasi-compacité de $X$. Voir
Faisceaux, exemple \ref{sheaves-example-conditions-needed-colimit}.
\end{remark}




\section{L'annulateur d'un faisceau de modules}
\label{section-annihilator}

\noindent
Soit $(X, \mathcal{O}_X)$ un espace annelé.
Soit $\mathcal{F}$ un $\mathcal{O}_X$-module.
Il existe un morphisme canonique de faisceaux de $\mathcal{O}_X$-modules
$$
\mathcal{O}_X
\longrightarrow
\SheafHom_{\mathcal{O}_X}(\mathcal{F}, \mathcal{F})
$$
qui envoie une section locale $f \in \mathcal{O}_X(U)$ sur le morphisme
$f : \mathcal{F}|_U \to \mathcal{F}|_U$ donné par la multiplication par $f$.

\begin{definition}
\label{definition-annihilator-sheaf}
Soit $(X, \mathcal{O}_X)$ un espace annelé et soit $\mathcal{F}$
un $\mathcal{O}_X$-module. L'{\it annulateur} de $\mathcal{F}$,
noté $\text{Ann}_{\mathcal{O}_X}(\mathcal{F})$,
est le noyau du morphisme
$\mathcal{O}_X \to \SheafHom_{\mathcal{O}_X}(\mathcal{F}, \mathcal{F})$
considéré ci-dessus.
\end{definition}

\noindent
Pour tout $x\in X$, on a une inclusion d'idéaux de $\mathcal{O}_{X, x}$ :
\begin{equation}
\label{equation-inclusion-of-annihilator-ideals}
(\text{Ann}_{\mathcal{O}_X}(\mathcal{F}))_x
\subset
\text{Ann}_{\mathcal{O}_{X, x}}(\mathcal{F}_x)
\end{equation}
car toute section de $\text{Ann}_{\mathcal{O}_X}(\mathcal{F})$
annule les germes de $\mathcal{F}$ en tout point où elle est
définie. Voici une situation simple dans laquelle
(\ref{equation-inclusion-of-annihilator-ideals}) est une égalité.

\begin{lemma}
\label{lemma-finite-type-annihilator-stalk}
Soit $(X, \mathcal{O}_X)$ un espace annelé et soit
$\mathcal{F}$ un faisceau de $\mathcal{O}_X$-modules.
Si $\mathcal{F}$ est de type fini, alors
$(\text{Ann}_{\mathcal{O}_X}(\mathcal{F}))_x =
\text{Ann}_{\mathcal{O}_{X, x}}(\mathcal{F}_x)$.
\end{lemma}

\begin{proof}
D'après le lemme \ref{lemma-stalk-internal-hom}, le morphisme
$$
\SheafHom_{\mathcal{O}_X}(\mathcal{F}, \mathcal{F})_x \longrightarrow
\Hom_{\mathcal{O}_{X,x}}(\mathcal{F}_x, \mathcal{F}_x)
$$
est injectif. Ainsi, toute section $f$ de $\mathcal{O}_X$ sur un voisinage
ouvert $U$ de $x$ qui agit par zéro sur $\mathcal{F}_x$
agit par zéro sur $\mathcal{F}|_V$ pour un ouvert $U \supset V \ni x$ convenable.
L'inclusion (\ref{equation-inclusion-of-annihilator-ideals})
est donc une égalité.
\end{proof}

\begin{lemma}
\label{lemma-induced-quotient-ring-sheaf-module-structure}
Soit $(X, \mathcal{O}_X)$ un espace annelé, soit $\mathcal{F}$ un
$\mathcal{O}_X$-module et soit $\mathcal{I} \subset \mathcal{O}_X$
un faisceau d'idéaux. Si
$\mathcal{I} \subset \text{Ann}_{\mathcal{O}_X}(\mathcal{F})$,
alors $\mathcal{F}$ possède une structure naturelle de
$\mathcal{O}_X/\mathcal{I}$-module qui, sur les germes, coïncide avec
la construction usuelle d'algèbre commutative.
\end{lemma}

\begin{proof}
En appliquant la propriété universelle du conoyau de l'inclusion
$\mathcal{I} \to \mathcal{O}_X$, nous obtenons un diagramme commutatif
$$
\xymatrix{
\mathcal{O}_X \ar[r] \ar[d] &
\SheafHom_{\mathcal{O}_X}(\mathcal{F}, \mathcal{F}) \\
\mathcal{O}_X/\mathcal{I} \ar@{-->}[ur]
}
$$
de $\mathcal{O}_X$-modules. D'après le lemme \ref{lemma-internal-hom},
le morphisme obtenu $\mathcal{O}_X/\mathcal{I} \to
\SheafHom_{\mathcal{O}_X}(\mathcal{F}, \mathcal{F})$
correspond à un morphisme de $\mathcal{O}_X$-modules
$$
\mathcal{O}_X/\mathcal{I} \otimes_{\mathcal{O}_X} \mathcal{F}
\longrightarrow
\mathcal{F}
$$
ce qui nous donne une structure de $\mathcal{O}_X/\mathcal{I}$-module
sur $\mathcal{F}$, compatible avec sa structure donnée de $\mathcal{O}_X$-module.
Nous omettons la vérification de l'assertion sur les germes.
\end{proof}

\begin{lemma}
\label{lemma-coherent-annihilator}
Soit $(X,\mathcal{O}_X)$ un espace annelé.
Si $\mathcal{O}_X$ et $\mathcal{F}$ sont cohérents,
alors $\text{Ann}_{\mathcal{O}_X}(\mathcal{F})$ l'est aussi.
\end{lemma}

\begin{proof}
Puisque $\text{Ann}_{\mathcal{O}_X}(\mathcal{F})$ est le noyau de
$\mathcal{O}_X \to \SheafHom_{\mathcal{O}_X}(\mathcal{F}, \mathcal{F})$
et d'après le lemme \ref{lemma-coherent-abelian}, il suffit
de montrer que $\SheafHom_{\mathcal{O}_X}(\mathcal{F}, \mathcal{F})$ est
cohérent. Cela résulte du
lemme \ref{lemma-internal-hom-locally-kernel-direct-sum}
et du fait que $\mathcal{F}$ est cohérent, donc a fortiori de présentation finie
(lemme \ref{lemma-coherent-finite-presentation}).
\end{proof}





\section{Complexes de Koszul}
\label{section-koszul-complex}

\noindent
Nous suggérons de lire d'abord la section consacrée aux complexes de Koszul dans
Compléments d'algèbre, section \ref{more-algebra-section-koszul}.
Nous définissons comme suit le complexe de Koszul dans la catégorie des
$\mathcal{O}_X$-modules.

\begin{definition}
\label{definition-koszul}
Soit $X$ un espace annelé. Soit $\varphi : \mathcal{E} \to \mathcal{O}_X$
un morphisme de $\mathcal{O}_X$-modules. Le
{\it complexe de Koszul} $K_\bullet(\varphi)$ associé à $\varphi$
est le faisceau d'algèbres différentielles graduées commutatives défini comme suit :
\begin{enumerate}
\item l'algèbre graduée sous-jacente est l'algèbre extérieure
$K_\bullet(\varphi) = \wedge(\mathcal{E})$,
\item la différentielle $d : K_\bullet(\varphi) \to K_\bullet(\varphi)$
est l'unique dérivation telle que $d(e) = \varphi(e)$ pour toute
section locale $e$ de $\mathcal{E} = K_1(\varphi)$.
\end{enumerate}
\end{definition}

\noindent
Explicitement, si $e_1 \wedge \ldots \wedge e_n$ est un produit extérieur de
sections locales de $\mathcal{E}$, alors
$$
d(e_1 \wedge \ldots \wedge e_n) =
\sum\nolimits_{i = 1, \ldots, n} (-1)^{i + 1}
\varphi(e_i)e_1 \wedge \ldots \wedge \widehat{e_i} \wedge \ldots \wedge e_n.
$$
On vérifie immédiatement que cette formule définit une dérivation bien définie
sur l'algèbre tensorielle, qui annule $e \wedge e$ et se factorise donc
par l'algèbre extérieure.

\begin{definition}
\label{definition-koszul-complex}
Soit $X$ un espace annelé et soient
$f_1, \ldots, f_n \in \Gamma(X, \mathcal{O}_X)$. Le
{\it complexe de Koszul associé à $f_1, \ldots, f_n$} est le complexe de Koszul
associé au morphisme
$(f_1, \ldots, f_n) : \mathcal{O}_X^{\oplus n} \to \mathcal{O}_X$.
On le note $K_\bullet(\mathcal{O}_X, f_1, \ldots, f_n)$
ou $K_\bullet(\mathcal{O}_X, f_\bullet)$.
\end{definition}

\noindent
Bien entendu, étant donné un morphisme de $\mathcal{O}_X$-modules
$\varphi : \mathcal{E} \to \mathcal{O}_X$,
si $\mathcal{E}$ est localement libre de type fini, alors
$K_\bullet(\varphi)$ est localement sur $X$ isomorphe à un complexe de Koszul
$K_\bullet(\mathcal{O}_X, f_1, \ldots, f_n)$.



\section{Modules inversibles}
\label{section-invertible}

\noindent
Comme dans le cas des modules sur un anneau
(Compléments d'algèbre, section \ref{more-algebra-section-picard}),
nous avons la définition suivante.

\begin{definition}
\label{definition-invertible}
Soit $(X, \mathcal{O}_X)$ un espace annelé. Un
{\it $\mathcal{O}_X$-module inversible} est un faisceau
de $\mathcal{O}_X$-modules $\mathcal{L}$ tel que
le foncteur
$$
\textit{Mod}(\mathcal{O}_X) \longrightarrow \textit{Mod}(\mathcal{O}_X),\quad
\mathcal{F} \longmapsto \mathcal{L} \otimes_{\mathcal{O}_X} \mathcal{F}
$$
est une équivalence de catégories. Nous disons que $\mathcal{L}$ est
{\it trivial} s'il est isomorphe à $\mathcal{O}_X$ en tant que
$\mathcal{O}_X$-module.
\end{definition}

\noindent
Le lemme \ref{lemma-invertible-is-locally-free-rank-1}
ci-dessous explique le lien avec les modules localement libres
de rang $1$.

\begin{lemma}
\label{lemma-invertible}
Soit $(X, \mathcal{O}_X)$ un espace annelé. Soit $\mathcal{L}$
un $\mathcal{O}_X$-module. Les conditions suivantes sont équivalentes :
\begin{enumerate}
\item $\mathcal{L}$ est inversible ;
\item il existe un $\mathcal{O}_X$-module $\mathcal{N}$
tel que
$\mathcal{L} \otimes_{\mathcal{O}_X} \mathcal{N} \cong \mathcal{O}_X$.
\end{enumerate}
Dans ce cas, $\mathcal{L}$ est localement facteur direct d'un
$\mathcal{O}_X$-module libre de type fini, et le module $\mathcal{N}$ de (2) est isomorphe à
$\SheafHom_{\mathcal{O}_X}(\mathcal{L}, \mathcal{O}_X)$.
\end{lemma}

\begin{proof}
Supposons (1). Le foncteur $- \otimes_{\mathcal{O}_X} \mathcal{L}$
est alors essentiellement surjectif ; il existe donc un $\mathcal{O}_X$-module
$\mathcal{N}$ comme en (2). Si (2) est vérifiée, le foncteur
$- \otimes_{\mathcal{O}_X} \mathcal{N}$ est un quasi-inverse
du foncteur $- \otimes_{\mathcal{O}_X} \mathcal{L}$ ;
on en déduit (1).

\medskip\noindent
Supposons (1) et (2) vérifiées. Notons
$\psi : \mathcal{L} \otimes_{\mathcal{O}_X} \mathcal{N} \to \mathcal{O}_X$
l'isomorphisme donné. Soit $x \in X$. Choisissons un voisinage ouvert
$U$ de $x$, un entier $n \geq 1$ et des sections $s_i \in \mathcal{L}(U)$,
$t_i \in \mathcal{N}(U)$ telles que $\psi(\sum s_i \otimes t_i) = 1$.
Considérons les isomorphismes
$$
\mathcal{L}|_U \to
\mathcal{L}|_U \otimes_{\mathcal{O}_U}
\mathcal{L}|_U \otimes_{\mathcal{O}_U} \mathcal{N}|_U \to \mathcal{L}|_U
$$
où la première flèche envoie $s$ sur $\sum s_i \otimes s \otimes t_i$
et la seconde envoie $s \otimes s' \otimes t$ sur $\psi(s' \otimes t)s$.
Nous en concluons que $s \mapsto \sum \psi(s \otimes t_i)s_i$ est
un automorphisme de $\mathcal{L}|_U$. Cet automorphisme se factorise sous la forme
$$
\mathcal{L}|_U \to \mathcal{O}_U^{\oplus n} \to \mathcal{L}|_U
$$
où la première flèche est donnée par
$s \mapsto (\psi(s \otimes t_1), \ldots, \psi(s \otimes t_n))$
et la seconde par $(a_1, \ldots, a_n) \mapsto \sum a_i s_i$.
Nous en concluons que $\mathcal{L}|_U$ est facteur direct
d'un $\mathcal{O}_U$-module libre de type fini.

\medskip\noindent
Supposons (1) et (2) vérifiées. Considérons le morphisme d'évaluation
$$
\mathcal{L} \otimes_{\mathcal{O}_X}
\SheafHom_{\mathcal{O}_X}(\mathcal{L}, \mathcal{O}_X)
\longrightarrow \mathcal{O}_X
$$
Pour achever la démonstration du lemme,
montrons que c'est un isomorphisme en vérifiant qu'il induit
des isomorphismes sur les germes. Soit $x \in X$.
Comme nous savons, d'après le paragraphe précédent,
que $\mathcal{L}$ est un $\mathcal{O}_X$-module
de présentation finie,
le lemme \ref{lemma-stalk-internal-hom}
montre qu'il suffit d'établir que
$$
\mathcal{L}_x \otimes_{\mathcal{O}_{X, x}}
\Hom_{\mathcal{O}_{X, x}}(\mathcal{L}_x, \mathcal{O}_{X, x})
\longrightarrow \mathcal{O}_{X, x}
$$
est un isomorphisme. Puisque
$\mathcal{L}_x \otimes_{\mathcal{O}_{X, x}} \mathcal{N}_x =
(\mathcal{L} \otimes_{\mathcal{O}_X} \mathcal{N})_x =
\mathcal{O}_{X, x}$ (lemme \ref{lemma-stalk-tensor-product}),
le résultat voulu découle de
Compléments d'algèbre, lemme \ref{more-algebra-lemma-invertible}.
\end{proof}

\begin{lemma}
\label{lemma-pullback-invertible}
Soit $f : (X, \mathcal{O}_X) \to (Y, \mathcal{O}_Y)$ un
morphisme d'espaces annelés. L'image inverse $f^*\mathcal{L}$ d'un
$\mathcal{O}_Y$-module inversible est inversible.
\end{lemma}

\begin{proof}
D'après le lemme \ref{lemma-invertible},
il existe un $\mathcal{O}_Y$-module $\mathcal{N}$ tel que
$\mathcal{L} \otimes_{\mathcal{O}_Y} \mathcal{N} \cong \mathcal{O}_Y$.
En prenant l'image inverse, nous obtenons
$f^*\mathcal{L} \otimes_{\mathcal{O}_X} f^*\mathcal{N} \cong \mathcal{O}_X$
d'après le lemme \ref{lemma-tensor-product-pullback}.
Ainsi $f^*\mathcal{L}$ est inversible d'après le lemme \ref{lemma-invertible}.
\end{proof}

\begin{lemma}
\label{lemma-invertible-is-locally-free-rank-1}
Soit $(X, \mathcal{O}_X)$ un espace annelé. Tout
$\mathcal{O}_X$-module localement libre de rang $1$ est inversible.
Si tous les anneaux $\mathcal{O}_{X, x}$ sont locaux, alors
la réciproque est également vraie (ce qui n'est pas le cas en général).
\end{lemma}

\begin{proof}
L'assertion entre parenthèses s'obtient en considérant un espace réduit à un point
$X$, muni du faisceau d'anneaux $\mathcal{O}_X$ défini par un anneau $R$.
Les $\mathcal{O}_X$-modules inversibles correspondent alors aux
$R$-modules inversibles ; dès que $\Pic(R)$ n'est pas le groupe trivial,
nous obtenons donc un exemple.

\medskip\noindent
Supposons $\mathcal{L}$ localement libre de rang $1$ et considérons le
morphisme d'évaluation
$$
\mathcal{L} \otimes_{\mathcal{O}_X}
\SheafHom_{\mathcal{O}_X}(\mathcal{L}, \mathcal{O}_X)
\longrightarrow \mathcal{O}_X
$$
Sur un recouvrement ouvert qui trivialise $\mathcal{L}$, nous voyons
que ce morphisme est un isomorphisme. Ainsi $\mathcal{L}$ est inversible
d'après le lemme \ref{lemma-invertible}.

\medskip\noindent
Supposons que tous les anneaux $\mathcal{O}_{X, x}$ soient locaux et que $\mathcal{L}$
soit inversible. Dans la démonstration du lemme \ref{lemma-invertible},
nous avons vu que $\mathcal{L}_x$ est un
$\mathcal{O}_{X, x}$-module inversible pour tout $x \in X$. Comme
$\mathcal{O}_{X, x}$ est local, nous voyons que
$\mathcal{L}_x \cong \mathcal{O}_{X, x}$
(Compléments d'algèbre, section \ref{more-algebra-section-picard}).
Puisque $\mathcal{L}$ est de présentation finie d'après le
lemme \ref{lemma-invertible}, nous concluons que $\mathcal{L}$
est localement libre de rang $1$ d'après le
lemme \ref{lemma-finite-presentation-stalk-free}.
\end{proof}

\begin{lemma}
\label{lemma-constructions-invertible}
Soit $(X, \mathcal{O}_X)$ un espace annelé.
\begin{enumerate}
\item Si $\mathcal{L}$ et $\mathcal{N}$ sont des
$\mathcal{O}_X$-modules inversibles, alors il en va de même de
$\mathcal{L} \otimes_{\mathcal{O}_X} \mathcal{N}$.
\item Si $\mathcal{L}$ est un $\mathcal{O}_X$-module inversible, alors
$\SheafHom_{\mathcal{O}_X}(\mathcal{L}, \mathcal{O}_X)$ l'est aussi, et le morphisme d'évaluation
$\mathcal{L} \otimes_{\mathcal{O}_X}
\SheafHom_{\mathcal{O}_X}(\mathcal{L}, \mathcal{O}_X) \to \mathcal{O}_X$
est un isomorphisme.
\end{enumerate}
\end{lemma}

\begin{proof}
L'assertion (1) est immédiate par définition, et (2) résulte du
lemme \ref{lemma-invertible} et de sa démonstration.
\end{proof}

\begin{definition}
\label{definition-powers}
Soit $(X, \mathcal{O}_X)$ un espace annelé. Étant donné un faisceau inversible
$\mathcal{L}$ sur $X$ et $n \in \mathbf{Z}$, nous définissons la
$n$-ième {\it puissance tensorielle} $\mathcal{L}^{\otimes n}$ de $\mathcal{L}$
comme l'image de $\mathcal{O}_X$ par l'application de l'équivalence
$\mathcal{F} \mapsto\mathcal{F} \otimes_{\mathcal{O}_X} \mathcal{L}$
exactement $n$ fois.
\end{definition}

\noindent
Cette définition a également un sens pour les entiers négatifs $n$, puisque nous avons
défini un $\mathcal{O}_X$-module inversible comme un module pour lequel le produit tensoriel
est une équivalence. Plus explicitement, nous avons
$$
\mathcal{L}^{\otimes n} =
\left\{
\begin{matrix}
\mathcal{O}_X & \text{si} & n = 0 \\
\SheafHom_{\mathcal{O}_X}(\mathcal{L}, \mathcal{O}_X) & \text{si} & n = -1\\
\mathcal{L} \otimes_{\mathcal{O}_X} \ldots \otimes_{\mathcal{O}_X} \mathcal{L}
& \text{si} & n > 0 \\
\mathcal{L}^{\otimes -1} \otimes_{\mathcal{O}_X} \ldots
\otimes_{\mathcal{O}_X} \mathcal{L}^{\otimes -1}
& \text{si} & n < -1
\end{matrix}
\right.
$$
voir le lemme \ref{lemma-constructions-invertible}.
Avec cette définition, nous avons des isomorphismes canoniques
$\mathcal{L}^{\otimes n} \otimes_{\mathcal{O}_X}
\mathcal{L}^{\otimes m} \to
\mathcal{L}^{\otimes n + m}$, et ces isomorphismes
satisfont à des contraintes de commutativité et d'associativité
(formulation omise).

\medskip\noindent
Soit $(X, \mathcal{O}_X)$ un espace annelé. Nous pouvons définir une structure
d'anneau $\mathbf{Z}$-gradué sur $\bigoplus \Gamma(X, \mathcal{L}^{\otimes n})$ en envoyant
$s \in \Gamma(X, \mathcal{L}^{\otimes n})$ et
$t \in \Gamma(X, \mathcal{L}^{\otimes m})$ sur la
section correspondant à $s \otimes t$ dans
$\Gamma(X, \mathcal{L}^{\otimes n + m})$.
Nous omettons de vérifier que cela définit un anneau commutatif
et associatif, d'unité $1$. Toutefois, selon nos conventions dans
Algèbre, section \ref{algebra-section-graded},
un anneau gradué ne possède aucun élément non nul de degré négatif.
Cela conduit à la définition suivante.

\begin{definition}
\label{definition-gamma-star}
Soit $(X, \mathcal{O}_X)$ un espace annelé.
Étant donné un faisceau inversible $\mathcal{L}$ sur $X$, nous définissons
l'{\it anneau gradué associé} par
$$
\Gamma_*(X, \mathcal{L})
=
\bigoplus\nolimits_{n \geq 0} \Gamma(X, \mathcal{L}^{\otimes n})
$$
Étant donné un faisceau de $\mathcal{O}_X$-modules $\mathcal{F}$, nous posons
$$
\Gamma_*(X, \mathcal{L}, \mathcal{F})
=
\bigoplus\nolimits_{n \in \mathbf{Z}} \Gamma(X,
\mathcal{F} \otimes_{\mathcal{O}_X} \mathcal{L}^{\otimes n})
$$
que nous considérons comme un $\Gamma_*(X, \mathcal{L})$-module gradué.
\end{definition}

\noindent
Nous écrivons souvent simplement $\Gamma_*(\mathcal{L})$ et $\Gamma_*(\mathcal{F})$
(bien que cette notation soit ambiguë si $\mathcal{F}$ est inversible).
La multiplication de $\Gamma_*(\mathcal{L})$ sur
$\Gamma_*(\mathcal{F})$ est définie au moyen des isomorphismes
ci-dessus. Si $\gamma : \mathcal{F} \to \mathcal{G}$ est un morphisme de
$\mathcal{O}_X$-modules, nous obtenons un morphisme de $\Gamma_*(\mathcal{L})$-modules
$\gamma : \Gamma_*(\mathcal{F}) \to \Gamma_*(\mathcal{G})$.
Si $\alpha : \mathcal{L} \to \mathcal{N}$ est un morphisme de
$\mathcal{O}_X$-modules entre des $\mathcal{O}_X$-modules inversibles, nous obtenons
un morphisme d'anneaux gradués $\Gamma_*(\mathcal{L}) \to \Gamma_*(\mathcal{N})$.
Si $f : (Y, \mathcal{O}_Y) \to (X, \mathcal{O}_X)$
est un morphisme d'espaces annelés et si $\mathcal{L}$ est inversible
sur $X$, nous obtenons un faisceau inversible $f^*\mathcal{L}$ sur $Y$
(lemme \ref{lemma-pullback-invertible})
et un morphisme induit d'anneaux gradués
$$
f^* :
\Gamma_*(X, \mathcal{L})
\longrightarrow
\Gamma_*(Y, f^*\mathcal{L})
$$
Il existe en outre certaines compatibilités entre les constructions
ci-dessus, dont nous omettons les énoncés.

\begin{lemma}
\label{lemma-pic-set}
Soit $(X, \mathcal{O}_X)$ un espace annelé.
Il existe un ensemble de modules inversibles $\{\mathcal{L}_i\}_{i \in I}$
tel que tout module inversible sur $X$ soit isomorphe à exactement
l'un des $\mathcal{L}_i$.
\end{lemma}

\begin{proof}
Rappelons que tout $\mathcal{O}_X$-module inversible est localement
facteur direct d'un $\mathcal{O}_X$-module libre de type fini ; voir le
lemme \ref{lemma-invertible}.
Pour tout recouvrement ouvert $\mathcal{U} : X = \bigcup_{j \in J} U_j$
et toute application $r : J \to \mathbf{N}$, considérons les faisceaux de
$\mathcal{O}_X$-modules $\mathcal{F}$ tels que
$\mathcal{F}_j = \mathcal{F}|_{U_j}$ soit facteur direct de
$\mathcal{O}_{U_j}^{\oplus r(j)}$.
La collection des classes d'isomorphisme de $\mathcal{F}_j$ est un ensemble, car
$\Hom_{\mathcal{O}_U}(\mathcal{O}_U^{\oplus r}, \mathcal{O}_U^{\oplus r})$
est un ensemble. Le faisceau $\mathcal{F}$ s'obtient par recollement des $\mathcal{F}_j$ ;
voir Faisceaux, section
\ref{sheaves-section-glueing-sheaves}. Remarquons que la collection de
toutes les données de recollement forme un ensemble. La collection de tous les recouvrements
$\mathcal{U} : X = \bigcup_{j \in J} U_j$ pour lesquels $J \to \mathcal{P}(X)$,
$j \mapsto U_j$ est injective forme elle aussi un ensemble. Pour chaque recouvrement,
il existe un ensemble d'applications $r : J \to \mathbf{N}$. La collection
de tous les $\mathcal{F}$ forme donc un ensemble.
\end{proof}

\noindent
Ce lemme affirme, en substance, que la collection des
classes d'isomorphisme de faisceaux inversibles forme un ensemble.
Le lemme \ref{lemma-constructions-invertible} affirme que
le produit tensoriel définit sur cet ensemble une structure de groupe
abélien.

\begin{definition}
\label{definition-pic}
Soit $(X, \mathcal{O}_X)$ un espace annelé.
Le {\it groupe de Picard} $\Pic(X)$ de $X$ est le
groupe abélien dont les éléments sont les classes d'isomorphisme de
$\mathcal{O}_X$-modules inversibles, l'addition
correspondant au produit tensoriel.
\end{definition}

\begin{lemma}
\label{lemma-s-open}
\begin{slogan}
Une trivialisation (locale) d'un fibré en droites
équivaut à une section (locale) ne s'annulant pas.
\end{slogan}
Soit $X$ un espace annelé. Supposons que chaque anneau local $\mathcal{O}_{X, x}$
soit un anneau local d'idéal maximal $\mathfrak m_x$.
Soit $\mathcal{L}$ un $\mathcal{O}_X$-module inversible.
Pour toute section $s \in \Gamma(X, \mathcal{L})$, l'ensemble
$$
X_s = \{x \in X \mid \text{image de }s \not\in \mathfrak m_x\mathcal{L}_x\}
$$
est ouvert dans $X$. Le morphisme $s : \mathcal{O}_{X_s} \to \mathcal{L}|_{X_s}$
est un isomorphisme, et il existe une section $s'$
de $\mathcal{L}^{\otimes -1}$ sur $X_s$ telle que $s' (s|_{X_s}) = 1$.
\end{lemma}

\begin{proof}
Supposons $x \in X_s$.
Nous avons un isomorphisme
$$
\mathcal{L}_x \otimes_{\mathcal{O}_{X, x}} (\mathcal{L}^{\otimes -1})_x
\longrightarrow
\mathcal{O}_{X, x}
$$
d'après le lemme \ref{lemma-constructions-invertible}.
Les modules $\mathcal{L}_x$ et $(\mathcal{L}^{\otimes -1})_x$
sont tous deux des $\mathcal{O}_{X, x}$-modules libres de rang $1$. Le
lemme de Nakayama (Algèbre, lemme \ref{algebra-lemma-NAK}) montre que
$s_x$ est une base de $\mathcal{L}_x$. Il existe donc
un élément de base $t_x \in (\mathcal{L}^{\otimes -1})_x$
tel que $s_x \otimes t_x$ s'envoie sur $1$.
Choisissons un voisinage ouvert $U$ de
$x$ tel que $t_x$ provienne d'une section $t$
de $\mathcal{L}^{\otimes -1}$ sur $U$ et que
$s \otimes t$ s'envoie sur $1 \in \mathcal{O}_X(U)$.
Il est clair que, pour tout $x' \in U$, la section $s$ engendre
le module $\mathcal{L}_{x'}$. Ainsi $U \subset X_s$.
Ceci prouve que $X_s$ est ouvert. De plus, la section
$t$ construite ci-dessus sur $U$ est unique ; ces sections
se recollent donc en la section $s'$ du lemme.
\end{proof}

\begin{remark}
\label{remark-pullback-s-open}
Étant donné un morphisme d'espaces localement annelés $f : Y \to X$
(voir Schémas, définition \ref{schemes-definition-locally-ringed-space}),
l'image inverse $f^{-1}(X_s)$ est égale à $Y_{f^*s}$, où
$f^*s \in \Gamma(Y, f^*\mathcal{L})$ est l'image inverse de $s$.
\end{remark}


\section{Rang et déterminant}
\label{section-rank-and-det}

\noindent
Soit $(X, \mathcal{O}_X)$ un espace annelé. Considérons la catégorie
$\textit{Vect}(X)$ des $\mathcal{O}_X$-modules localement libres de type fini.
C'est une catégorie exacte
(voir Injectifs, remarque \ref{injectives-remark-embed-exact-category})
dont les épimorphismes admissibles sont les
surjections et dont les monomorphismes admissibles sont les noyaux de
surjections. De plus, les classes d'isomorphisme des objets de
$\textit{Vect}(X)$ forment un ensemble (démonstration omise). Nous pouvons donc former
le $K$-groupe de Grothendieck d'indice zéro $K_0(\textit{Vect}(X))$.
Explicitement, dans ce cas, $K_0(\textit{Vect}(X))$
est le groupe abélien engendré par $[\mathcal{E}]$, où $\mathcal{E}$
est un $\mathcal{O}_X$-module localement libre de type fini, soumis aux relations
$$
[\mathcal{E}] = [\mathcal{E}'] + [\mathcal{E}'']
$$
pour toute suite exacte courte
$0 \to \mathcal{E}' \to \mathcal{E} \to \mathcal{E}'' \to 0$
de $\mathcal{O}_X$-modules localement libres de type fini.

\medskip\noindent
{\bf Rangs.} Supposons que tous les anneaux locaux $\mathcal{O}_{X, x}$ soient non nuls.
Étant donné un $\mathcal{O}_X$-module localement libre de type fini $\mathcal{E}$,
le {\it rang} est la fonction localement constante
$$
\text{rang}_\mathcal{E} : X \longrightarrow \mathbf{Z}_{\geq 0},\quad
x \longmapsto \text{rang}_{\mathcal{O}_{X, x}} \mathcal{E}_x
$$
Voir le lemme \ref{lemma-rank}. Par définition des modules localement
libres, la fonction $\text{rang}_\mathcal{E}$ est localement constante. Si
$0 \to \mathcal{E}' \to \mathcal{E} \to \mathcal{E}'' \to 0$
est une suite exacte courte de $\mathcal{O}_X$-modules localement libres de type fini,
alors $\text{rang}_\mathcal{E} = \text{rang}_{\mathcal{E}'} +
\text{rang}_{\mathcal{E}''}$.
Le rang définit donc un morphisme
$$
K_0(\textit{Vect}(X)) \longrightarrow \text{Map}_{cont}(X, \mathbf{Z}),\quad
[\mathcal{E}] \longmapsto \text{rang}_\mathcal{E}
$$

\medskip\noindent
{\bf Déterminants.} Étant donné un
$\mathcal{O}_X$-module localement libre de type fini $\mathcal{E}$, nous obtenons une
décomposition en union disjointe
$$
X = X_0 \amalg X_1 \amalg X_2 \amalg \ldots
$$
où les $X_i$ sont ouverts et fermés, de telle sorte que $\mathcal{E}$ soit localement
libre de type fini et de rang $i$ sur $X_i$ (cela revient exactement à dire que
$\text{rang}_\mathcal{E}$ est localement constante). Dans ce cas, nous définissons
$\det(\mathcal{E})$ comme le faisceau inversible sur $X$ qui est égal à
$\wedge^i(\mathcal{E}|_{X_i})$ sur $X_i$ pour tout $i \geq 0$.
La décomposition ci-dessus étant disjointe, il n'y a aucune condition de recollement
à vérifier. D'après le lemme \ref{lemma-det-ses} ci-dessous,
ceci définit un morphisme
$$
\det : K_0(\textit{Vect}(X)) \longrightarrow \Pic(X),\quad
[\mathcal{E}] \longmapsto \det(\mathcal{E})
$$
de groupes abéliens. Les éléments de $\Pic(X)$ ainsi obtenus
sont localement libres de rang $1$ (voir, après le lemme, une généralisation).

\begin{lemma}
\label{lemma-det-ses}
Soit $X$ un espace annelé. Soit
$0 \to \mathcal{E}' \to \mathcal{E} \to \mathcal{E}'' \to 0$
une suite exacte courte de $\mathcal{O}_X$-modules localement libres de type fini.
Il existe alors un isomorphisme canonique
$$
\det(\mathcal{E}') \otimes_{\mathcal{O}_X}\det(\mathcal{E}'')
\longrightarrow
\det(\mathcal{E})
$$
de $\mathcal{O}_X$-modules.
\end{lemma}

\begin{proof}
Nous pouvons décomposer $X$ en sous-ensembles ouverts et fermés disjoints sur lesquels
$\mathcal{E}'$ et $\mathcal{E}''$ sont tous deux de rang constant.
Nous nous ramenons donc au cas où $\mathcal{E}'$ et $\mathcal{E}''$
sont de rang constant, disons $r'$ et $r''$. Dans cette situation, nous
définissons
$$
\wedge^{r'}(\mathcal{E}') \otimes_{\mathcal{O}_X} \wedge^{r''}(\mathcal{E}'')
\longrightarrow
\wedge^{r' + r''}(\mathcal{E})
$$
comme suit. Étant données des sections locales $s'_1, \ldots, s'_{r'}$ de $\mathcal{E}'$
et des sections locales $s''_1, \ldots, s''_{r''}$ de $\mathcal{E}''$,
nous envoyons
$$
s'_1 \wedge \ldots \wedge s'_{r'} \otimes
s''_1 \wedge \ldots \wedge s''_{r''}
\quad\text{sur}\quad
s'_1 \wedge \ldots \wedge s'_{r'} \wedge
\tilde s''_1 \wedge \ldots \wedge \tilde s''_{r''}
$$
où $\tilde s''_i$ est un relèvement local de la section
$s''_i$ en une section de $\mathcal{E}$. Nous omettons les détails.
\end{proof}

\noindent
Soit $(X, \mathcal{O}_X)$ un espace annelé. Au lieu de considérer les
$\mathcal{O}_X$-modules localement libres de type fini, nous pouvons considérer les
$\mathcal{O}_X$-modules $\mathcal{F}$ qui sont localement sur $X$
facteurs directs d'un $\mathcal{O}_X$-module libre de type fini. Cela revient à
demander que $\mathcal{F}$ soit un $\mathcal{O}_X$-module plat et de
présentation finie ; voir le lemme \ref{lemma-flat-locally-finite-presentation}.
Si tous les anneaux locaux $\mathcal{O}_{X, x}$ sont locaux, alors un tel module
$\mathcal{F}$ est localement libre de type fini ; voir le
lemme \ref{lemma-direct-summand-of-locally-free-is-locally-free}.
Ce n'est toutefois pas le cas en général ; par
exemple, $X$ pourrait être un point, $\Gamma(X, \mathcal{O}_X)$ pourrait
être le produit $A \times B$ de deux anneaux non nuls et $\mathcal{F}$
pourrait correspondre à $A \times 0$.
Pour un tel module, la fonction rang n'est donc pas définie.
Néanmoins, il est encore possible de définir $\det(\mathcal{F})$,
qui sera un $\mathcal{O}_X$-module inversible au sens de la
définition \ref{definition-invertible} (sans être nécessairement localement
libre de rang $1$). Notre construction coïncidera avec celle donnée ci-dessus
lorsque $\mathcal{F}$ est localement libre de type fini.
Nous recommandons au lecteur de sauter le reste de cette section.

\begin{lemma}
\label{lemma-determinant-as-socle}
Soit $(X, \mathcal{O}_X)$ un espace annelé. Soit $\mathcal{F}$
un $\mathcal{O}_X$-module plat et de présentation finie.
Notons
$$
\det(\mathcal{F}) \subset
\wedge^*_{\mathcal{O}_X}(\mathcal{F})
$$
l'annulateur de $\mathcal{F} \subset \wedge^*_{\mathcal{O}_X}(\mathcal{F})$.
Alors $\det(\mathcal{F})$ est un $\mathcal{O}_X$-module inversible.
\end{lemma}

\begin{proof}
Pour le démontrer, nous pouvons travailler localement sur $X$. Nous pouvons donc supposer que $\mathcal{F}$
est facteur direct d'un module libre de type fini ; voir le
lemme \ref{lemma-flat-locally-finite-presentation}. Écrivons
$\mathcal{F} \oplus \mathcal{G} = \mathcal{O}_X^{\oplus n}$.
Posons $R = \mathcal{O}_X(X)$. Nous avons alors
$\mathcal{F}(X) \oplus \mathcal{G}(X) = R^{\oplus n}$
et, de même,
$\mathcal{F}(U) \oplus \mathcal{G}(U) = \mathcal{O}_X(U)^{\oplus n}$
pour tout ouvert $U \subset X$.
Nous en concluons que $\mathcal{F} = \mathcal{F}_M$ comme dans le
lemme \ref{lemma-construct-quasi-coherent-sheaves},
où $M = \mathcal{F}(X)$ est un $R$-module projectif de type fini.
Autrement dit, nous avons $\mathcal{F}(U) = M \otimes_R \mathcal{O}_X(U)$.
Ceci implique que
$\det(M) \otimes_R \mathcal{O}_X(U) = \det(\mathcal{F}(U))$
pour tout ouvert $U \subset X$, où $\det$ est celui de
Compléments d'algèbre, section \ref{more-algebra-section-determinants}. D'après
Compléments d'algèbre, remarque \ref{more-algebra-remark-determinant-as-socle},
nous voyons que
$$
\det(M) \otimes_R \mathcal{O}_X(U) =
\det(\mathcal{F}(U)) \subset
\wedge^*_{\mathcal{O}_X(U)}(\mathcal{F}(U))
$$
est l'annulateur de $\mathcal{F}(U)$. Nous en concluons que
$\det(\mathcal{F})$, tel qu'il est défini dans l'énoncé du lemme,
est égal à $\mathcal{F}_{\det(M)}$. Nous omettons certains détails : il faut
prendre garde au fait que les annulateurs ne peuvent pas être définis en faisceautisant
la prise des annulateurs sur les sections au-dessus des ouverts. Ainsi $\det(\mathcal{F})$
est l'image inverse d'un module inversible, ce qui conclut la démonstration.
\end{proof}
\section{Localisation des faisceaux d'anneaux}
\label{section-localizing-sheaves-rings}

\noindent
Soit $X$ un espace topologique et soit $\mathcal{O}_X$ un
préfaisceau d'anneaux. Soit $\mathcal{S} \subset \mathcal{O}_X$
un préfaisceau d'ensembles contenu dans $\mathcal{O}_X$.
Supposons que, pour tout ouvert $U \subset X$, l'ensemble
$\mathcal{S}(U) \subset \mathcal{O}_X(U)$ soit une partie multiplicative;
voir Algèbre, définition \ref{algebra-definition-multiplicative-subset}.
On peut alors considérer le préfaisceau d'anneaux
$$
\mathcal{S}^{-1}\mathcal{O}_X :
U \longmapsto \mathcal{S}(U)^{-1}\mathcal{O}_X(U).
$$
L'application de restriction envoie la section $f/s$, où
$f \in \mathcal{O}_X(U)$ et $s \in \mathcal{S}(U)$, sur
$(f|_V)/(s|_V)$ lorsque $V \subset U$ sont des ouverts de $X$.

\begin{lemma}
\label{lemma-simple-invert}
Soit $X$ un espace topologique et soit $\mathcal{O}_X$ un
préfaisceau d'anneaux. Soit $\mathcal{S} \subset \mathcal{O}_X$
un préfaisceau d'ensembles contenu dans $\mathcal{O}_X$.
Supposons que, pour tout ouvert $U \subset X$, l'ensemble
$\mathcal{S}(U) \subset \mathcal{O}_X(U)$ soit une partie multiplicative.
\begin{enumerate}
\item Il existe un morphisme de préfaisceaux d'anneaux
$\mathcal{O}_X \to \mathcal{S}^{-1}\mathcal{O}_X$
tel que toute section locale de $\mathcal{S}$ ait pour image une section
inversible de $\mathcal{S}^{-1}\mathcal{O}_X$.
\item Pour tout homomorphisme de préfaisceaux d'anneaux
$\mathcal{O}_X \to \mathcal{A}$ tel que toute section locale
de $\mathcal{S}$ ait pour image une section inversible de $\mathcal{A}$,
il existe une unique factorisation
$\mathcal{S}^{-1}\mathcal{O}_X \to \mathcal{A}$.
\item Pour tout $x \in X$, on a
$$
(\mathcal{S}^{-1}\mathcal{O}_X)_x = \mathcal{S}_x^{-1} \mathcal{O}_{X, x}.
$$
\item Le faisceau associé $(\mathcal{S}^{-1}\mathcal{O}_X)^\#$ est un
faisceau d'anneaux muni d'un morphisme de faisceaux d'anneaux
$(\mathcal{O}_X)^\# \to (\mathcal{S}^{-1}\mathcal{O}_X)^\#$
qui est universel parmi les morphismes de $(\mathcal{O}_X)^\#$ vers des
faisceaux d'anneaux envoyant toute section locale de $\mathcal{S}$ sur une
section inversible.
\item Pour tout $x \in X$, on a
$$
(\mathcal{S}^{-1}\mathcal{O}_X)^\#_x = \mathcal{S}_x^{-1} \mathcal{O}_{X, x}.
$$
\end{enumerate}
\end{lemma}

\begin{proof}
Démonstration omise.
\end{proof}

\noindent
Soit $X$ un espace topologique et soit $\mathcal{O}_X$ un
préfaisceau d'anneaux. Soit $\mathcal{S} \subset \mathcal{O}_X$
un préfaisceau d'ensembles contenu dans $\mathcal{O}_X$.
Supposons que, pour tout ouvert $U \subset X$, l'ensemble
$\mathcal{S}(U) \subset \mathcal{O}_X(U)$ soit une partie multiplicative.
Soit $\mathcal{F}$ un préfaisceau de $\mathcal{O}_X$-modules.
On peut alors considérer le préfaisceau de
$\mathcal{S}^{-1}\mathcal{O}_X$-modules
$$
\mathcal{S}^{-1}\mathcal{F} :
U \longmapsto \mathcal{S}(U)^{-1}\mathcal{F}(U).
$$
L'application de restriction envoie la section $t/s$, où
$t \in \mathcal{F}(U)$ et $s \in \mathcal{S}(U)$, sur
$(t|_V)/(s|_V)$ lorsque $V \subset U$ sont des ouverts de $X$.

\begin{lemma}
\label{lemma-simple-invert-module}
Soit $X$ un espace topologique.
Soit $\mathcal{O}_X$ un préfaisceau d'anneaux.
Soit $\mathcal{S} \subset \mathcal{O}_X$ un préfaisceau d'ensembles contenu
dans $\mathcal{O}_X$. Supposons que, pour tout ouvert $U \subset X$,
l'ensemble $\mathcal{S}(U) \subset \mathcal{O}_X(U)$ soit une partie
multiplicative. Pour tout préfaisceau de $\mathcal{O}_X$-modules
$\mathcal{F}$, on a
$$
\mathcal{S}^{-1}\mathcal{F}
=
\mathcal{S}^{-1}\mathcal{O}_X \otimes_{p, \mathcal{O}_X} \mathcal{F}
$$
(voir Faisceaux, section \ref{sheaves-section-presheaves-modules} pour la
notation) et, si $\mathcal{F}$ et $\mathcal{O}_X$ sont des faisceaux, alors
$$
(\mathcal{S}^{-1}\mathcal{F})^\#
=
(\mathcal{S}^{-1}\mathcal{O}_X)^\# \otimes_{\mathcal{O}_X} \mathcal{F}
$$
(voir Faisceaux, section
\ref{sheaves-section-sheafification-presheaves-modules} pour la notation).
\end{lemma}

\begin{proof}
Démonstration omise.
\end{proof}








\section{Modules de différentielles}
\label{section-differentials}

\noindent
Dans cette section, nous expliquons brièvement comment définir le module des
différentielles relatives d'un morphisme d'espaces annelés. Nous conseillons
au lecteur de consulter la section correspondante du chapitre d'algèbre
commutative (Algèbre, section \ref{algebra-section-differentials}).

\begin{definition}
\label{definition-derivation}
Soit $X$ un espace topologique. Soit
$\varphi : \mathcal{O}_1 \to \mathcal{O}_2$ un homomorphisme de faisceaux
d'anneaux. Soit $\mathcal{F}$ un $\mathcal{O}_2$-module. Une
{\it $\mathcal{O}_1$-dérivation}, ou plus précisément une
{\it $\varphi$-dérivation}, à valeurs dans $\mathcal{F}$ est une application
$D : \mathcal{O}_2 \to \mathcal{F}$ additive, nulle sur l'image de
$\mathcal{O}_1 \to \mathcal{O}_2$, et satisfaisant la
{\it règle de Leibniz}
$$
D(ab) = aD(b) + D(a)b
$$
pour toutes sections locales $a, b$ de $\mathcal{O}_2$ (sur tout ouvert où
elles sont toutes deux définies). On note
$\text{Der}_{\mathcal{O}_1}(\mathcal{O}_2, \mathcal{F})$ l'ensemble des
$\varphi$-dérivations à valeurs dans $\mathcal{F}$.
\end{definition}

\noindent
C'est l'analogue faisceautique d'Algèbre, définition
\ref{algebra-definition-derivation}. Étant donnée une dérivation
$D : \mathcal{O}_2 \to \mathcal{F}$ comme dans la définition,
l'application induite sur les sections globales
$$
D : \Gamma(X, \mathcal{O}_2) \longrightarrow \Gamma(X, \mathcal{F})
$$
est une $\Gamma(X, \mathcal{O}_1)$-dérivation au sens de la définition
algébrique. Remarquons que, si
$\alpha : \mathcal{F} \to \mathcal{G}$ est un morphisme de
$\mathcal{O}_2$-modules, il existe une application induite
$$
\text{Der}_{\mathcal{O}_1}(\mathcal{O}_2, \mathcal{F})
\longrightarrow
\text{Der}_{\mathcal{O}_1}(\mathcal{O}_2, \mathcal{G})
$$
donnée par $D \mapsto \alpha \circ D$. Autrement dit, on obtient un
foncteur.

\begin{lemma}
\label{lemma-universal-module}
Soit $X$ un espace topologique. Soit
$\varphi : \mathcal{O}_1 \to \mathcal{O}_2$ un homomorphisme de faisceaux
d'anneaux. Le foncteur
$$
\textit{Mod}(\mathcal{O}_2) \longrightarrow \textit{Ab}, \quad
\mathcal{F} \longmapsto \text{Der}_{\mathcal{O}_1}(\mathcal{O}_2, \mathcal{F})
$$
est représentable.
\end{lemma}

\begin{proof}
La démonstration est exactement la même que celle de l'énoncé analogue en
algèbre. Au cours de cette démonstration, pour tout faisceau d'ensembles
$\mathcal{F}$ sur $X$, notons $\mathcal{O}_2[\mathcal{F}]$ le faisceau
associé au préfaisceau $U \mapsto \mathcal{O}_2(U)[\mathcal{F}(U)]$, où
cette dernière expression désigne le $\mathcal{O}_2(U)$-module libre de base
l'ensemble $\mathcal{F}(U)$. Pour $s \in \mathcal{F}(U)$, on note $[s]$ la
section correspondante de $\mathcal{O}_2[\mathcal{F}]$ sur $U$. Si
$\mathcal{F}$ est un faisceau de $\mathcal{O}_2$-modules, il existe un
morphisme canonique
$$
c : \mathcal{O}_2[\mathcal{F}] \longrightarrow \mathcal{F}
$$
qui, au niveau des préfaisceaux, est donné par la règle
$\sum f_s[s] \mapsto \sum f_s s$. Nous emploierons l'abréviation
$[s] \mapsto s$ pour décrire ce morphisme, et des abréviations analogues
pour les morphismes ci-dessous. Considérons le morphisme de
$\mathcal{O}_2$-modules
\begin{equation}
\label{equation-define-module-differentials}
\begin{matrix}
\mathcal{O}_2[\mathcal{O}_2 \times \mathcal{O}_2] \oplus
\mathcal{O}_2[\mathcal{O}_2 \times \mathcal{O}_2] \oplus
\mathcal{O}_2[\mathcal{O}_1] &
\longrightarrow &
\mathcal{O}_2[\mathcal{O}_2] \\
[(a, b)] \oplus [(f, g)] \oplus [h] & \longmapsto & [a + b] - [a] - [b] + \\
& & [fg] - g[f] - f[g] + \\
& & [\varphi(h)]
\end{matrix}
\end{equation}
avec la notation abrégée ci-dessus. Posons
$\Omega_{\mathcal{O}_2/\mathcal{O}_1}$ égal au conoyau de ce morphisme.
Il est alors clair qu'il existe un morphisme de faisceaux d'ensembles
$$
\text{d} : \mathcal{O}_2 \longrightarrow \Omega_{\mathcal{O}_2/\mathcal{O}_1}
$$
qui envoie une section locale $f$ sur l'image de $[f]$ dans
$\Omega_{\mathcal{O}_2/\mathcal{O}_1}$. Par construction, $\text{d}$ est une
$\mathcal{O}_1$-dérivation. Soit ensuite $\mathcal{F}$ un faisceau de
$\mathcal{O}_2$-modules et soit
$D : \mathcal{O}_2 \to \mathcal{F}$ une $\mathcal{O}_1$-dérivation.
On peut considérer le morphisme $\mathcal{O}_2$-linéaire
$\mathcal{O}_2[\mathcal{O}_2] \to \mathcal{F}$ qui envoie $[g]$ sur $D(g)$.
La définition d'une dérivation montre que ce morphisme annule les sections
appartenant à l'image du morphisme
(\ref{equation-define-module-differentials}); il définit donc un morphisme
$$
\alpha_D : \Omega_{\mathcal{O}_2/\mathcal{O}_1} \longrightarrow \mathcal{F}
$$
Comme $D = \alpha_D \circ \text{d}$, le lemme est démontré.
\end{proof}

\begin{definition}
\label{definition-module-differentials}
Soit $X$ un espace topologique. Soit
$\varphi : \mathcal{O}_1 \to \mathcal{O}_2$ un homomorphisme de faisceaux
d'anneaux sur $X$. Le {\it module des différentielles} de $\varphi$ est
l'objet représentant le foncteur
$\mathcal{F} \mapsto \text{Der}_{\mathcal{O}_1}(\mathcal{O}_2, \mathcal{F})$
dont l'existence résulte du lemme \ref{lemma-universal-module}.
On le note $\Omega_{\mathcal{O}_2/\mathcal{O}_1}$, et la
{\it $\varphi$-dérivation universelle} se note
$\text{d} : \mathcal{O}_2 \to \Omega_{\mathcal{O}_2/\mathcal{O}_1}$.
\end{definition}

\noindent
Remarquons que $\Omega_{\mathcal{O}_2/\mathcal{O}_1}$ est le conoyau du
morphisme de $\mathcal{O}_2$-modules
(\ref{equation-define-module-differentials}). De plus, le morphisme
$\text{d}$ est caractérisé par le fait que $\text{d}f$ est l'image de la
section locale $[f]$.

\begin{lemma}
\label{lemma-differentials-sheafify}
Soit $X$ un espace topologique. Soit
$\varphi : \mathcal{O}_1 \to \mathcal{O}_2$ un homomorphisme de faisceaux
d'anneaux sur $X$. Alors $\Omega_{\mathcal{O}_2/\mathcal{O}_1}$ est le
faisceau associé au préfaisceau
$U \mapsto \Omega_{\mathcal{O}_2(U)/\mathcal{O}_1(U)}$.
\end{lemma}

\begin{proof}
Considérons le morphisme (\ref{equation-define-module-differentials}). Il
existe un morphisme analogue de préfaisceaux dont la valeur sur l'ouvert
$U$ est
$$
\mathcal{O}_2(U)[\mathcal{O}_2(U) \times \mathcal{O}_2(U)] \oplus
\mathcal{O}_2(U)[\mathcal{O}_2(U) \times \mathcal{O}_2(U)] \oplus
\mathcal{O}_2(U)[\mathcal{O}_1(U)]
\longrightarrow
\mathcal{O}_2(U)[\mathcal{O}_2(U)]
$$
Par la construction du module des différentielles dans Algèbre, définition
\ref{algebra-definition-differentials}, le conoyau de ce morphisme a pour
valeur $\Omega_{\mathcal{O}_2(U)/\mathcal{O}_1(U)}$ sur $U$. D'autre part,
les faisceaux figurant dans (\ref{equation-define-module-differentials}) sont
les faisceaux associés aux préfaisceaux ci-dessus. Le résultat découle donc
de l'exactitude de la faisceautisation.
\end{proof}

\begin{lemma}
\label{lemma-localize-differentials}
Soit $X$ un espace topologique. Soit
$\varphi : \mathcal{O}_1 \to \mathcal{O}_2$ un homomorphisme de faisceaux
d'anneaux. Pour tout ouvert $U \subset X$, il existe un isomorphisme canonique
$$
\Omega_{\mathcal{O}_2/\mathcal{O}_1}|_U =
\Omega_{(\mathcal{O}_2|_U)/(\mathcal{O}_1|_U)}
$$
compatible aux dérivations universelles.
\end{lemma}

\begin{proof}
Cela résulte du fait que $\Omega_{\mathcal{O}_2/\mathcal{O}_1}$ est le
conoyau du morphisme (\ref{equation-define-module-differentials}).
\end{proof}

\begin{lemma}
\label{lemma-pullback-differentials}
Soit $f : Y \to X$ une application continue d'espaces topologiques.
Soit $\varphi : \mathcal{O}_1 \to \mathcal{O}_2$ un homomorphisme de
faisceaux d'anneaux sur $X$. Il existe alors une identification canonique
$f^{-1}\Omega_{\mathcal{O}_2/\mathcal{O}_1} =
\Omega_{f^{-1}\mathcal{O}_2/f^{-1}\mathcal{O}_1}$
compatible aux dérivations universelles.
\end{lemma}

\begin{proof}
Cela résulte du fait que le faisceau
$\Omega_{\mathcal{O}_2/\mathcal{O}_1}$ est le conoyau du morphisme
(\ref{equation-define-module-differentials}) et qu'un énoncé analogue vaut
pour $\Omega_{f^{-1}\mathcal{O}_2/f^{-1}\mathcal{O}_1}$, puisque le foncteur
$f^{-1}$ est exact et que
$f^{-1}(\mathcal{O}_2[\mathcal{O}_2]) =
f^{-1}\mathcal{O}_2[f^{-1}\mathcal{O}_2]$,
$f^{-1}(\mathcal{O}_2[\mathcal{O}_2 \times \mathcal{O}_2]) =
f^{-1}\mathcal{O}_2[f^{-1}\mathcal{O}_2 \times f^{-1}\mathcal{O}_2]$, and
$f^{-1}(\mathcal{O}_2[\mathcal{O}_1]) =
f^{-1}\mathcal{O}_2[f^{-1}\mathcal{O}_1]$.
\end{proof}

\begin{lemma}
\label{lemma-stalk-module-differentials}
Soit $X$ un espace topologique. Soit
$\mathcal{O}_1 \to \mathcal{O}_2$ un homomorphisme de faisceaux d'anneaux
sur $X$. Pour $x \in X$, on a
$\Omega_{\mathcal{O}_2/\mathcal{O}_1, x} =
\Omega_{\mathcal{O}_{2, x}/\mathcal{O}_{1, x}}$.
\end{lemma}

\begin{proof}
C'est le cas particulier du lemme \ref{lemma-pullback-differentials}
correspondant à l'inclusion $\{x\} \to X$. On peut aussi utiliser le lemme
\ref{lemma-differentials-sheafify}, Faisceaux, lemme
\ref{sheaves-lemma-stalk-sheafification}, et Algèbre, lemme
\ref{algebra-lemma-colimit-differentials}.
\end{proof}

\begin{lemma}
\label{lemma-functoriality-differentials}
Soit $X$ un espace topologique. Soit
$$
\xymatrix{
\mathcal{O}_2 \ar[r]_\varphi & \mathcal{O}_2' \\
\mathcal{O}_1 \ar[r] \ar[u] & \mathcal{O}'_1 \ar[u]
}
$$
un diagramme commutatif de faisceaux d'anneaux sur $X$. La composée de
$\mathcal{O}_2 \to \mathcal{O}'_2$ avec le morphisme
$\text{d} : \mathcal{O}'_2 \to \Omega_{\mathcal{O}'_2/\mathcal{O}'_1}$
est une $\mathcal{O}_1$-dérivation. On obtient donc un morphisme canonique
de $\mathcal{O}_2$-modules
$\Omega_{\mathcal{O}_2/\mathcal{O}_1} \to
\Omega_{\mathcal{O}'_2/\mathcal{O}'_1}$.
Il est uniquement caractérisé par la propriété
$\text{d}(f) \mapsto \text{d}(\varphi(f))$
pour toute section locale $f$ de $\mathcal{O}_2$. Ainsi,
$\Omega_{-/-}$ devient un foncteur sur la catégorie des flèches de faisceaux
d'anneaux.
\end{lemma}

\begin{proof}
L'énoncé découle immédiatement des définitions.
\end{proof}

\begin{lemma}
\label{lemma-differential-seq}
Dans le lemme \ref{lemma-functoriality-differentials}, supposons que
$\mathcal{O}_2 \to \mathcal{O}'_2$ soit surjectif, de noyau
$\mathcal{I} \subset \mathcal{O}_2$, et que
$\mathcal{O}_1 = \mathcal{O}'_1$. Il existe alors une suite exacte canonique
de $\mathcal{O}'_2$-modules
$$
\mathcal{I}/\mathcal{I}^2
\longrightarrow
\Omega_{\mathcal{O}_2/\mathcal{O}_1} \otimes_{\mathcal{O}_2} \mathcal{O}'_2
\longrightarrow
\Omega_{\mathcal{O}'_2/\mathcal{O}_1}
\longrightarrow
0
$$
Le morphisme de gauche est caractérisé par le fait qu'une section locale
$f$ de $\mathcal{I}$ est envoyée sur $\text{d}f \otimes 1$.
\end{lemma}

\begin{proof}
Pour une section locale $f$ de $\mathcal{I}$, notons $\overline{f}$ l'image
de $f$ dans $\mathcal{I}/\mathcal{I}^2$. Pour montrer que l'application
$\overline{f} \mapsto \text{d}f \otimes 1$ est bien définie, il suffit de
vérifier que $\text{d} f_1f_2 \otimes 1 = 0$ lorsque $f_1, f_2$ sont des
sections locales de $\mathcal{I}$. Cela découle de la règle de Leibniz :
$\text{d} f_1f_2 \otimes 1 =
(f_1 \text{d}f_2 + f_2 \text{d} f_1 )\otimes 1 =
\text{d}f_2 \otimes f_1 + \text{d}f_1 \otimes f_2 = 0$.
Un calcul analogue montre que cette application est
$\mathcal{O}'_2 = \mathcal{O}_2/\mathcal{I}$-linéaire. Le morphisme de droite
est celui du lemme \ref{lemma-functoriality-differentials}. Pour vérifier
l'exactitude de la suite, on peut raisonner sur les fibres (lemme
\ref{lemma-abelian}). Le lemme \ref{lemma-stalk-module-differentials} ramène
alors l'assertion à Algèbre, lemme \ref{algebra-lemma-differential-seq}.
\end{proof}

\begin{definition}
\label{definition-differentials}
Soit $(f, f^\sharp) : (X, \mathcal{O}_X) \to (S, \mathcal{O}_S)$
un morphisme d'espaces annelés.
\begin{enumerate}
\item Soit $\mathcal{F}$ un $\mathcal{O}_X$-module. Une
{\it $S$-dérivation} à valeurs dans $\mathcal{F}$ est une
$f^{-1}\mathcal{O}_S$-dérivation, ou plus précisément une
$f^\sharp$-dérivation au sens de la définition
\ref{definition-derivation}. On note
$\text{Der}_S(\mathcal{O}_X, \mathcal{F})$ l'ensemble des $S$-dérivations
à valeurs dans $\mathcal{F}$.
\item Le {\it faisceau des différentielles $\Omega_{X/S}$ de $X$ sur $S$}
est le module des différentielles
$\Omega_{\mathcal{O}_X/f^{-1}\mathcal{O}_S}$ muni de sa $S$-dérivation
universelle $\text{d}_{X/S} : \mathcal{O}_X \to \Omega_{X/S}$.
\end{enumerate}
\end{definition}

\noindent
Voici une situation particulière dans laquelle les dérivations apparaissent
naturellement.

\begin{lemma}
\label{lemma-double-structure-gives-derivation}
Soit $(f, f^\sharp) : (X, \mathcal{O}_X) \to (S, \mathcal{O}_S)$
un morphisme d'espaces annelés. Considérons une suite exacte courte
$$
0 \to \mathcal{I} \to \mathcal{A} \to \mathcal{O}_X \to 0
$$
Ici, $\mathcal{A}$ est un faisceau de $f^{-1}\mathcal{O}_S$-algèbres,
$\pi : \mathcal{A} \to \mathcal{O}_X$ est un morphisme surjectif de
faisceaux de $f^{-1}\mathcal{O}_S$-algèbres, et
$\mathcal{I} = \Ker(\pi)$ est son noyau. Supposons que $\mathcal{I}$ soit
un faisceau d'idéaux de carré nul dans $\mathcal{A}$. Alors $\mathcal{I}$
possède une structure naturelle de $\mathcal{O}_X$-module.
Une section $s : \mathcal{O}_X \to \mathcal{A}$ de $\pi$ est un morphisme
de $f^{-1}\mathcal{O}_S$-algèbres tel que $\pi \circ s = \text{id}$.
Étant données une section $s : \mathcal{O}_X \to \mathcal{A}$ de $\pi$
et une $S$-dérivation $D : \mathcal{O}_X \to \mathcal{I}$, l'application
$$
s + D : \mathcal{O}_X \to \mathcal{A}
$$
est une section de $\pi$, et toute section $s'$ est de la forme $s + D$ pour
une unique $S$-dérivation $D$.
\end{lemma}

\begin{proof}
Rappelons que la structure de $\mathcal{O}_X$-module sur $\mathcal{I}$ est
donnée par $h \tau = \tilde h \tau$ (produit dans $\mathcal{A}$), où $h$
est une section locale de $\mathcal{O}_X$, où $\tilde h$ est un relevé
local de $h$ à $\mathcal{A}$, et où $\tau$ est une section locale de
$\mathcal{I}$. En particulier, une fois $s$ fixée, on peut prendre
$\tilde h = s(h)$. Pour vérifier que $s + D$ est un homomorphisme de
faisceaux d'anneaux, calculons
\begin{eqnarray*}
(s + D)(ab) & = & s(ab) + D(ab) \\
& = & s(a)s(b) + aD(b) + D(a)b \\
& = & s(a) s(b) + s(a)D(b) + D(a)s(b) \\
& = & (s(a) + D(a))(s(b) + D(b))
\end{eqnarray*}
par la règle de Leibniz. De même, on montre que $s + D$ est un morphisme de
$f^{-1}\mathcal{O}_S$-algèbres, puisque $D$ est une $S$-dérivation.
Réciproquement, étant donnée $s'$, on pose $D = s' - s$. Les détails sont omis.
\end{proof}

\begin{lemma}
\label{lemma-functoriality-differentials-ringed-spaces}
Soit
$$
\xymatrix{
X' \ar[d]_{h'} \ar[r]_f & X \ar[d]^h \\
S' \ar[r]^g & S
}
$$
un diagramme commutatif d'espaces annelés.
\begin{enumerate}
\item La composée du morphisme canonique
$\mathcal{O}_X \to f_*\mathcal{O}_{X'}$ avec
$f_*\text{d}_{X'/S'} : f_*\mathcal{O}_{X'} \to f_*\Omega_{X'/S'}$ est une
$S$-dérivation; on obtient donc un morphisme canonique de
$\mathcal{O}_X$-modules
$\Omega_{X/S} \to f_*\Omega_{X'/S'}$.
\item Le diagramme commutatif
$$
\xymatrix{
f^{-1}\mathcal{O}_X \ar[r] & \mathcal{O}_{X'} \\
f^{-1}h^{-1}\mathcal{O}_S \ar[u] \ar[r] & (h')^{-1}\mathcal{O}_{S'} \ar[u]
}
$$
induit, par les lemmes \ref{lemma-pullback-differentials} et
\ref{lemma-functoriality-differentials}, un morphisme canonique
$f^{-1}\Omega_{X/S} \to \Omega_{X'/S'}$.
\end{enumerate}
Ces deux morphismes correspondent (par adjonction entre $f_*$ et $f^*$, et
en utilisant $f^*\Omega_{X/S} =
f^{-1}\Omega_{X/S} \otimes_{f^{-1}\mathcal{O}_X} \mathcal{O}_{X'}$ et
Faisceaux, lemme \ref{sheaves-lemma-adjointness-tensor-restrict})
au même homomorphisme de $\mathcal{O}_{X'}$-modules
$$
c_f : f^*\Omega_{X/S} \longrightarrow \Omega_{X'/S'}
$$
uniquement caractérisé par la propriété que
$f^*\text{d}_{X/S}(a)$ soit envoyé sur $\text{d}_{X'/S'}(f^*a)$
pour toute section locale $a$ de $\mathcal{O}_X$.
\end{lemma}

\begin{proof}
Démonstration omise.
\end{proof}

\begin{lemma}
\label{lemma-check-functoriality-differentials}
Soit
$$
\xymatrix{
X'' \ar[d] \ar[r]_g & X' \ar[d] \ar[r]_f & X \ar[d] \\
S'' \ar[r] & S' \ar[r] & S
}
$$
un diagramme commutatif d'espaces annelés. Avec les notations du lemme
\ref{lemma-functoriality-differentials-ringed-spaces}, on a
$$
c_{f \circ g} = c_g \circ g^* c_f
$$
comme morphismes $(f \circ g)^*\Omega_{X/S} \to \Omega_{X''/S''}$.
\end{lemma}

\begin{proof}
Démonstration omise.
\end{proof}

\begin{lemma}
\label{lemma-triangle-differentials}
Soient $f : X \to Y$ et $g : Y \to S$ des morphismes d'espaces annelés.
Il existe alors une suite exacte canonique
$$
f^*\Omega_{Y/S} \to \Omega_{X/S} \to \Omega_{X/Y} \to 0
$$
dont les flèches proviennent d'applications du lemme
\ref{lemma-functoriality-differentials-ringed-spaces}.
\end{lemma}

\begin{proof}
En prenant les morphismes induits sur les fibres en $x \in X$ et en utilisant
le lemme \ref{lemma-stalk-module-differentials}, on obtient la suite
$$
\mathcal{O}_{X, x} \otimes_{\mathcal{O}_{Y, f(x)}}
\Omega_{\mathcal{O}_{Y, f(x)}/\mathcal{O}_{S, g(f(x))}} \to
\Omega_{\mathcal{O}_{X, x}/\mathcal{O}_{S, g(f(x))}} \to
\Omega_{\mathcal{O}_{X, x}/\mathcal{O}_{Y, f(x)}} \to 0
$$
Il suffit de vérifier que les flèches de cette suite sont celles d'Algèbre,
lemme \ref{algebra-lemma-exact-sequence-differentials}. Cela résulte de la
caractérisation des flèches dans le lemme
\ref{lemma-functoriality-differentials-ringed-spaces} et du fait que,
via l'isomorphisme de Faisceaux, lemme
\ref{sheaves-lemma-stalk-pullback-modules}, on a
$(f^*s)_x = s_x \otimes 1$ pour toute section locale $s$ d'un faisceau de
$\mathcal{O}_Y$-modules.
\end{proof}


















\section{Opérateurs différentiels d'ordre fini}
\label{section-differential-operators}

\noindent
Dans cette section, nous introduisons les opérateurs différentiels d'ordre
fini. Nous conseillons au lecteur de consulter la section correspondante du
chapitre d'algèbre commutative (Algèbre, section
\ref{algebra-section-differential-operators}).

\begin{definition}
\label{definition-differential-operators}
Soit $X$ un espace topologique. Soit
$\varphi : \mathcal{O}_1 \to \mathcal{O}_2$ un homomorphisme de faisceaux
d'anneaux sur $X$. Soit $k \geq 0$ un entier. Soient $\mathcal{F}$ et
$\mathcal{G}$ des faisceaux de $\mathcal{O}_2$-modules. Un
{\it opérateur différentiel $D : \mathcal{F} \to \mathcal{G}$ d'ordre $k$}
est une application $\mathcal{O}_1$-linéaire telle que, pour toute section
locale $g$ de $\mathcal{O}_2$, l'application
$s \mapsto D(gs) - gD(s)$ soit un opérateur différentiel d'ordre $k - 1$.
Dans le cas initial $k = 0$, on définit un opérateur différentiel d'ordre $0$
comme une application $\mathcal{O}_2$-linéaire.
\end{definition}

\noindent
Si $D : \mathcal{F} \to \mathcal{G}$ est un opérateur différentiel d'ordre
$k$, alors, pour toute section locale $g$ de $\mathcal{O}_2$, l'application
$gD$ est encore un opérateur différentiel d'ordre $k$. La somme de deux
opérateurs différentiels d'ordre $k$ en est encore un. Par conséquent,
l'ensemble de tous ces opérateurs
$$
\text{Diff}^k(\mathcal{F}, \mathcal{G}) =
\text{Diff}^k_{\mathcal{O}_2/\mathcal{O}_1}(\mathcal{F}, \mathcal{G})
$$
est un $\Gamma(X, \mathcal{O}_2)$-module. On a
$$
\text{Diff}^0(\mathcal{F}, \mathcal{G}) \subset
\text{Diff}^1(\mathcal{F}, \mathcal{G}) \subset
\text{Diff}^2(\mathcal{F}, \mathcal{G}) \subset \ldots
$$
La règle qui, à tout ouvert $U \subset X$, associe le module des opérateurs
différentiels $D : \mathcal{F}|_U \to \mathcal{G}|_U$ d'ordre $k$ définit
un faisceau de $\mathcal{O}_2$-modules sur $X$. On obtient ainsi un faisceau
d'opérateurs différentiels (si nous en avons besoin, nous ajouterons ici une
définition).

\begin{lemma}
\label{lemma-composition-differential-operators}
Soit $X$ un espace topologique. Soit
$\mathcal{O}_1 \to \mathcal{O}_2$ un morphisme de faisceaux d'anneaux sur
$X$. Soient $\mathcal{E}, \mathcal{F}, \mathcal{G}$ des faisceaux de
$\mathcal{O}_2$-modules. Si $D : \mathcal{E} \to \mathcal{F}$ et
$D' : \mathcal{F} \to \mathcal{G}$ sont des opérateurs différentiels
d'ordres respectifs $k$ et $k'$, alors $D' \circ D$ est un opérateur
différentiel d'ordre $k + k'$.
\end{lemma}

\begin{proof}
Soit $g$ une section locale de $\mathcal{O}_2$. L'application qui envoie une
section locale $x$ de $\mathcal{E}$ sur
$$
D'(D(gx)) - gD'(D(x)) = D'(D(gx)) - D'(gD(x)) + D'(gD(x)) - gD'(D(x))
$$
est la somme de deux composées d'opérateurs différentiels d'ordre inférieur.
Le lemme en résulte par récurrence sur $k + k'$.
\end{proof}

\begin{lemma}
\label{lemma-module-principal-parts}
Soit $X$ un espace topologique. Soit
$\mathcal{O}_1 \to \mathcal{O}_2$ un morphisme de faisceaux d'anneaux sur
$X$. Soit $\mathcal{F}$ un faisceau de $\mathcal{O}_2$-modules et soit
$k \geq 0$. Il existe un faisceau de $\mathcal{O}_2$-modules
$\mathcal{P}^k_{\mathcal{O}_2/\mathcal{O}_1}(\mathcal{F})$
et un isomorphisme canonique
$$
\text{Diff}^k_{\mathcal{O}_2/\mathcal{O}_1}(\mathcal{F}, \mathcal{G}) =
\Hom_{\mathcal{O}_2}(
\mathcal{P}^k_{\mathcal{O}_2/\mathcal{O}_1}(\mathcal{F}), \mathcal{G})
$$
fonctoriel par rapport au $\mathcal{O}_2$-module $\mathcal{G}$.
\end{lemma}

\begin{proof}
L'existence résulte d'arguments généraux de théorie des catégories
(insérer ici une référence ultérieure), mais nous donnerons aussi une
construction directe, qui sera utile dans les démonstrations à venir. Nous
emploierons librement les notations introduites dans la démonstration du
lemme \ref{lemma-universal-module}. À tout opérateur différentiel
$D : \mathcal{F} \to \mathcal{G}$ est associé un morphisme
$\mathcal{O}_2$-linéaire
$L_D : \mathcal{O}_2[\mathcal{F}] \to \mathcal{G}$
envoyant $[m]$ sur $D(m)$. Si $D$ est d'ordre $0$, alors $L_D$ annule les
sections locales
$$
[m + m'] - [m] - [m'],\quad
g_0[m] - [g_0m]
$$
où $g_0$ est une section locale de $\mathcal{O}_2$ et $m, m'$ sont des
sections locales de $\mathcal{F}$. Si $D$ est d'ordre $1$, alors $L_D$
annule les sections locales
$$
[m + m'] - [m] - [m'],\quad
f[m] - [fm], \quad
g_0g_1[m] - g_0[g_1m] - g_1[g_0m] + [g_1g_0m]
$$
où $f$ est une section locale de $\mathcal{O}_1$, où $g_0, g_1$ sont des
sections locales de $\mathcal{O}_2$, et où $m, m'$ sont des sections locales
de $\mathcal{F}$. Si $D$ est d'ordre $k$, alors $L_D$ annule les sections
locales $[m + m'] - [m] - [m']$, $f[m] - [fm]$, ainsi que les sections locales
$$
g_0g_1\ldots g_k[m] - \sum g_0 \ldots \hat g_i \ldots g_k[g_im] + \ldots
+(-1)^{k + 1}[g_0\ldots g_km]
$$
Réciproquement, si $L : \mathcal{O}_2[\mathcal{F}] \to \mathcal{G}$ est un
morphisme $\mathcal{O}_2$-linéaire annulant toutes les sections locales
énumérées dans la phrase précédente, alors $m \mapsto L([m])$ est un
opérateur différentiel d'ordre $k$. Ainsi,
$\mathcal{P}^k_{\mathcal{O}_2/\mathcal{O}_1}(\mathcal{F})$
est le quotient de $\mathcal{O}_2[\mathcal{F}]$ par le
$\mathcal{O}_2$-sous-module engendré par ces sections locales.
\end{proof}

\begin{definition}
\label{definition-module-principal-parts}
Soit $X$ un espace topologique. Soit
$\mathcal{O}_1 \to \mathcal{O}_2$ un morphisme de faisceaux d'anneaux sur
$X$. Soit $\mathcal{F}$ un faisceau de $\mathcal{O}_2$-modules. Le module
$\mathcal{P}^k_{\mathcal{O}_2/\mathcal{O}_1}(\mathcal{F})$ construit dans le
lemme \ref{lemma-module-principal-parts} est appelé le
{\it module des parties principales d'ordre $k$} de $\mathcal{F}$.
\end{definition}

\noindent
Remarquons que les inclusions
$$
\text{Diff}^0(\mathcal{F}, \mathcal{G}) \subset
\text{Diff}^1(\mathcal{F}, \mathcal{G}) \subset
\text{Diff}^2(\mathcal{F}, \mathcal{G}) \subset \ldots
$$
correspondent, par le lemme de Yoneda (Catégories, lemme
\ref{categories-lemma-yoneda}), aux surjections
$$
\ldots \to \mathcal{P}^2_{\mathcal{O}_2/\mathcal{O}_1}(\mathcal{F})
\to \mathcal{P}^1_{\mathcal{O}_2/\mathcal{O}_1}(\mathcal{F})
\to \mathcal{P}^0_{\mathcal{O}_2/\mathcal{O}_1}(\mathcal{F}) = \mathcal{F}
$$

\begin{lemma}
\label{lemma-differential-operators-sheafify}
Soit $X$ un espace topologique. Soit
$\mathcal{O}_1 \to \mathcal{O}_2$ un homomorphisme de préfaisceaux d'anneaux
sur $X$. Soit $\mathcal{F}$ un préfaisceau de $\mathcal{O}_2$-modules. Alors
$\mathcal{P}^k_{\mathcal{O}_2^\#/\mathcal{O}_1^\#}(\mathcal{F}^\#)$
est le faisceau associé au préfaisceau
$U \mapsto P^k_{\mathcal{O}_2(U)/\mathcal{O}_1(U)}(\mathcal{F}(U))$.
\end{lemma}

\begin{proof}
On peut le démontrer exactement comme pour le faisceau des différentielles
dans le lemme \ref{lemma-differentials-sheafify}. Une méthode peut-être plus
satisfaisante consiste à utiliser directement la propriété universelle du
lemme \ref{lemma-module-principal-parts}. Nous omettons les détails.
\end{proof}

\begin{lemma}
\label{lemma-sequence-of-principal-parts}
Soit $X$ un espace topologique. Soit
$\mathcal{O}_1 \to \mathcal{O}_2$ un homomorphisme de faisceaux d'anneaux
sur $X$. Soit $\mathcal{F}$ un faisceau de $\mathcal{O}_2$-modules. Il
existe une suite exacte courte canonique
$$
0 \to
\Omega_{\mathcal{O}_2/\mathcal{O}_1} \otimes_{\mathcal{O}_2} \mathcal{F} \to
\mathcal{P}^1_{\mathcal{O}_2/\mathcal{O}_1}(\mathcal{F}) \to
\mathcal{F} \to 0
$$
fonctorielle en $\mathcal{F}$, appelée la
{\it suite des parties principales}.
\end{lemma}

\begin{proof}
Cela résulte de la version en algèbre commutative (Algèbre, lemme
\ref{algebra-lemma-sequence-of-principal-parts}) et des lemmes
\ref{lemma-differentials-sheafify} et
\ref{lemma-differential-operators-sheafify}.
\end{proof}

\begin{remark}
\label{remark-functoriality-principal-parts}
Soit $X$ un espace topologique. Supposons donné un diagramme commutatif de
faisceaux d'anneaux
$$
\xymatrix{
\mathcal{B} \ar[r] & \mathcal{B}' \\
\mathcal{A} \ar[u] \ar[r] & \mathcal{A}' \ar[u]
}
$$
sur $X$, ainsi qu'un $\mathcal{B}$-module $\mathcal{F}$, un
$\mathcal{B}'$-module $\mathcal{F}'$ et un morphisme $\mathcal{B}$-linéaire
$\mathcal{F} \to \mathcal{F}'$. On obtient alors un système compatible de
morphismes de modules
$$
\xymatrix{
\ldots \ar[r] &
\mathcal{P}^2_{\mathcal{B}'/\mathcal{A}'}(\mathcal{F}') \ar[r] &
\mathcal{P}^1_{\mathcal{B}'/\mathcal{A}'}(\mathcal{F}') \ar[r] &
\mathcal{P}^0_{\mathcal{B}'/\mathcal{A}'}(\mathcal{F}') \\
\ldots \ar[r] &
\mathcal{P}^2_{\mathcal{B}/\mathcal{A}}(\mathcal{F}) \ar[r] \ar[u] &
\mathcal{P}^1_{\mathcal{B}/\mathcal{A}}(\mathcal{F}) \ar[r] \ar[u] &
\mathcal{P}^0_{\mathcal{B}/\mathcal{A}}(\mathcal{F}) \ar[u]
}
$$
Ces morphismes sont compatibles à toute composition ultérieure de
morphismes de ce type. Le moyen le plus simple de le voir consiste à utiliser
la description des modules
$\mathcal{P}^k_{\mathcal{B}/\mathcal{A}}(\mathcal{M})$ par générateurs
(locaux) et relations donnée dans la démonstration du lemme
\ref{lemma-module-principal-parts}; on peut aussi le déduire directement de
leur propriété universelle. De plus, ces morphismes sont compatibles aux
suites exactes courtes du lemme \ref{lemma-sequence-of-principal-parts}.
\end{remark}

\noindent
Étendons maintenant notre définition aux morphismes d'espaces annelés.

\begin{definition}
\label{definition-relative-differential-operators}
Soit $(f, f^\sharp) : (X, \mathcal{O}_X) \to (S, \mathcal{O}_S)$
un morphisme d'espaces annelés. Soient $\mathcal{F}$ et $\mathcal{G}$ des
$\mathcal{O}_X$-modules et soit $k \geq 0$ un entier. Un
{\it opérateur différentiel d'ordre $k$ sur $X/S$} est un opérateur
différentiel $D : \mathcal{F} \to \mathcal{G}$ relativement à
$f^\sharp : f^{-1}\mathcal{O}_S \to \mathcal{O}_X$. On note
$\text{Diff}^k_{X/S}(\mathcal{F}, \mathcal{G})$ l'ensemble de ces opérateurs
différentiels.
\end{definition}












\section{Le complexe de de Rham}
\label{section-de-rham-complex}

\noindent
Cette section est l'analogue d'Algèbre, section
\ref{algebra-section-de-rham-complex}, pour les morphismes d'espaces annelés.
Nous recommandons vivement au lecteur de commencer par cette section.

\medskip\noindent
Soit $X$ un espace topologique. Soit $\mathcal{A} \to \mathcal{B}$ un
homomorphisme de faisceaux d'anneaux. Notons
$\text{d} : \mathcal{B} \to \Omega_{\mathcal{B}/\mathcal{A}}$
le module des différentielles muni de sa $\mathcal{A}$-dérivation universelle,
construit dans la section \ref{section-differentials}. Posons
$$
\Omega_{\mathcal{B}/\mathcal{A}}^i =
\wedge^i_\mathcal{B}(\Omega_{\mathcal{B}/\mathcal{A}})
$$
pour $i \geq 0$; c'est la $i$-ième puissance extérieure définie dans la
section \ref{section-symmetric-exterior}.

\begin{definition}
\label{definition-de-rham-complex}
Dans la situation ci-dessus, le
{\it complexe de de Rham de $\mathcal{B}$ sur $\mathcal{A}$} est l'unique
complexe
$$
\Omega_{\mathcal{B}/\mathcal{A}}^0 \to
\Omega_{\mathcal{B}/\mathcal{A}}^1 \to
\Omega_{\mathcal{B}/\mathcal{A}}^2 \to \ldots
$$
de faisceaux de $\mathcal{A}$-modules dont la différentielle en degré $0$
est $\text{d} : \mathcal{B} \to \Omega_{\mathcal{B}/\mathcal{A}}$ et dont
les différentielles en degrés supérieurs satisfont la propriété suivante :
\begin{equation}
\label{equation-rule}
\text{d}\left(b_0\text{d}b_1 \wedge \ldots \wedge \text{d}b_p\right) =
\text{d}b_0 \wedge \text{d}b_1 \wedge \ldots \wedge \text{d}b_p
\end{equation}
où $b_0, \ldots, b_p \in \mathcal{B}(U)$ sont des sections définies sur un même
ouvert $U \subset X$.
\end{definition}

\noindent
On pourrait construire ce complexe en reprenant les arguments assez lourds
d'Algèbre, section \ref{algebra-section-de-rham-complex}. Rappelons plutôt
que $\Omega_{\mathcal{B}/\mathcal{A}}$ est le faisceau associé au préfaisceau
$U \mapsto \Omega_{\mathcal{B}(U)/\mathcal{A}(U)}$; voir le lemme
\ref{lemma-differentials-sheafify}. Ainsi,
$\Omega_{\mathcal{B}/\mathcal{A}}^i$ est le faisceau associé au préfaisceau
$U \mapsto \Omega^i_{\mathcal{B}(U)/\mathcal{A}(U)}$; voir le lemme
\ref{lemma-local-tensor-algebra}. On peut donc définir le complexe de de Rham
en faisceautisant la règle
$$
U \longmapsto
\Omega^\bullet_{\mathcal{B}(U)/\mathcal{A}(U)}
$$

\begin{lemma}
\label{lemma-pullback-de-rham-complex}
Soit $f : Y \to X$ une application continue d'espaces topologiques. Soit
$\mathcal{A} \to \mathcal{B}$ un homomorphisme de faisceaux d'anneaux sur
$X$. Il existe alors une identification canonique
$f^{-1}\Omega^\bullet_{\mathcal{B}/\mathcal{A}} =
\Omega^\bullet_{f^{-1}\mathcal{B}/f^{-1}\mathcal{A}}$
de complexes de de Rham.
\end{lemma}

\begin{proof}
Démonstration omise. Indication : comparer avec le lemme
\ref{lemma-pullback-differentials}.
\end{proof}

\begin{lemma}
\label{lemma-differentials-de-rham-complex-order-1}
Soit $X$ un espace topologique. Soit $\mathcal{A} \to \mathcal{B}$ un
homomorphisme de faisceaux d'anneaux sur $X$. Les différentielles
$\text{d} : \Omega^i_{\mathcal{B}/\mathcal{A}}  \to
\Omega^{i + 1}_{\mathcal{B}/\mathcal{A}}$
sont des opérateurs différentiels d'ordre $1$.
\end{lemma}

\begin{proof}
Compte tenu de la construction ci-dessus du complexe de de Rham comme
faisceau associé à la règle
$U \mapsto \Omega^\bullet_{\mathcal{B}(U)/\mathcal{A}(U)}$, cela résulte
d'Algèbre, lemme
\ref{algebra-lemma-differentials-de-rham-complex-order-1}.
\end{proof}

\noindent
Soit $X$ un espace topologique. Soit
$$
\xymatrix{
\mathcal{B} \ar[r] & \mathcal{B}' \\
\mathcal{A} \ar[r] \ar[u] & \mathcal{A}' \ar[u]
}
$$
un diagramme commutatif de faisceaux d'anneaux sur $X$. Il existe un
morphisme naturel de complexes de de Rham
$$
\Omega^\bullet_{\mathcal{B}/\mathcal{A}} \longrightarrow
\Omega^\bullet_{\mathcal{B}'/\mathcal{A}'}
$$
En effet, en degré $0$, c'est le morphisme
$\mathcal{B} \to \mathcal{B}'$; en degré $1$, c'est le morphisme
$\Omega_{\mathcal{B}/\mathcal{A}} \to \Omega_{\mathcal{B}'/\mathcal{A}'}$
construit dans la section \ref{section-differentials}; pour $p \geq 2$, c'est
le morphisme induit
$\Omega^p_{\mathcal{B}/\mathcal{A}} =
\wedge^p_\mathcal{B}(\Omega_{\mathcal{B}/\mathcal{A}}) \to
\wedge^p_{\mathcal{B}'}(\Omega_{\mathcal{B}'/\mathcal{A}'}) =
\Omega^p_{\mathcal{B}'/\mathcal{A}'}$.
La compatibilité aux différentielles résulte de leur caractérisation par la
formule (\ref{equation-rule}).

\begin{definition}
\label{definition-de-rham-complex-morphism-ringed-spaces}
Soit $f : (X, \mathcal{O}_X) \to (Y, \mathcal{O}_Y)$ un morphisme
d'espaces annelés. Le {\it complexe de de Rham} de $f$, ou de $X$ sur $Y$,
est le complexe
$$
\Omega^\bullet_{X/Y} = \Omega^\bullet_{\mathcal{O}_X/f^{-1}\mathcal{O}_Y}
$$
\end{definition}

\noindent
Considérons un diagramme commutatif d'espaces annelés
$$
\xymatrix{
X' \ar[d]_{h'} \ar[r]_f & X \ar[d]^h \\
S' \ar[r]^g & S
}
$$
On obtient alors un morphisme canonique
$$
\Omega^\bullet_{X/S} \to f_*\Omega^\bullet_{X'/S'}
$$
de complexes de de Rham. En effet, le diagramme commutatif de faisceaux
d'anneaux
$$
\xymatrix{
f^{-1}\mathcal{O}_X \ar[r] & \mathcal{O}_{X'} \\
f^{-1}h^{-1}\mathcal{O}_S \ar[u] \ar[r] & (h')^{-1}\mathcal{O}_{S'} \ar[u]
}
$$
sur $X'$ fournit un morphisme de complexes (voir ci-dessus)
$$
f^{-1}\Omega^\bullet_{X/S} =
\Omega^\bullet_{f^{-1}\mathcal{O}_X/f^{-1}h^{-1}\mathcal{O}_S}
\longrightarrow
\Omega^\bullet_{\mathcal{O}_{X'}/(h')^{-1}\mathcal{O}_{S'}} =
\Omega^\bullet_{X'/S'}
$$
(en utilisant le lemme \ref{lemma-pullback-de-rham-complex} pour la première
égalité), puis on applique l'adjonction.

\begin{lemma}
\label{lemma-differentials-relative-de-rham-complex-order-1}
Soit $f : X \to Y$ un morphisme d'espaces annelés. Les différentielles
$\text{d} : \Omega^i_{X/Y}  \to \Omega^{i + 1}_{X/Y}$
sont des opérateurs différentiels d'ordre $1$ sur $X/Y$.
\end{lemma}

\begin{proof}
Cela résulte immédiatement du lemme
\ref{lemma-differentials-de-rham-complex-order-1} et de la définition.
\end{proof}





















\section{Le complexe cotangent naïf}
\label{section-netherlander}

\noindent
Cette section est l'analogue d'Algèbre, section
\ref{algebra-section-netherlander}, pour les morphismes d'espaces annelés.
Nous recommandons vivement au lecteur de commencer par cette section.

\medskip\noindent
Soit $X$ un espace topologique. Soit $\mathcal{A} \to \mathcal{B}$ un
homomorphisme de faisceaux d'anneaux. Dans cette section, pour tout faisceau
d'ensembles $\mathcal{E}$ sur $X$, on note $\mathcal{A}[\mathcal{E}]$ le
faisceau associé au préfaisceau
$U \mapsto \mathcal{A}(U)[\mathcal{E}(U)]$. Ici,
$\mathcal{A}(U)[\mathcal{E}(U)]$
désigne l'algèbre de polynômes sur $\mathcal{A}(U)$ dont les variables
correspondent aux éléments de $\mathcal{E}(U)$. On note
$[e] \in \mathcal{A}(U)[\mathcal{E}(U)]$ la variable correspondant à
$e \in \mathcal{E}(U)$. Il existe une surjection canonique de
$\mathcal{A}$-algèbres
\begin{equation}
\label{equation-canonical-presentation}
\mathcal{A}[\mathcal{B}] \longrightarrow \mathcal{B},\quad [b] \longmapsto b
\end{equation}
dont on note le noyau $\mathcal{I} \subset \mathcal{A}[\mathcal{B}]$.
On observe immédiatement que $\mathcal{I}$ est engendré par les sections
locales $[b][b'] - [bb']$ et $[a] - a$. D'après le lemme
\ref{lemma-differential-seq}, il existe un morphisme canonique
\begin{equation}
\label{equation-naive-cotangent-complex}
\mathcal{I}/\mathcal{I}^2
\longrightarrow
\Omega_{\mathcal{A}[\mathcal{B}]/\mathcal{A}}
\otimes_{\mathcal{A}[\mathcal{B}]} \mathcal{B}
\end{equation}
dont le conoyau est canoniquement isomorphe à
$\Omega_{\mathcal{B}/\mathcal{A}}$.

\begin{definition}
\label{definition-naive-cotangent-complex}
Soit $X$ un espace topologique. Soit $\mathcal{A} \to \mathcal{B}$ un
homomorphisme de faisceaux d'anneaux. Le
{\it complexe cotangent naïf} $\NL_{\mathcal{B}/\mathcal{A}}$ est le
complexe de chaînes
(\ref{equation-naive-cotangent-complex})
$$
\NL_{\mathcal{B}/\mathcal{A}} =
\left(\mathcal{I}/\mathcal{I}^2
\longrightarrow
\Omega_{\mathcal{A}[\mathcal{B}]/\mathcal{A}}
\otimes_{\mathcal{A}[\mathcal{B}]} \mathcal{B}\right)
$$
où $\mathcal{I}/\mathcal{I}^2$ est placé en degré $-1$ et où
$\Omega_{\mathcal{A}[\mathcal{B}]/\mathcal{A}}
\otimes_{\mathcal{A}[\mathcal{B}]} \mathcal{B}$
est placé en degré $0$.
\end{definition}

\noindent
Cette construction possède une fonctorialité analogue à celle étudiée dans
le lemme \ref{lemma-functoriality-differentials} pour les modules de
différentielles. Plus précisément, étant donné un diagramme commutatif
\begin{equation}
\label{equation-commutative-square-sheaves}
\vcenter{
\xymatrix{
\mathcal{B} \ar[r] & \mathcal{B}' \\
\mathcal{A} \ar[u] \ar[r] & \mathcal{A}' \ar[u]
}
}
\end{equation}
de faisceaux d'anneaux sur $X$, il existe un morphisme canonique
$\mathcal{B}$-linéaire de complexes
$$
\NL_{\mathcal{B}/\mathcal{A}} \longrightarrow \NL_{\mathcal{B}'/\mathcal{A}'}
$$
En effet, les morphismes du diagramme commutatif donnent un morphisme
canonique $\mathcal{A}[\mathcal{B}] \to \mathcal{A}'[\mathcal{B}']$ qui
envoie $\mathcal{I}$ dans
$\mathcal{I}' = \Ker(\mathcal{A}'[\mathcal{B}'] \to \mathcal{B}')$.
On obtient donc un morphisme
$\mathcal{I}/\mathcal{I}^2 \to \mathcal{I}'/(\mathcal{I}')^2$ et un
morphisme entre les modules de différentielles; ensemble, ils donnent le
morphisme voulu entre les complexes cotangents naïfs. Ce morphisme est
compatible aux compositions au sens suivant : étant donné un diagramme
commutatif
$$
\xymatrix{
\mathcal{B} \ar[r] &
\mathcal{B}' \ar[r] &
\mathcal{B}'' \\
\mathcal{A} \ar[u] \ar[r] &
\mathcal{A}' \ar[u] \ar[r] &
\mathcal{A}'' \ar[u]
}
$$
de faisceaux d'anneaux, la composée
$$
\NL_{\mathcal{B}/\mathcal{A}} \longrightarrow
\NL_{\mathcal{B}'/\mathcal{A}'} \longrightarrow
\NL_{\mathcal{B}''/\mathcal{A}''}
$$
est le morphisme associé au rectangle extérieur.

\medskip\noindent
On peut choisir une autre présentation de $\mathcal{B}$ comme quotient d'une
algèbre de polynômes sur $\mathcal{A}$ et obtenir néanmoins le même objet de
$D(\mathcal{B})$. Pour l'expliquer, soit $\mathcal{E}$ un faisceau
d'ensembles sur $X$ et soit $\alpha : \mathcal{E} \to \mathcal{B}$ un
morphisme de faisceaux d'ensembles. On obtient alors un homomorphisme de
$\mathcal{A}$-algèbres $\mathcal{A}[\mathcal{E}] \to \mathcal{B}$. Si ce
morphisme est surjectif, c'est-à-dire si $\alpha(\mathcal{E})$ engendre
$\mathcal{B}$ comme $\mathcal{A}$-algèbre, on pose
$$
\NL(\alpha) = \left(
\mathcal{J}/\mathcal{J}^2
\longrightarrow
\Omega_{\mathcal{A}[\mathcal{E}]/\mathcal{A}}
\otimes_{\mathcal{A}[\mathcal{E}]} \mathcal{B}\right)
$$
où $\mathcal{J} \subset \mathcal{A}[\mathcal{E}]$ est le noyau de la
surjection $\mathcal{A}[\mathcal{E}] \to \mathcal{B}$. Voici le résultat.

\begin{lemma}
\label{lemma-NL-up-to-qis}
Dans la situation ci-dessus, il existe un isomorphisme canonique
$\NL(\alpha) = \NL_{\mathcal{B}/\mathcal{A}}$ in $D(\mathcal{B})$.
\end{lemma}

\begin{proof}
Remarquons que
$\NL_{\mathcal{B}/\mathcal{A}} = \NL(\text{id}_\mathcal{B})$. Il suffit
donc de montrer que, pour deux morphismes
$\alpha_i : \mathcal{E}_i \to \mathcal{B}$ comme ci-dessus, il existe un
quasi-isomorphisme canonique
$\NL(\alpha_1) = \NL(\alpha_2)$ dans $D(\mathcal{B})$. Pour cela, posons
$\mathcal{E} = \mathcal{E}_1 \amalg \mathcal{E}_2$ et
$\alpha = \alpha_1 \amalg \alpha_2 : \mathcal{E} \to \mathcal{B}$.
Posons
$\mathcal{J}_i = \Ker(\mathcal{A}[\mathcal{E}_i] \to \mathcal{B})$
and
$\mathcal{J} = \Ker(\mathcal{A}[\mathcal{E}] \to \mathcal{B})$.
On obtient des morphismes
$\mathcal{A}[\mathcal{E}_i] \to \mathcal{A}[\mathcal{E}]$ qui envoient
$\mathcal{J}_i$ dans $\mathcal{J}$. On obtient donc des morphismes
canoniques de complexes
$$
\NL(\alpha_i) \longrightarrow \NL(\alpha)
$$
et il suffit de montrer que ce sont des quasi-isomorphismes. Il suffit pour
cela de le vérifier sur les fibres (lemme \ref{lemma-abelian}). Si $x \in X$,
la fibre de $\NL(\alpha)$ est le complexe $\NL(\alpha_x)$ d'Algèbre,
section \ref{algebra-section-netherlander}, associé à la présentation
$\mathcal{A}_x[\mathcal{E}_x] \to \mathcal{B}_x$ provenant du morphisme
$\alpha_x : \mathcal{E}_x \to \mathcal{B}_x$. (Certains détails sont omis;
utiliser le lemme \ref{lemma-stalk-module-differentials} pour vérifier la
compatibilité entre la formation des différentielles et le passage aux
fibres.) Le résultat découle d'Algèbre, lemme
\ref{algebra-lemma-NL-homotopy}.
\end{proof}

\begin{lemma}
\label{lemma-pullback-NL}
Soit $f : X \to Y$ une application continue d'espaces topologiques. Soit
$\mathcal{A} \to \mathcal{B}$ un homomorphisme de faisceaux d'anneaux sur
$Y$. Alors $f^{-1}\NL_{\mathcal{B}/\mathcal{A}} =
\NL_{f^{-1}\mathcal{B}/f^{-1}\mathcal{A}}$.
\end{lemma}

\begin{proof}
Démonstration omise. Indication : utiliser le lemme
\ref{lemma-pullback-differentials}.
\end{proof}

\begin{lemma}
\label{lemma-stalk-NL}
Soit $X$ un espace topologique. Soit $\mathcal{A} \to \mathcal{B}$ un
homomorphisme de faisceaux d'anneaux sur $X$. Pour $x \in X$, on a
$\NL_{\mathcal{B}/\mathcal{A}, x} =
\NL_{\mathcal{B}_x/\mathcal{A}_x}$.
\end{lemma}

\begin{proof}
C'est le cas particulier du lemme \ref{lemma-pullback-NL} correspondant à
l'inclusion $\{x\} \to X$.
\end{proof}

\begin{lemma}
\label{lemma-exact-sequence-NL}
Soit $X$ un espace topologique. Soient
$\mathcal{A} \to \mathcal{B} \to \mathcal{C}$ des morphismes de faisceaux
d'anneaux. Soit $C$ le cône (Catégories dérivées, définition
\ref{derived-definition-cone}) du morphisme de complexes
$\NL_{\mathcal{C}/\mathcal{A}} \to \NL_{\mathcal{C}/\mathcal{B}}$.
Il existe un morphisme canonique
$$
c :
\NL_{\mathcal{B}/\mathcal{A}} \otimes_\mathcal{B} \mathcal{C}
\longrightarrow
C[-1]
$$
de complexes de $\mathcal{C}$-modules qui fournit une suite exacte canonique
à six termes
$$
\xymatrix{
H^0(\NL_{\mathcal{B}/\mathcal{A}} \otimes_\mathcal{B} \mathcal{C}) \ar[r] &
H^0(\NL_{\mathcal{C}/\mathcal{A}}) \ar[r] &
H^0(\NL_{\mathcal{C}/\mathcal{B}}) \ar[r] &
0 \\
H^{-1}(\NL_{\mathcal{B}/\mathcal{A}} \otimes_\mathcal{B} \mathcal{C}) \ar[r] &
H^{-1}(\NL_{\mathcal{C}/\mathcal{A}}) \ar[r] &
H^{-1}(\NL_{\mathcal{C}/\mathcal{B}}) \ar[llu]
}
$$
de faisceaux de cohomologie.
\end{lemma}

\begin{proof}
Pour définir le morphisme $c$, il faut définir un morphisme
$c_1 : \NL_{\mathcal{B}/\mathcal{A}} \otimes_\mathcal{B} \mathcal{C}
\to \NL_{\mathcal{C}/\mathcal{A}}$ ainsi qu'une homotopie explicite entre
la composée
$$
\NL_{\mathcal{B}/\mathcal{A}} \otimes_\mathcal{B} \mathcal{C} \to
\NL_{\mathcal{C}/\mathcal{A}} \to \NL_{\mathcal{C}/\mathcal{B}}
$$
et le morphisme nul; voir Catégories dérivées, lemme
\ref{derived-lemma-map-from-cone}. Pour $c_1$, on utilise la fonctorialité
décrite ci-dessus appliquée au diagramme évident. Pour l'homotopie, on
utilise le morphisme
$$
\NL_{\mathcal{B}/\mathcal{A}}^0 \otimes_\mathcal{B} \mathcal{C}
\longrightarrow
\NL_{\mathcal{C}/\mathcal{B}}^{-1},\quad
\text{d}[b] \otimes 1 \longmapsto [\varphi(b)] - b[1]
$$
où $\varphi : \mathcal{B} \to \mathcal{C}$ est le morphisme donné. Comparer
avec Algèbre, remarque
\ref{algebra-remark-composition-homotopy-equivalent-to-zero}. Pour en déduire
l'assertion sur les faisceaux de cohomologie, il suffit de montrer que
$H^0(c)$ est un isomorphisme et que $H^{-1}(c)$ est surjectif. On peut le
vérifier sur les fibres, en utilisant le lemme \ref{lemma-stalk-NL}, puis
appliquer le résultat correspondant d'algèbre commutative; voir Algèbre,
lemme \ref{algebra-lemma-exact-sequence-NL}. Certains détails sont omis.
\end{proof}

\noindent
Le complexe cotangent naïf d'un morphisme d'espaces annelés se définit au
moyen du complexe construit ci-dessus.

\begin{definition}
\label{definition-cotangent-complex-morphism-ringed-topoi}
Le {\it complexe cotangent naïf} $\NL_f = \NL_{X/Y}$ d'un morphisme d'espaces
annelés $f : (X, \mathcal{O}_X) \to (Y, \mathcal{O}_Y)$ est
$\NL_{\mathcal{O}_X/f^{-1}\mathcal{O}_Y}$.
\end{definition}

\noindent
Étant donné un diagramme commutatif
$$
\xymatrix{
X' \ar[r]_g \ar[d]_{f'} & X \ar[d]^f \\
Y' \ar[r]^h & Y
}
$$
d'espaces annelés, il existe un morphisme canonique
$c : g^*\NL_{X/Y} \to \NL_{X'/Y'}$. C'est le morphisme
$$
g^*\NL_{X/Y} =
\mathcal{O}_{X'} \otimes_{g^{-1}\mathcal{O}_X}
\NL_{g^{-1}\mathcal{O}_X/g^{-1}f^{-1}\mathcal{O}_Y}
\longrightarrow
\NL_{\mathcal{O}_{X'}/(f')^{-1}\mathcal{O}_{Y'}} = \NL_{X'/Y'}
$$
où la flèche provient du diagramme commutatif de faisceaux d'anneaux
$$
\xymatrix{
g^{-1}\mathcal{O}_X \ar[r]_{g^\sharp} &
\mathcal{O}_{X'} \\
g^{-1}f^{-1}\mathcal{O}_Y
\ar[r]^{g^{-1}h^\sharp}
\ar[u]^{g^{-1}f^\sharp} &
(f')^{-1}\mathcal{O}_{Y'} \ar[u]_{(f')^\sharp}
}
$$
comme dans (\ref{equation-commutative-square-sheaves}) ci-dessus. Étant
donné un second diagramme de ce type
$$
\xymatrix{
X'' \ar[r]_{g'} \ar[d] & X' \ar[d] \\
Y'' \ar[r] & Y'
}
$$
la composée de $(g')^*c$ et du morphisme
$c' : (g')^*\NL_{X'/Y'} \to \NL_{X''/Y''}$ est le morphisme
$(g \circ g')^*\NL_{X/Y} \to \NL_{X''/Y''}$.

\begin{lemma}
\label{lemma-exact-sequence-NL-ringed-topoi}
Soient $f : X \to Y$ et $g : Y \to Z$ des morphismes d'espaces annelés.
Soit $C$ le cône du morphisme $\NL_{X/Z} \to \NL_{X/Y}$ de complexes de
$\mathcal{O}_X$-modules. Il existe un morphisme canonique
$$
f^*\NL_{Y/Z} \to C[-1]
$$
qui fournit une suite exacte canonique à six termes
$$
\xymatrix{
H^0(f^*\NL_{Y/Z}) \ar[r] &
H^0(\NL_{X/Z}) \ar[r] &
H^0(\NL_{X/Y}) \ar[r] &
0 \\
H^{-1}(f^*\NL_{Y/Z}) \ar[r] &
H^{-1}(\NL_{X/Z}) \ar[r] &
H^{-1}(\NL_{X/Y}) \ar[llu]
}
$$
de faisceaux de cohomologie.
\end{lemma}

\begin{proof}
Considérons les morphismes de faisceaux d'anneaux
$$
(g \circ f)^{-1}\mathcal{O}_Z \to f^{-1}\mathcal{O}_Y \to \mathcal{O}_X
$$
et appliquons le lemme \ref{lemma-exact-sequence-NL}.
\end{proof}







\begin{multicols}{2}[\section{Autres chapitres}]
\noindent
Préliminaires
\begin{enumerate}
\item \hyperref[introduction-section-phantom]{Introduction}
\item \hyperref[conventions-section-phantom]{Conventions}
\item \hyperref[sets-section-phantom]{Théorie des ensembles}
\item \hyperref[categories-section-phantom]{Catégories}
\item \hyperref[topology-section-phantom]{Topologie}
\item \hyperref[sheaves-section-phantom]{Faisceaux sur les espaces}
\item \hyperref[sites-section-phantom]{Sites et faisceaux}
\item \hyperref[stacks-section-phantom]{Champs}
\item \hyperref[fields-section-phantom]{Corps}
\item \hyperref[algebra-section-phantom]{Algèbre commutative}
\item \hyperref[brauer-section-phantom]{Groupes de Brauer}
\item \hyperref[homology-section-phantom]{Algèbre homologique}
\item \hyperref[derived-section-phantom]{Catégories dérivées}
\item \hyperref[simplicial-section-phantom]{Méthodes simpliciales}
\item \hyperref[more-algebra-section-phantom]{Compléments d'algèbre}
\item \hyperref[smoothing-section-phantom]{Lissage des morphismes d'anneaux}
\item \hyperref[modules-section-phantom]{Faisceaux de modules}
\item \hyperref[sites-modules-section-phantom]{Modules sur les sites}
\item \hyperref[injectives-section-phantom]{Objets injectifs}
\item \hyperref[cohomology-section-phantom]{Cohomologie des faisceaux}
\item \hyperref[sites-cohomology-section-phantom]{Cohomologie sur les sites}
\item \hyperref[dga-section-phantom]{Algèbre différentielle graduée}
\item \hyperref[dpa-section-phantom]{Algèbre à puissances divisées}
\item \hyperref[sdga-section-phantom]{Faisceaux différentiels gradués}
\item \hyperref[hypercovering-section-phantom]{Hyperrecouvrements}
\end{enumerate}
Schémas
\begin{enumerate}
\setcounter{enumi}{25}
\item \hyperref[schemes-section-phantom]{Schémas}
\item \hyperref[constructions-section-phantom]{Constructions de schémas}
\item \hyperref[properties-section-phantom]{Propriétés des schémas}
\item \hyperref[morphisms-section-phantom]{Morphismes de schémas}
\item \hyperref[coherent-section-phantom]{Cohomologie des schémas}
\item \hyperref[divisors-section-phantom]{Diviseurs}
\item \hyperref[limits-section-phantom]{Limites de schémas}
\item \hyperref[varieties-section-phantom]{Variétés}
\item \hyperref[topologies-section-phantom]{Topologies sur les schémas}
\item \hyperref[descent-section-phantom]{Descente}
\item \hyperref[perfect-section-phantom]{Catégories dérivées de schémas}
\item \hyperref[more-morphisms-section-phantom]{Compléments sur les morphismes}
\item \hyperref[flat-section-phantom]{Compléments sur la platitude}
\item \hyperref[groupoids-section-phantom]{Schémas en groupoïdes}
\item \hyperref[more-groupoids-section-phantom]{Compléments sur les schémas en groupoïdes}
\item \hyperref[etale-section-phantom]{Morphismes étales de schémas}
\end{enumerate}
Thèmes de théorie des schémas
\begin{enumerate}
\setcounter{enumi}{41}
\item \hyperref[chow-section-phantom]{Homologie de Chow}
\item \hyperref[intersection-section-phantom]{Théorie de l'intersection}
\item \hyperref[pic-section-phantom]{Schémas de Picard des courbes}
\item \hyperref[weil-section-phantom]{Théories cohomologiques de Weil}
\item \hyperref[adequate-section-phantom]{Modules adéquats}
\item \hyperref[dualizing-section-phantom]{Complexes dualisants}
\item \hyperref[duality-section-phantom]{Dualité pour les schémas}
\item \hyperref[discriminant-section-phantom]{Discriminants et différentes}
\item \hyperref[derham-section-phantom]{Cohomologie de de Rham}
\item \hyperref[local-cohomology-section-phantom]{Cohomologie locale}
\item \hyperref[algebraization-section-phantom]{Géométrie algébrique et formelle}
\item \hyperref[curves-section-phantom]{Courbes algébriques}
\item \hyperref[resolve-section-phantom]{Résolution des surfaces}
\item \hyperref[models-section-phantom]{Réduction semi-stable}
\item \hyperref[functors-section-phantom]{Foncteurs et morphismes}
\item \hyperref[equiv-section-phantom]{Catégories dérivées de variétés}
\item \hyperref[pione-section-phantom]{Groupes fondamentaux de schémas}
\item \hyperref[etale-cohomology-section-phantom]{Cohomologie étale}
\item \hyperref[crystalline-section-phantom]{Cohomologie cristalline}
\item \hyperref[proetale-section-phantom]{Cohomologie pro-étale}
\item \hyperref[relative-cycles-section-phantom]{Cycles relatifs}
\item \hyperref[more-etale-section-phantom]{Compléments de cohomologie étale}
\item \hyperref[trace-section-phantom]{La formule des traces}
\end{enumerate}
Espaces algébriques
\begin{enumerate}
\setcounter{enumi}{64}
\item \hyperref[spaces-section-phantom]{Espaces algébriques}
\item \hyperref[spaces-properties-section-phantom]{Propriétés des espaces algébriques}
\item \hyperref[spaces-morphisms-section-phantom]{Morphismes d'espaces algébriques}
\item \hyperref[decent-spaces-section-phantom]{Espaces algébriques décents}
\item \hyperref[spaces-cohomology-section-phantom]{Cohomologie des espaces algébriques}
\item \hyperref[spaces-limits-section-phantom]{Limites d'espaces algébriques}
\item \hyperref[spaces-divisors-section-phantom]{Diviseurs sur les espaces algébriques}
\item \hyperref[spaces-over-fields-section-phantom]{Espaces algébriques sur des corps}
\item \hyperref[spaces-topologies-section-phantom]{Topologies sur les espaces algébriques}
\item \hyperref[spaces-descent-section-phantom]{Descente et espaces algébriques}
\item \hyperref[spaces-perfect-section-phantom]{Catégories dérivées d'espaces}
\item \hyperref[spaces-more-morphisms-section-phantom]{Compléments sur les morphismes d'espaces}
\item \hyperref[spaces-flat-section-phantom]{Platitude sur les espaces algébriques}
\item \hyperref[spaces-groupoids-section-phantom]{Groupoïdes dans les espaces algébriques}
\item \hyperref[spaces-more-groupoids-section-phantom]{Compléments sur les groupoïdes dans les espaces}
\item \hyperref[bootstrap-section-phantom]{Amorçage}
\item \hyperref[spaces-pushouts-section-phantom]{Sommes amalgamées d'espaces algébriques}
\end{enumerate}
Thèmes de géométrie
\begin{enumerate}
\setcounter{enumi}{81}
\item \hyperref[spaces-chow-section-phantom]{Groupes de Chow des espaces}
\item \hyperref[groupoids-quotients-section-phantom]{Quotients de groupoïdes}
\item \hyperref[spaces-more-cohomology-section-phantom]{Compléments sur la cohomologie des espaces}
\item \hyperref[spaces-simplicial-section-phantom]{Espaces simpliciaux}
\item \hyperref[spaces-duality-section-phantom]{Dualité pour les espaces}
\item \hyperref[formal-spaces-section-phantom]{Espaces algébriques formels}
\item \hyperref[restricted-section-phantom]{Algébrisation des espaces formels}
\item \hyperref[spaces-resolve-section-phantom]{Résolution des surfaces revisitée}
\end{enumerate}
Théorie des déformations
\begin{enumerate}
\setcounter{enumi}{89}
\item \hyperref[formal-defos-section-phantom]{Théorie formelle des déformations}
\item \hyperref[defos-section-phantom]{Théorie des déformations}
\item \hyperref[cotangent-section-phantom]{Le complexe cotangent}
\item \hyperref[examples-defos-section-phantom]{Problèmes de déformation}
\end{enumerate}
Champs algébriques
\begin{enumerate}
\setcounter{enumi}{93}
\item \hyperref[algebraic-section-phantom]{Champs algébriques}
\item \hyperref[examples-stacks-section-phantom]{Exemples de champs}
\item \hyperref[stacks-sheaves-section-phantom]{Faisceaux sur les champs algébriques}
\item \hyperref[criteria-section-phantom]{Critères de représentabilité}
\item \hyperref[artin-section-phantom]{Axiomes d'Artin}
\item \hyperref[quot-section-phantom]{Espaces de Quot et de Hilbert}
\item \hyperref[stacks-properties-section-phantom]{Propriétés des champs algébriques}
\item \hyperref[stacks-morphisms-section-phantom]{Morphismes de champs algébriques}
\item \hyperref[stacks-limits-section-phantom]{Limites de champs algébriques}
\item \hyperref[stacks-cohomology-section-phantom]{Cohomologie des champs algébriques}
\item \hyperref[stacks-perfect-section-phantom]{Catégories dérivées de champs}
\item \hyperref[stacks-introduction-section-phantom]{Introduction aux champs algébriques}
\item \hyperref[stacks-more-morphisms-section-phantom]{Compléments sur les morphismes de champs}
\item \hyperref[stacks-geometry-section-phantom]{La géométrie des champs}
\end{enumerate}
Thèmes de théorie des espaces de modules
\begin{enumerate}
\setcounter{enumi}{107}
\item \hyperref[moduli-section-phantom]{Champs de modules}
\item \hyperref[moduli-curves-section-phantom]{Espaces de modules de courbes}
\end{enumerate}
Divers
\begin{enumerate}
\setcounter{enumi}{109}
\item \hyperref[examples-section-phantom]{Exemples}
\item \hyperref[exercises-section-phantom]{Exercices}
\item \hyperref[guide-section-phantom]{Guide bibliographique}
\item \hyperref[desirables-section-phantom]{Desiderata}
\item \hyperref[coding-section-phantom]{Style de programmation}
\item \hyperref[obsolete-section-phantom]{Obsolète}
\item \hyperref[fdl-section-phantom]{Licence de documentation libre GNU}
\item \hyperref[index-section-phantom]{Index généré automatiquement}
\end{enumerate}
\end{multicols}


\bibliography{my}
\bibliographystyle{amsalpha}

\end{document}
